\documentclass{book}

\usepackage[utf8]{inputenc}  % para que funcionen las tildes
\usepackage{amsmath}
\usepackage{amssymb}
\usepackage{amsthm}
\usepackage{cancel}
\usepackage{letltxmacro}
\LetLtxMacro\amsproof\proof
\LetLtxMacro\amsendproof\endproof
\usepackage{txfonts}
\usepackage{lmodern}
\usepackage[dvipsnames]{xcolor}
\usepackage{thmtools, thm-restate}
\usepackage{datetime} % para la hora de compilación
\usepackage{graphicx}
\graphicspath{{images/}}
\usepackage[spanish,es-noquoting]{babel} % es-noquoting es para que funcione tikz
\usepackage{mathabx} % para \divides
\usepackage{centernot} % para \centernot\inculdes
\usepackage{wrapfig}
\usepackage{multicol}
\usepackage{subcaption}
\usepackage{etoolbox}
\usepackage{xfrac}
\usepackage{glossaries}
\usepackage{standalone}
\usepackage{pgfmath}
\usepackage{pgfplots}
\usepackage{pdfpages}
\usepackage{tikz}
\usepackage{mathtools}
\usepackage{hyperref}
\usepackage{cancel}
\usepackage[shortlabels]{enumitem}
%\usepackage{blkarray} % Labeled matrix :D
\usetikzlibrary{arrows.meta}
\usetikzlibrary{shapes}
\usetikzlibrary{decorations.text}
\usetikzlibrary{positioning}
\usetikzlibrary{external}
\tikzexternalize[prefix=./tikzbuild/]
\tikzexternalize % activate!

\AtBeginDocument{%
  \LetLtxMacro\proof\amsproof
  \LetLtxMacro\endproof\amsendproof
}

\definecolor{coolblack}{rgb}{0.0, 0.18, 0.39}
%\usepackage[a6paper,margin=5mm]{geometry}
\usepackage[a4paper, top=2.5cm,bottom=2.5cm,left=2.5cm,right=2.5cm]{geometry}

%\setlength{\parindent}{0pt}
\usepackage{parskip}

\usepackage{hyperref}

% ESTILOS DE DEFINICIONES Y TEOREMAS
\declaretheoremstyle[
	bodyfont=\normalfont,
	shaded={
		margin=8pt,
		bgcolor=White,
		rulecolor=Black,
		rulewidth=1pt
}]{mythm}

\declaretheoremstyle[
	bodyfont=\normalfont,
	shaded={
		margin=1em,
		bgcolor={rgb}{0.9,0.9,0.9}
}]{mydfn}

\declaretheoremstyle[
bodyfont=\normalfont,
spacebelow=1em,
spaceabove=1em,
]{myej}

\declaretheoremstyle[
bodyfont=\normalfont,
spacebelow=1em,
spaceabove=1em,
shaded={
    margin=8pt,
    bgcolor=White,
    rulecolor=coolblack,
    rulewidth=1pt
}]{thej}

\declaretheoremstyle[
bodyfont=\normalfont,
spacebelow=1em,
spaceabove=1em,
thmbox=M
]{myeg}

% DEFINICIONES DE ENTORNOS DE TEOREMAS
\declaretheorem[
	name=Teorema,
	refname={teorema,teoremas},
	Refname={Teorema,Teoremas},
	style=mythm
]{thm}

\declaretheorem[
	name=Corolario,
	refname={corolario,corolarios},
	Refname={Corolario,Corolarios},
	style=myej
]{cor}

\declaretheorem[
	name=Proposici\'{o}n,
	refname={proposici\'{o}n,proposiciones},
	Refname={Proposici\'{o}n, Proposiciones},
	sharenumber=thm,
	style=myej
]{pro}

\declaretheorem[
name=Lema,
refname={lema,lemas},
Refname={Lema, Lemas},
sharenumber=thm,
style=myej
]{lm}

%\newtheorem*{cor}{Corolario}
\newtheorem*{lem}{Lema}

\declaretheorem[
	name=Definici\'{o}n,
	refname={definici\'{o}n,definiciones},
	Refname={Definici\'{o}n,Definiciones},
	style=mydfn
]{dfn}

\declaretheorem[
name=Ejercicio,
refname={ejercicio,ejercicios},
Refname={Ejercicio,Ejercicios},
style=myej,
numbered=no
]{ex}

\declaretheorem[
name=Ejercicio propuesto,
refname={ejercicio,ejercicios},
Refname={Ejercicio,Ejercicios},
style=thej,
]{th_ex}

\declaretheorem[
name=Ejemplo,
refname={ejemplo,ejemplos},
Refname={Ejemplo,Ejemplos},
style=myeg
]{eg}

\declaretheorem[
name=Observaci\'{o}n,
refname={observaci\'{o}n,observaciones},
Refname={Observaci\'{o}n,Observaciones},
style=myej,
numbered=no
]{obs}

\makeglossaries
\newglossaryentry{di}{name={DI},description={Dominio de integridad (ver \autoref{dfn:dominiointegridad})}}
\newglossaryentry{dip}{name={DIP},description={Dominio de ideales principales (ver \autoref{dfn:dominioidealesprincipales})}}
\newglossaryentry{dfu}{name={DFU},description={Dominio de factorización única (ver \autoref{dfn:dfu})}}
\newglossaryentry{de}{name={DE},description={Dominio euclídeo (ver \autoref{dfn:de})}}


%\setcounter{chapter}{-1}
%COMANDOS ÚTILES PARA DERIVADAS
\AtBeginDocument{\renewcommand{\d}{\mathrm{d}}}
\newcommand{\Dd}[1]{\frac{\mathrm{d}}{\mathrm{d}#1}}
\newcommand{\dd}[2]{\frac{\mathrm{d}#1}{\mathrm{d}#2}}
\newcommand{\Pd}[1]{\frac{\partial}{\partial #1}}
\newcommand{\pd}[2]{\frac{\partial #1}{\partial #2}}
\newcommand{\pdn}[3]{\frac{\partial^{#3} #1}{\partial #2^{#3}}}
\newcommand{\mbf}[1]{\mathbf{#1}}
% COMANDOS ÚTILES PARA LA TEORÍA DE GRUPOS
%\newcommand{\normsub}{\mathbin{\triangleleft}}
\newcommand{\normsub}{\lhd}
\newcommand{\normsubeq}{\trianglelefteq}
\newcommand{\uds}[1]{\mathcal{U}(#1)}
\newcommand{\inv}[1]{#1^{-1}}
\newcommand{\ima}{\text{Im }}
\newcommand{\isom}{\simeq}
\newcommand{\autom}[1]{\text{Aut}(#1)}
\newcommand{\gen}[1]{\langle#1\rangle}
\newcommand{\biy}[1]{\text{Biy}(#1)}
\newcommand{\quo}[2]{{\left.\raisebox{.2em}{$#1$}\middle/\raisebox{-.2em}{$#2$}\right.}}
\newcommand{\appad}[1]{\mathrel{\mkern#1}}
\newcommand{\setdef}[1]{\left\{#1\right\}}
\newcommand{\sdf}[1]{\left\{#1\right\}}
\newcommand{\car}[1]{\mathrm{car}(#1)}
\newcommand{\aut}{\mathrm{Aut}}
\newcommand{\gal}{\mathrm{Gal}}
\newcommand{\actar}{\curvearrowright}
\newcommand{\vabs}[1]{\left|#1\right|}
\newcommand{\rest}[1]{\left.#1\right|}
\renewcommand{\exp}[1]{\mathrm{exp}(#1)}
\newcommand{\mcm}[1]{\mathrm{mcm}\sdf{#1}}
\newcommand{\mtc}[1]{\mathcal{#1}}
%\DeclareMathSymbol{\varprod}{\mathop}{largesymbolsA}{16}

\newcommand{\hr}{\rule{\textwidth}{.4pt}}
\newcommand{\N}{\mathbb{N}}
\newcommand{\Z}{\mathbb{Z}}
\newcommand{\R}{\mathbb{R}}
\newcommand{\Q}{\mathbb{Q}}
\newcommand{\C}{\mathbb{C}}
\newcommand{\A}{\mathbb{A}}
\newcommand{\F}{\mathbb{F}}
\newcommand{\J}{\mathbb{J}}
\newcommand{\B}{\mathbb{B}}
\DeclareMathOperator{\diag}{diag}
\DeclareMathOperator{\traza}{traza}
\renewcommand{\P}{\mathcal{P}}
\renewcommand{\S}{\mathbb{S}}
\newcommand{\ZnZ}{\quo{\Z}{n\Z}}
\newcommand{\ZmZ}{\quo{\Z}{m\Z}}
\newcommand{\0}{\mathbf{0}}
\newcommand{\1}{\mathbf{1}}
\newcommand{\gr}{\text{gr }}
\newcommand\ninf[1]{\left|\left|#1\right|\right|_\infty}
\newcommand\cnj[1]{\overline{#1}}
% \renewcommand\Re{\mathrm{Re}}
% \renewcommand\Im{\mathrm{Im}}
\newcommand\fn[5]{
    \begin{center}
        \begin{tikzpicture}[node distance=0.5cm and 2cm, auto]
            \node (A) [above of =B] {#2};
            \node (B) {#4};
            \node (E) [right= of B] {#5};
            \node (D) [above of =E] {#3};
            \draw[->] (A) to node {#1}(D);
            \draw[|->] (B) to node {}(E);
        \end{tikzpicture}\\
    \end{center}
}

\newcommand{\pfrac}[2]{\left(\frac{#1}{#2}\right)}

%%%% MATRICES
\newcommand\mx[1]{\begin{bmatrix*}#1\end{bmatrix*}}
\newcommand\px[1]{\begin{pmatrix*}#1\end{pmatrix*}}
\newcommand\Mx[1]{\begin{matrix*}#1\end{matrix*}}
\newcommand\dmx[1]{\begin{pmatrix*}#1\end{pmatrix*}}
\newcommand\Dmx[1]{
    \bordermatrix{#1}}

\newcommand\Gdmx[1]{
    \dmx{
        {#1}_{11} & \cdots & {#1}_{m1}\\
        \vdots & \ddots & \vdots\\
        {#1}_{1n} & \cdots & {#1}_{mn}
    }
}

\newcommand\GDmx[1]{
    \Dmx{~ & q_1 & \cdots & q_m \cr
        L_1 & {#1}_{11} & \cdots & {#1}_{m1}\cr
        \vdots & \vdots & \ddots & \vdots\cr
        L_n & {#1}_{1n} & \cdots & {#1}_{mn}
    }
}

\newcommand{\Suc}[4][1]{{#2}_{#1}^{#4_{#1}}, \ldots, {#2}_{#3}^{#4_{#3}}}
\newcommand{\suc}[3][1]{{#2}_{#1}, \ldots, {#2}_{#3}}
\newcommand{\sucn}[1]{\suc{#1}{n}}
\newcommand{\Sucn}[2]{\Suc{#1}{n}{#2}}
\newcommand{\sucm}[1]{\suc{#1}{m}}
\newcommand{\Sucm}[2]{\Suc{#1}{m}{#2}}

\newcommand{\Mult}[4][1]{{#2}_{#1}^{#4_{#1}} \cdots {#2}_{#3}^{#4_{#3}}}
\newcommand{\mult}[3][1]{{#2}_{#1} \cdots {#2}_{#3}}
\newcommand{\multn}[1]{\mult{#1}{n}}
\newcommand{\Multn}[2]{\Mult{#1}{n}{#2}}
\newcommand{\multm}[1]{\mult{#1}{m}}
\newcommand{\Multm}[2]{\Mult{#1}{m}{#2}}

\renewcommand\qedsymbol{$\diamondsuit$}

\title{Apuntes de Modelizaci\'on}
\author{Rafael S\'{a}nchez S\'{a}nchez}
\begin{document}
\begin{titlepage}
	\includepdf{portada.pdf}
\end{titlepage}


Revisión del \today $ $ a las \currenttime.

\begin{center}

\end{center}

\tableofcontents
\part{Primer parcial}
% !TeX root = ../modelizacion.tex

\chapter{An\'alisis dimensional}

El análisis dimensional es una herramienta que nos permite simplificar el estudio de cualquier fenómeno que involucre varias magnitudes físicas para tratarlas como variables independientes. Esto nos ayudará a simplificar los modelos matemáticos de lo que queramos estudiar.\\

Vamos a comenzar con un ejemplo introductorio a la asignatura, con el que se busca de alguna forma introducir conceptos que si bien no son del todo matemáticos o formales serán de utilidad en el desarrollo del curso.

\begin{eg}[Segunda Ley de Newton - Ley física]\label{eg:1}
    La segunda ley de Newton se puede escribir como la ecuación diferencial:
    $$
        m \ddot x(t) = F(x, t),\ t\in[0, T]
    $$
    donde $m$ representa la masa de un objeto, $x(t)$ la posición del mismo respecto del tiempo, $F(x, t)$ la fuerza que se ejerce sobre él y $T$ es el tiempo final.\\
    Para completar el problema daremos un par de condiciones iniciales:
    \begin{align*}
        x(0) = x_0 \text{ la posición inicial }\\
        \dot x(0) = v_0 \text{ la velocidad inicial}
    \end{align*}

    En el análisis dimensional analizaremos que \textbf{magnitudes} entran en juego en la \textbf{ley}. En este caso tenemos:
    \begin{itemize}
        \item $m$ - masa.
        \item $x$ - posición.
        \item $F$ - fuerza.
        \item $T$ - tiempo final.
        \item $x_0$ - posición inicial.
        \item $v_0$ - velocidad inicial.
    \end{itemize}

    Por tanto nuestra función final será de la forma:
    $$
        f(m, x, F, T, x_0, v_0) = 0
    $$
    que es otra forma de expresar la \textbf{ley}. Además, querremos ver de qué \textbf{magnitudes} dependen estos $6$ parámetros. Esto lo expresaremos con la notación:
    $$
        \mx{p} = M
    $$
    donde $p$ representa un parámetro y $M$ una magnitud (también puede ser un producto de ellas). En nuestro caso tenemos:
    \begin{itemize}
        \item $\mx{m} = M$. Masa, una magnitud elemental.
        \item $\mx{x} = L$. Longitud, una magnitud elemental.
        \item $\mx{T} = \tau$. Tiempo, una magnitud elemental.
        \item $\mx{x_0} = L$. Longitud.
        \item $\mx{v_0} = L \cdot \tau^{-1}$. Velocidad, longitud $\times$ tiempo$^{-1}$.
        \item $\mx{F} = \mx{m \cdot \ddot x} = M \cdot L \cdot \tau^{-2}$. Fuerza, masa $\times$ longitud $\times$ tiempo$^{-2}$.
    \end{itemize}
\end{eg}

\section{Magnitudes. Teorema $\Pi$.}

Vamos a suponer la existencia de $\sucn{L}$ magnitudes elementales con $n < \infty \in \N$, es decir, cada $L_i$ es independiente de cada magnitud de $\mtc L \setminus L_i$. Diremos que una colección de magnitudes conforman un sistema.

\begin{dfn}[Dimensión de una magnitud]
    Sea $a \in \R$ una medida de una magnitud $A$ en un sistema $\sucn{L}$. Si cambiamos a un sistema $\sucn{L'}$ con $L'_i = \lambda_i L_i$ y sea $a'$ la medida de $A$ en el nuevo sistema, entonces si se cumple que:
    $$
        a' = a \cdot \lambda_1^{a_1} \cdot \cdots \cdot \lambda_n^{a_n}
    $$
    para una serie de escalares $\sucn{a}$, entonces diremos que la magnitud $A$ tiene \textbf{dimensión} $L_{1}^{a_1} \cdots L_{n}^{a_n}$ y lo expresamos por:
    $$
        \mx{A} = L_{1}^{a_1} \cdots L_{n}^{a_n}
    $$

    Cuando el sistema $\sucn{L}$ esté fijado podremos identificar la dimensión de la magnitud $A$ con el vector de escalares $(\sucn{a})$.
\end{dfn}

Recordando el \hyperref[eg:1]{ejemplo de la segunda ley de Newton}, donde teníamos tres magnitudes elementales $(L, \tau, M)$, si consideramos que nuestro sistema $\sucn{L}$ es dicha $3$-tupla, entonces podemos expresar las magnitudes no elementales como $3$-tuplas (o vectores de $\R^3$):
\begin{itemize}
    \item $\mx{v_0} = L \cdot \tau^{-1} = (1, -1, 0)$
    \item $\mx{F} = \mx{m \ddot x} = M \cdot L \cdot \tau^{-1} =  (1, -2, 1)$
\end{itemize}

\begin{eg}[Dimensión de una magnitud]
    Sea $L_1 = \sdf{m} $ y $L_2=\sdf{s}$ un sistema de magnitudes (longitud en metros y tiempo en segundos  respectivamente), consideramos la magnitud de la velocidad $V$ que tiene dimensión:
    $$
        \mx{V} = L_1 \cdot L_2^{-1} = (1, -1)
    $$
    entonces, si tenemos una medida $v = 30 \sfrac{m}{s}$ y queremos ver su medida $L'$ en el sistema:
    \begin{align*}
        L_1' &= 10^{-3} L_1 \text{ con $L_1'$ longitud en km }\\
        L_2' &= \frac{1}{3600} L_2 \text{ con $L_2'$ tiempo en h }
    \end{align*}
    entonces:
    $$
        v' = v \cdot \lambda_1 \cdot \lambda_2^{-1} = \frac{30 \cdot 3.6 \cdot \cancel{10^3}}{\cancel{10^3}} = 108 \sfrac{km}{h}
    $$
\end{eg}

\begin{pro}[Expresión de una magnitud dependiente]
    Sean $A, B$ dos magnitudes tales que:
    \begin{align*}
        \mx{A} &= \Multn{L}{a}\\
        \mx{B} &= \Multn{L}{b}
    \end{align*}
    Sea $C$ otra magnitud dependiente de $A$ y $B$, tal que si $a, b$ son medidas de $A, B$ y $C$ y $\exists p, q, d$ tales que $c = d\cdot a^p + b^q$ con $p, q, d$ independientes de las unidades $\sucn{L}$, entonces:
    $$
        \mx{C} = L_1^{a_1p + b_1q} \cdots L_n^{a_n p + b_n q}
    $$
\end{pro}

\begin{proof}
    Sean $L'_i = \lambda_i L_i$ un nuevo sistema, entonces:
    $$
        a' = a \cdot \Multn{\lambda}{a},\ \ b' = b \cdot \Multn{\lambda}{b}
    $$
    y por tanto $c'$ es:
    \begin{align*}
        c' &= da'^p + b'^q = d(a\Multn{\lambda}{a})^p + (b\Multn{\lambda}{b})\\
          &= (da^p + b^q) \cdot (\lambda_1^{a_1p+b_1q} \cdots \lambda_n^{a_np+b_nq})\\
          &=  c \cdot (\lambda_1^{a_1p+b_1q} \cdots \lambda_n^{a_np+b_nq}) \implies\\
          \mx{C} &= L_1^{a_1p + b_1q} \cdots L_n^{a_n p + b_n q}
    \end{align*}
\end{proof}

\begin{dfn}[Matriz de dimensiones]
    Dados $\sucm{q}$ magnitudes, tales que su dimensión es:
    $$
        \mx{q_i} = \Multn{L}{{a_i}}
    $$
    llamamos \textbf{matriz de dimensiones} a la matriz ($n \times m$):
    $$
        \dmx{
            a_{11} & a_{21} & \cdots & a_{m1}\\
            \vdots & \vdots & \ddots & \vdots\\
            a_{1n} & a_{2n} & \cdots & a_{mn}
            }
    $$
    que tiene $n$ filas (una por cada \textit{magnitud elemental} $L_i$) y $m$ columnas (una por cada magnitud del problema $q_i$)
\end{dfn}

\begin{obs}
    Los índices de los elementos de la matriz quedan permutados respecto de la notación habitual.
\end{obs}

Retomando el \hyperref[eg:1]{ejemplo de la segunda ley de Newton} tendríamos la \textit{matriz de dimensiones}:
$$
    \Dmx{~ & x_0 & v_0 & T & F & M & x \cr
        L & 1 & 1 & 0 & 1 & 0 & 1\cr
        \tau & 0 & -1 & 1 & -2 & 0 & 0\cr
        M & 0 & 0 & 0 & 1 & 1 & 0
    }
$$

\begin{dfn}[Magnitud adimensional]
    Una magnitud $\Pi$ se dice adimensional si $\mx{\Pi} = 1$.
\end{dfn}

Hallar magnitudes adimensionales en los distintos problemas nos ayuda a simplificar el estudio de los mismos. Retomemos el \hyperref[eg:1]{ejemplo de la segunda ley de Newton}, vamos a intentar reducir la dimensión del problema.

\begin{eg}[Reduciendo la dimensión del ejemplo \ref{eg:1}]
    Recordemos que teníamos $6$ parámetros ($x, x_0, v_0, T, F, m$). Una forma de reducir los parámetros es intentar \textit{enmascarar} los valores iniciales en nuevas variables.\\ Recordemos que tanto $x$ como $x_0$ tenían la misma dimensión. Gracias a ello podemos definir un nuevo parámetro $y$ sin dimensión:
    $$
        y = \frac{x}{x_0}
    $$
    Además, como $x(0) = x_0$ tendremos que $y(0) = 1$ y como $\dot x(0) = v_0$ entonces $\dot y(0) = \frac{v_0}{x_0}$.\\
    Podemos también hacer lo mismo con el tiempo, recordemos que en la fórmula original la variable $t$ pertencía a $[0, T]$. Podemos definir entonces:
    $$
        \frac{t}{T} = \tau
    $$
    y por tanto $\tau \in [0, 1]$, con lo que hemos \textit{eliminado} $T$. Sin embargo, cambiar $t$ tiene consecuencias debido que es la variable respecto de la que se diferencia $x$ (y por tanto $y$), tenemos que ver como afectan estos cambios a nuestras variables.\\
    Usaremos la notación $\dot x$ para referirnos a $\pd{x}{t}$ y $x'$ para referirnos a $\pd{x}{\tau}$.\\
    Entonces obtenemos:
    $$
        y' = \pd{y}{\tau} = \frac{1}{x_0} \pd{x}{\tau} = \frac{T}{x_0} \pd{x}{t} = \frac{T}{x_0} \dot x \implies y'(0) = \frac{T}{x_0} \cdot v_0 = \tilde q
    $$
    y de nuevo $\mx{\tilde q} = 1$.\\\\

    Recordemos que nuestro problema comenzaba con $m\ddot x = F$, vamos a usar esto para encontrar otro parámetro adimensional.
    \begin{align*}
        &\pd{x}{t} = \frac{x_0}{T} \pd{y}{\tau}\\
        &\pdn{x}{t}{2} = \frac{x_0}{T} \Pd{t} \left(\pd{y}{\tau}\right) = \frac{x_0}{T} \pd{\tau}{t} \pdn{y}{\tau}{2} = \frac{x_0}{T} \cdot y''\\
        & \frac{mx_0}{T^2} \cdot y'' = F \implies y'' = \frac{T^2}{mx_0}\cdot F = f
    \end{align*}
    y podemos comprobar que $\mx{f} = 1$. Recapitulado, hemos conseguido encontrar nuevos parámetros adimensionales $y$, $\tilde q = y'$, $f = y''$ haciendo algunos cambios en el problema. Con esto, podemos reescribir el problema de valores iniciales con los nuevos parámetros adimensionales:
    \begin{align*}
        &y'' = f = \frac{T^2}{mx_0}\cdot F\\
        &y(0) = 1\\
        &y'(0) = \tilde q = \frac{T v_0}{x_0}
    \end{align*}
\end{eg}

\begin{dfn}[Ley invariante]
    Sea una ley $f(\sucm{q}) = 0$, se dice que es \textbf{invariante} frente al cambio de unidades $L'_1 = \lambda_1 L_1, \ldots L'_n = \lambda_n L_n$ si verifica que $f(\sucm{q'}) = 0$ para $\sucm{q'}$ las medidas de $\sucm{q}$ en las nuevas unidades $\sucn{L'}$. Informalmente:\\
    \begin{center}
        Una ley es invariante cuando sigue siendo cierta tras el cambio de variables del problema
    \end{center}
\end{dfn}

\begin{thm}[Teorema $\Pi$]\label{thm:pi}
    Sea $f(\sucm{q}) = 0$ una ley invariante con $\sucm{q}$ magnitudes con matriz de dimensiones:
    $$
        D = \dmx{
            a_{11} & \cdots & a_{m1}\\
            \vdots & \ddots & \vdots\\
            a_{1n} & \cdots & a_{mn}
        }
    $$
    tal que $n < m$ y el rango de $D$ es $r \leq n$. Entonces existen $m - r$ cantidades $\suc{\Pi}{m-r}$ que van a ser magnitudes adimensionales tales que la ley invariante es equivalente a una relación $F(\suc{\Pi}{m-r}) = 0$.
\end{thm}

\begin{proof}$ $
    \begin{enumerate}[(1)]
        \item Vamos a demostrar que existen $\suc{\Pi}{m-r}$ magnitudes adimensionales independientes entre sí. Partimos de la matriz de dimensiones:
        $$
            \GDmx{a}
        $$
        Y entonces tenemos:
        $$\mx{\Pi} = 1 \text{ con } \Pi = \Multm{q}{\alpha}$$
        $$
            \mx{\Pi} = (\Multn{L}{{a_1}})^{\alpha_1} \cdot (\Multn{L}{{a_2}})^{\alpha_2} \cdots (\Multn{L}{{a_m}})^{\alpha_m} = L_1^{\alpha_1 a_{11} + \cdot + \alpha_m a_{m1}} \cdots L_n^{\alpha_1 a_{1n} + \cdots \alpha_m a_{m1}}
        $$
        De donde surge el sistema:
        \begin{align*}
            \alpha_1 a_{11} + \alpha_2 a_{21} + \cdots + \alpha_m a_{m1} &= 0\\
            \alpha_1 a_{12} + \alpha_2 a_{22} + \cdots + \alpha_m a_{m2} &= 0\\
            (\cdots)\\
            \alpha_1 a_{1n} + \alpha_2 a_{2n} + \cdots + \alpha_m a_{mn} &= 0\\
        \end{align*}
        Y la demostración se reduce a resolver un sistema homogéneo de $n$ ecuaciones con $m$ incógnitas. Como la matriz $D$ es de rango $r$, entonces existen $m-r$ soluciones linealmente independientes.

        \item La demostración de que la ley invariante es equivalente a otra que solo comprende a las magnitudes adimensionales para el caso general resulta muy difícil de escribir. Un caso particular para $m=4, n=2$ y $r=2$ se encuentra en el siguiente ejemplo.
    \end{enumerate}
\end{proof}

\begin{eg}[Cálculo de las magnitudes adimensionales y la ley invariante asociada]
    Vamos a hacer un caso particular del sistema de la demostración. Sea $m=4$, $n=2$ y $r=2$, es decir:
    \begin{align*}
        \alpha_1 a_{11} + \alpha_2 a_{21} + \alpha_3 a_{31} + \alpha_4 a_{41} &= 0\\
        \alpha_1 a_{12} + \alpha_2 a_{22} + \alpha_3 a_{32} + \alpha_4 a_{42} &= 0\\
    \end{align*}
    Como el rango es $2$, sin perdida de generalidad puedo suponer que:
    $$
        \det \mx{a_{11} & a_{21}\\
                 a_{12} & a_{22}} = 0
    $$
    Entonces por Rouche-Frobenius existen $C_{34}, C_{32}, C_{41}, C_{42}$ tales que:
    \begin{align*}
        \alpha_1 = - \alpha_3 C_{31} - \alpha_4 C_{41} &\text{ con } \alpha_3 = 1,\ \alpha_4=0\\
        \alpha_2 = - \alpha_3 C_{32} - \alpha_4 C_{42} &\text{ con } \alpha_3 = 0,\ \alpha_4=1
    \end{align*}

    Y entonces puedo encontrar las magnitudes adimensionales:
    \begin{align*}
        \Pi_1 = q_1^{-C_{31}} q_2^{-C_{32}} q_3^{1} q_4^{0} \implies q_3 = \Pi_1 q_1^{C_{31}} q_2^{C_{32}}\\
        \Pi_2 = q_1^{-C_{41}} q_2^{-C_{42}} q_3^{0} q_4^{1} \implies q_4 = \Pi_2 q_1^{C_{41}} q_2^{C_{42}}
    \end{align*}

    Recordemos que partíamos de una ley:
    $$
        f(q_1, q_2, q_3, q_4) = 0
    $$
    Y además, hemos encontrado que:
    $$
        f(q_1, q_2, q_3, q_4) = f(q_1, q_2, \Pi_1 q_1^{C_{31}} q_2^{C_{32}}, \Pi_2 q_1^{C_{41}} q_2^{C_{42}}) = G(q_1, q_2, \Pi_1, \Pi_2)
    $$
    Además, vamos a hacer un cambio del sistema de magnitudes de la forma:
    $$
        L_1' = \lambda_1 L_1,\ L_2' = \lambda_2 L_2
    $$
    para conseguir:
    \begin{align*}
        q_1' &= q_1 \lambda_1^{a_11} \lambda_2^{a_12}\\
        q_2' &= q_2 \lambda_1^{a_21} \lambda_2^{a_22}\\
        \Pi'_1 &= \Pi_1\\
        \Pi'_2 &= \Pi_2
    \end{align*}
    Entonces tenemos una nueva ley invariante al cambio de escala:
    $$
        0 = G(q_1, q_2, \Pi_1, \Pi_2) = G(q_1', q_2', \Pi_1, \Pi_2)
    $$
    Finalmente, querremos hacer un cambio de variables tal que $q_1 = q_2 = 1$ en nuestra nueva ley.
    \begin{align*}
        L_nq_1 + a_{11} L_n \lambda_1 + a_{12} L_n + \lambda_2 &= 0\\
        L_nq_2 + a_{21} L_n \lambda_1 + a_{22} L_n + \lambda_2 &= 0
    \end{align*}
    En el sistema anterior hacemos el cambio: $y_i = L_n \lambda_i$ y tomando logaritmos obtenemos:
    \begin{align*}
        a_{11} y_1 + a_{12} y_2 = - \log(q_1)\\
        a_{21} y_1 + a_{22} y_2 = - \log(q_2)
    \end{align*}
    Y sabemos que dichos $q_1, q_2$ existen porque la matriz es invertible, y por tanto:
    $$
        0 = f(q_1, q_2, q_3, q_4) = G(1, 1, \Pi_1, \Pi_2) = F(\Pi_1, \Pi_2)
    $$
\end{eg}

\begin{eg}[H1.4]$ $
    \begin{enumerate}[a)]
        \item Magnitudes. Magnitudes elementales. Magnitudes adimensionales.\\\\
        Nuestras magnitudes $q_i$ serán ($h, t, g, R, v$). Podemos analizarlas en función de las magnitudes elementales:
        \begin{itemize}
            \item $\mx{h} = L = $ \{longitud\}.
            \item $\mx{t} = T = $ \{tiempo\}.
            \item $\mx{g} = L T^{-2}$.
            \item $\mx{v} = L T^{-1}$.
            \item $\mx{R} = L$.
        \end{itemize}
        Que tiene como matriz de dimensiones:
        $$
            \dmx{
                1 & 0 &  1 &  1 & 1\\
                0 & 1 & -2 & -1 & 0
            }
        $$
        Donde podemos encontrar las siguientes $3$ magnitudes adimensionales (ya que el rango de la matriz es $2$):
        $$
            \frac{h}{R},\ \frac{tv}{R},\ \frac{v}{\sqrt{gR}}
        $$
        \item Relación de la altura máxima a alcanzar el proyectil con respecto a $v$, $g$, $R$.\\\\
        Partimos de $0 = f(h, t, v, g, R)$ y por el \hyperref[thm:pi]{teorema \Pi} sabemos que tenemos una ley equivalente:
        $$
            0 = F\left(\frac{h}{R},\ \frac{tv}{R},\ \frac{v}{\sqrt{gR}}\right)
        $$
        Además, sabemos que la altura máxima $h_{max}$ se va a tomar en un tiempo determinado $t_{max}$. Por tanto tenemos la relación:
        $$
            0 = F\left(\frac{h_{max}}{R},\ \frac{t_{max}v}{R},\ \frac{v}{\sqrt{gR}}\right)
        $$
        Por el Teorema de la función implícita (TFI):
        $$
            \frac{h}{R} = G\left( \frac{tV}{R}, \frac{V}{\sqrt{gR}} \right)
        $$
        Como alcanzamos la altura máxima:
        $$
            h'(t_{max}) = 0 \implies \pd{G}{\Pi_2}\left( \frac{t_{max}V}{R}, \frac{V}{\sqrt{gR}} \right) = 0
        $$
        Aplicando de nuevo el TFI:
        $$
            \frac{t_{max} v}{R} = \varphi\left( \frac{V}{\sqrt{gR}} \right)
        $$
        Y llegamos a la relación que nos pedían:
        $$
            \frac{h_{max}}{R} = G\left( \varphi\left( \frac{V}{\sqrt{gR}} \right), \frac{V}{\sqrt{gR}} \right) = \varphi^{\star}\left( \frac{V}{\sqrt{gR}} \right)
        $$
    \end{enumerate}
\end{eg}

% !TeX root = ../modelizacion.tex

\chapter{Modelos matriciales discretos}

\section{Modelo discreto unidimensional}

En esta sección trataremos problemas con ecuaciones en diferencias finitas. Podemos dar un ejemplo con el interés de una tarjeta de crédito.

\begin{eg}[Interés en una tarjeta de crédito]
    Tomamos $x_0$ la cantidad que debemos, e $i = 2\%$ el tipo de interés.\\
    $x(t)$ es la cantidad de dinero que debo cuando han pasado $t$ días. Vamos a intentar dar una expresión para $x$.
    \begin{align*}
        x(0) &= x_0\\
        x(1) &= (1 + i)\cdot x(0) = x(0) + i \cdot x(0)\\
        x(2) &= (1+i) \cdot x(1) = (1+i)^2\cdot x(0)
    \end{align*}
    y por inducción hallamos:
    $$
        x(k) = (1+i)^{k} \cdot x(0)
    $$
    Por tanto obtenemos el sistema:
    \begin{align*}
        x(k+1) &= ax(k)\\
        x(0) &= x_0
    \end{align*}
    donde $a = (1+i)$, y estamos ante un sistema de ecuaciones en diferenciadas finitas homogéneas y autónomas.
\end{eg}

% \part{Segundo parcial}
% % !TeX root = ../galois.tex

\chapter{Extensiones de Galois}

\section{Cuerpos de escisión}

\begin{dfn}[Escisión de un polinomio]
    Sea $f \in K[x]$, $K$ cuerpo, decimos que $f$ se \textbf{escinde} si $\exists a, a_i\in K$ si:
    $$
        f(x) = a(x - a_1) \ldots (x - a_n)
    $$
\end{dfn}

\begin{obs}
    Si $f(x) \in K[x]$, $\delta f = 1$, entonces $f$ se escinde en $K$.
    $$
        f(x) = ax + b = a \left(x - \left(- \frac{b}{a}\right)\right),\ \ \ a, -\frac{b}{a} \in K
    $$
\end{obs}

\begin{obs}
    Si $f \in K[x]$, se escinde en $E$, sea $E/K$ una extensión:
    $$
        f(x) = a (x - a_1) \ldots (x - a_n) \in E[x]
    $$
    Entonces las raíces de $f$ en cualquier extensión de $E$ son $a_1, \ldots, a_n$
\end{obs}

\begin{dfn}[Cuerpo de escisión]
Sea $f \in K[x]$, entonces $E$ es un cuerpo de escisión de $f$ sobre $K$ si:
    \begin{itemize}
        \item $E/K$ es una extensión.
        \item $f$ se escinde en $E$.
        \item Si $f$ se escinde en $L$, con $L$ un cuerpo intermedio $K \subseteq L \subseteq E$, entonces $L=E$, es decir, $E = K(f)$. (Es el menor subcuerpo con la propiedad de que $f$ se escinda en él).
    \end{itemize}
\end{dfn}

\begin{eg}
    $x^2 + 1$ se escinde en $\C$ pero su cuerpo de escisión es $\Q(i)$.
\end{eg}

\begin{obs}
    Si $f \in K[x]$ se escinde en $E$, entonces $E$ tiene todas las raíces de $f$. Para construir su cuerpo de escisión basta adjuntar todas las raíces al cuerpo sobre el que está definido:
    $$
        f(x) = a (x - a_1) \ldots (x - a_n) \in E[x], \text{ entonces su cuerpo de escisión es } K(f) = K(a_1, \ldots, a_n)
    $$
\end{obs}

\begin{lm}[Escisión de polinomios no constantes]\label{lm:3.1}
    Sea $K$ un cuerpo:\\
    \begin{enumerate}
        \item[($i$)] Si $f \in K[x]$ no constante se escinde en $K$ si y sólo si todos los polinomios en la descomposición en términos de factores irreducibles tienen grado $1$.
        \item[($ii$)] Si $f \in K[x]$ no constante que se escinde en $K$ y sea $p \in K[x]$ tal que $p \divides f$, entonces $p$ se escinde en $K$.
    \end{enumerate}
\end{lm}

\begin{proof}$ $
    \begin{enumerate}
        \item[($i$)] Como $f \in K[x]$ no constante suponemos $\delta(f) = n \geq 1$, lo que quiere decir que $f$ tiene $n$ raíces. Entonces:
        \begin{itemize}
            \item[$\implies$] Si $f$ se escinde en $K$, sea $\Omega = \sdf{\omega \in K \mid f(\omega) = 0, \forall \omega \in \Omega}$ el conjunto de raíces de $f$, sabemos que $\Omega \subseteq K$, por tanto, por Ruffini podemos expresar $f$ como $\prod_i^n (x-\omega_i)^{m(i)},\ \forall \omega_i \in \Omega$ (donde $m(i)$ es la multiplicidad de cada raíz $\omega_i$) y entonces todos los términos de la descomposición tienen grado $1$.
            \item[$\Longleftarrow$] Partimos de que $f$ se puede expresar en $K[x]$ como $\prod_i^n (x-\omega_i)^{m(i)}$ con $\omega_i$ como en el apartado anterior. Entonces $\omega_i \in K[x]\ \forall i \implies f$ se escinde en $K$.
        \end{itemize}
        \item[($ii$)] Ya hemos visto que si $f$ se escinde se puede expresar como producto de polinomios de grado $1$. Si $p \divides f$ entonces sea $\delta(p) = l \leq n$, $p$ es $\prod_i^l (x-\omega_i)^{m(i)}, \forall \omega_i \in \Omega' \subseteq \Omega$. Por el resultado anterior, como es producto de polinomios de grado $1$, se escinde en $K$.
    \end{enumerate}
\end{proof}

\begin{lm}[Cuerpos de escisión y cuerpos intermedios]\label{lm:3.2}
    Sean $E/L$ y $L/K$ extensiones, y $f \in K[x]$ un polinomio no constante. Si $E$ es el cuerpo de escisión de $f$ sobre $K$, entonces $E$ es el cuerpo de escisión de $f$ sobre $L$.
\end{lm}

\begin{proof}
    Como $E = K(f)$, tenemos que $f(x) = a (x - a_1) \ldots (x - a_n) \in E[x]$ y $E = K(a_1, \ldots, a_n)$. Además, $f \in K[x] \subseteq L[x]$. Si el cuerpo de escisión de $f$ sobre $L$ es $E'$, fácilmente llegamos a que $E' = L(f) \subseteq E$ y por tanto: $E = E'$ ya que $E$ era el menor cuerpo de escisión de $f$ sobre $K$.
\end{proof}

Si $f$ es un polinomio sobre $\Q[x]$ de grado $n$, \textit{sabemos} que $\exists \alpha_i \in \C$ con $i \in \setdef{1, \ldots n}$ tal que $f(x) = a (x - \alpha_1) \ldots (x - \alpha_n)$, y hemos visto para calcular su cuerpo de escisión basta con adjuntar las raíces de $f$, es decir, $\Q(f) = \Q(\alpha_1, \ldots, \alpha_n)$.\\

Vamos a ver resultados para generalizar este proceso.

\begin{lm}[Teorema de Kronecker]\label{lm:3.3}
    Sea $p \in K[x]$ un polinomio irreducible, entonces existe $E/K$ tal que $p$ tiene una raíz en $E$.
\end{lm}

\begin{proof}
    Sea $L = \quo{K[x]}{\gen{p}}$, ya vimos que $L$ es un cuerpo por ser $p$ irreducible. Además se puede comprobar que:
    $$
        \bar x = x + \gen{p} \in L \text{ es raíz de } \bar p \in \bar K[y]
    $$
    Donde $\bar p = \overline{a_0} + \ldots + \overline{a_n}x^n$ y $K \isom \bar K = \setdef{\bar K \mid k \in K}$.

    Hay que comprobar que $\bar K \subseteq L$. Tras ello, habría que ver que $L = \bar K(\bar x)$.
\end{proof}

\begin{thm}[Existencia de cuerpos de escisión]\label{thm:3.4}
    Sea $K$ un cuerpo, $f \in K[x]$ un polinomio no constante. Entonces existe un cuerpo de escisión $E$ de $f$ sobre $K$.
\end{thm}

\begin{proof}
    Basta construir una extensión en la que $f$ se escinda. La demostración sigue por inducción sobre $\delta(f)$.
    \begin{itemize}
        \item[$\delta(f) = 1$] Entonces por el lema \ref{lm:3.1}, $K$ es un cuerpo de escisión de $f$ sobre $K$.
        \item[$\delta(f)>1$] Sea $p \divides f$ un factor irreducible de $f$ en $K[x]$, por el lema \ref{lm:3.3}, sabemos que existe $E/K$ en la que $p$ tiene una raíz $\alpha \in E$, en particular $f(\alpha) = 0$ y por Ruffini:
        $$
            f(x) = (x - \alpha) g(x),\text{ con } g(x) \in K(\alpha)[x]
        $$
        Como $1 \leq \delta(g) \leq \delta(f) - 1 \leq \delta(f)$, por inducción sabemos que existe un cuerpo de escisión $L$ de $g$ sobre $K(\alpha)$. Si $\alpha_1, \ldots \alpha_n \in L$ son las raíces de $g$, entonces:
        $$
            L = K(\alpha)(\alpha_1, \ldots \alpha_n) = K(\alpha, \alpha_1, \ldots, \alpha_n)
        $$
        y por tanto $f(x) = (x - \alpha) g(x) = (x - \alpha) p (x - \alpha_1) \ldots (x - \alpha_n) \in L[x]$, así que $f$ se escinde sobre $L$, y de hecho $L$ es un cuerpo de escisión de $f$ sobre $K$.
    \end{itemize}
\end{proof}

\begin{ex}[Cálculo del cuerpo de escisión]
    Describir el cuerpo de escisión de $f(x) = x^4 - 4x^2 + 2$ sobre $\Q$.\\\\
    Comenzamos hallando las raíces, que resulta fácil pues $f(x) = 0$ es una ecuación bicuadrática. Tenemos entonces las raíces:
    $$
        \alpha = \sqrt{2 + \sqrt 2},\ \beta = \sqrt{2 - \sqrt{2}},\ -\alpha,\ -\beta
    $$
    De donde vemos que $f$ se escinde en $R$ ya que todas las raíces son reales. Para ver su cuerpo de escisión vamos a adjuntar las raíces a $\Q$.
    $$
        E = \Q(\alpha,  \beta, -\alpha, -\beta) = \Q(\alpha, \beta)
    $$
    Sin embargo, podemos reducirlo aún más si vemos que $\Q(\alpha, \beta) = \Q(\alpha)$.
    $$
        \alpha\beta = \sqrt{(2 + \sqrt 2)(2 - \sqrt 2)} = \sqrt 2,\ \alpha^{-1} = \left(\frac{\beta}{\sqrt 2}\right) \in \Q(\alpha) \implies \beta = \frac{\beta}{\sqrt 2} \cdot \sqrt 2 \in \Q(\alpha)
    $$
    Con lo que el cuerpo de escisión es: $\Q(\sqrt{2 + \sqrt 2})$.
\end{ex}

\begin{obs}
    Si $E = K(f)$, entonces $E/K$ es finita por el teorema \ref{thm:extalg}. Anteriormente, vimos que si $f \in K[x]$, entonces existe un cuerpo de escisión de $f$ sobre $K$.\\
    Sea $p$ un polinomio irreducible tal que $p \divides f$ y $\delta p > 1$, entonces consideramos $\quo{K[x]}{\gen{p}}$ y en ese cuerpo hay una raíz de $p$ (que también es raíz de $f$), y después aplicamos inducción.
\end{obs}

\begin{thm}\label{thm:3.5}
    Sea $K_1, K_2$ cuerpos, y $\sigma: K_1 \to K_2$ un isomorfismo de cuerpos. Sea $f_1 \in K_1[x]$ no constante y $\sigma(f_1) = f_2$.\\ Supongamos que $E_i$ son cuerpos de escisión de $f_1$ sobre $K_i$ con $i \in \setdef{1, 2}$. Entonces existe un isomorfismo $\tau: E_1 \to E_2$ que extiende a $\sigma$, es decir, $\left. \tau \right|_{K_1} = \sigma$.
\end{thm}

\begin{proof}
    Por inducción sobre $|E_1:K_1| < \infty$.\\\\
    \begin{itemize}
        \item[$|E_1:K_1|=1$] En este caso, $\vabs{E_1:K_1} \iff f_1$ se escinde en $K_1[x] \implies \sigma(f_1) = f_2$ se escinde en $K_2[x] \implies E_2 = K_2$ y podemos tomar $\tau = \sigma$.
        \item[$\vabs{E_1:K_1}>1$] En este caso $f_1$ no se escinde en $K_1[x]$, y podemos escoger $p$ un factor de $f_1$ irreducible con $\delta(p) > 1$. Además, sea $q = \sigma(p)$ tal que $q \divides f_2$ y es irreducible con $\delta(q) > 1$.
        Escogemos las raíces de $p$ en $E_1(a)$ y las de $q$ en $E_2(b)$ (por el lema \ref{lm:3.1} sabemos que $p$ y $q$ se escinden en $E_1$ y $E_2$ respectivamente).
        \begin{center}
            \begin{tikzpicture}[node distance=1cm and 2cm, auto]
                \node (A) {$\tau: E_1 = K_1(f_1)$};
                \node (B) [right= of A] {$E_2 = K_2(f_2)$};
                \node (C) [below of =A] {$\rest{\tau}_{K_1(a)} = \theta: K_1(a)$};
                \node (D) [below of =B] {$K_2(b)$};
                \node (E) [below of =C] {$\rest{\theta}_{K_1} = \sigma: K_1$};
                \node (F) [below of =D] {$K_2$};
                \node (p) [below of =E] {$p$};
                \node (q) [below of =F] {$q$};
                \draw[->] (C) to node {$\isom$}(D);
                \draw[->] (E) to node {$\isom$}(F);
                \draw[|->] (p) to node {}(q);
                \draw[-] (A) to node {}(C);
                \draw[-] (B) to node {}(D);
                \draw[-] (C) to node {}(E);
                \draw[-] (D) to node {}(F);
                \draw[-] (E) to node {}(p);
                \draw[-] (F) to node {}(q);
            \end{tikzpicture}\\
        \end{center}
        Por el teorema \ref{thm:2.5} existe un isomorfismo $\theta: K_1(a) \to K_2(b)$ que extiende a $\sigma$.\\\\
        Por el teorema \ref{thm:2.3} $\vabs{K_1(a):K_1} = o(p) > 1$, luego por el teorema \ref{thm:2.1} $\vabs{E_1:K_1(a)} < \vabs{E_1:K_1}$ y además $E_1 = K_1(a)(f_1)$ (por el lema \ref{lm:3.2}). Del mismo modo $E_2$ es un cuerpo de escisión de $f_2$ sobre $K_2(b)$.
        Por inducción, existe un isomorfismo $\tau: E_1 \to E_2$ tal que $\rest{\tau}_{K_1(a)} = \theta$. En particular, $\rest{\tau}_{K_1} = \rest{\left( \rest{\tau}_{K_1(a)} \right)}_{K_1} = \rest{\theta}_{K_1} = \sigma$.
    \end{itemize}
\end{proof}

\begin{cor}[Unicidad de los cuerpos de escisión]
    Sea $K$ un cuerpo y $f \in K[x]$ no constante. Si $E_1$ y $E_2$ son cuerpos de escisión de $f$ sobre $K$ entonces $\exists \tau$ un isomorfismo $\tau: E_1 \to E_2$ tal que $\tau(a) = a$ $\forall a \in K$.
\end{cor}
\begin{proof}
    Basta tomar $\sigma = id_K$ en el teorema \ref{thm:3.5}.
\end{proof}

\begin{obs}
    Un cuerpo $K$ en el que todo polinomio $f \in K[x]$ se escinde se dice que es \textit{algebraico cerrado}. Por ejemplo $\C$ es un cuerpo algebraico cerrado.
\end{obs}

\section{Extensiones normales}

\begin{dfn}[Extensión normal]
    Sea $K$ un cuerpo y $E/K$ una extensión, diremos que $E/K$ es una \textbf{extensión normal} si $E$ es el cuerpo de escisión de algún polinomio $f \in K[x]$.
\end{dfn}

\begin{obs}
    Como consecuencia del teorema \ref{thm:2.2} y del teorema \ref{thm:2.4}, si $E/K$ es normal $\implies E/K$ es finita $\implies E/K$ es algebraica.
\end{obs}

\begin{ex}[H3.1]
    Construye los cuerpos de escisión $\Q(x^3-1)$, $\Q(x^4 + 5x^2 + 5)$, $\Q(x^6-8)$ y calcula los grados de las extensiones que resultan.\\\\

    \textbf{Solución}\\
    \begin{itemize}
        \item[$f(x) = x^3-1$] Las raíces son las raíces cúbicas de la unidad: $1, \omega, \bar\omega = \omega^2$, por lo que $\Q(f) = \Q(1, \omega, \omega^2) = \Q(\omega)$.
        \item[$f(x) = x^4+5x^2+5$] Muy parecido a un ejemplo ya resuelto.
        \item[$f(x) = x^6 - 8$] Las raíces son las complejas de $8$.
        $$
            \sqrt[6] 8 = (8^\frac{1}{3})^\frac{1}{2} = \sqrt 2
        $$
        Sea $\xi = e^\frac{2 \pi i}{6} = e^\frac{\pi i}{3}$, las raíces de $f$ son:
        $$
            \sdf{\sqrt 2, \sqrt 2 \xi, \sqrt 2 \xi^2, \sqrt 2 \xi^3, \sqrt 2 \xi^4, \sqrt 2 \xi^5}
        $$
        Por tanto $\Q(f) = \Q(\sqrt 2, \sqrt 2 \xi, \sqrt 2 \xi^2, \sqrt 2 \xi^3, \sqrt 2 \xi^4, \sqrt 2 \xi^5) $.Faltaría comprobar que esta extensión es la misma que $\Q(\sqrt 2 \xi)$. Procediendo de forma similar a ejercicios anteriores, para hallar el grado de la extensión calculamos el polinomio $Irr(\Q, \sqrt 2 \xi)$ y vemos que es $x^4 + 2x^2 + 4$ con lo que el grado de la extensión es $4$.
    \end{itemize}
\end{ex}

% \begin{ex}[H3.5]
%     Decide si las extensiones $\Q(\sqrt 5 i)/\Q$, $\Q(\sqrt 5)/\Q$, $\Q(\sqrt[4] 5)/\Q$ son normales.\\\\
%
%     \textbf{Solución}\\
%     \begin{itemize}
%         \item[$\Q(\sqrt 5 i)/\Q$] Es fácil comprobar que $\Q(\sqrt 5 i) = \Q(x^2 + 5)$ y por tanto son normales.
%         \item[$\Q(\sqrt 5)/\Q$] Es fácil comprobar que $\Q(\sqrt 5) = \Q(x^2 - 5)$ y por tanto son normales.
%         \item[$\Q(5^{\frac{1}{4}})/\Q$] $x^4 - 5$ tiene una raíz en $\Q(\sqrt[4]{5})$ pero no se escinde, por tanto no es normal.
%     \end{itemize}
% \end{ex}

%%%%%%%%%%%%%%%%%%%%%%%%%%%%%%%%%%%%%%%%%%%%%%%%%%%%%%%%%%%%%%%%%%%%%%%%%% 24/10

\begin{lm}[Normalidad de extensiones en cuerpos intermedios]\label{lm:3.7}
    $K \subseteq L \subseteq E$ extensiones, entonces:
    $$
        E/K \text{ normal } \implies E/L \text{ normal.}
    $$
\end{lm}

\begin{proof}
    Como $E/K$ es normal $\implies \exists f \in K[x]$ tal que $E$ es el cuerpo de escisión de $f$ sobre $K$. En particular, $f \in K[x] \subseteq L[x]$. Por el lema \ref{lm:3.2}, $E$ también es el cuerpo de escisión de $f$ sobre $L$, por lo que $E/L$ es normal.
\end{proof}

\begin{cor}[al teorema \ref{thm:3.5}]\label{cor:3.8}
    Sea $E/K$ una extensión normal, $K \subseteq M_1, M_2 \subseteq E$ con $M_1, M_2$ cuerpos intermedios.\\\\
    Si existe un isomorfismo $\sigma: M_1 \to M_2$ que fija a $K$, es decir, $\sigma(a) = a,\ \forall a \in K$, entonces $\sigma$ se extiende a un isomorfismo $\tau: E \to E$ ($\left. \tau \right|_{M_1} = \sigma$).
\end{cor}
\begin{proof}
    Sea $E = K(f)$ con $f \in K[x]$ no constante. En particular por el lema \ref{lm:3.2}, $E = M_1(f) = M_2(f)$. Como $\sigma(f) = f$, el resultado se sigue del teorema \ref{thm:3.5}.
\end{proof}
\begin{obs}
    La hipótesis $\sigma(a) = a,\ \forall a \in K$ es necesaria. La idea es que sin esta hipótesis:
    $$
        E = M_1(f) = M_2(f)
    $$
    pero no sabríamos si $E = M_2(\sigma(f))$.
\end{obs}

\begin{thm}[Condición necesaria y suficiente para la normalidad de una extensión]\label{thm:3.9}
    Sea $E/K$ una extensión finita, $E/K$ es normal si y sólo si todo polinomio $p \in K[x]$ irreducible con una raíz en $E$ se escinde en $E$.
\end{thm}

\begin{proof}$ $
    \begin{itemize}
        \item[$\ \Longleftarrow\ $] Sabemos que $E/K$ es finita, digamos $|E:K| = n$, podemos tomar una $K$-base $\sdf{a_1, \ldots, a_n}$. En particular $E=K(a_1, \ldots, a_n)$.\\
        Como $E/K$ es finita, por el teorema \ref{thm:2.2}, los $a_i$ en son algebraicos sobre $K$. Sea $p_i = Irr(k, a_i) \in K[x]$, por hipótesis cada $p_i$ se escinde en $E$.\\

        Tomamos $f = \prod_{i=1}^{n} p_i \in K[x]$ se escinde en $E$. De hecho, $E = K(f)$. De esta forma hemos probado que:
        $$
            K(f) \subseteq E = K(a_1, \ldots, a_n) \subseteq K(f)
        $$
        \item[$\implies$] Ahora partimos de que $E/K$ es finita y normal. Entonces $\exists f \in K[x]$ tal que $E = K(f)$.\\
        De hecho, si $f(x) = c(x - b_1) \ldots (x-b_n) \in E[x]$, $E = K(b_1, \ldots, b_n)$.\\
        Sea $p\in K[x]$ irreducible con una raíz $a\in E$, por Ruffini (corolario \ref{cor:ruffini}) $p(x) = (x-a)g(x)\in E[x]$. Además, por el teorema \ref{thm:3.4}, existe $M/E$ un cuerpo de escisión de $p$ sobre $E$, por tanto:
        $$
            p(x) = d(x-d_1) \cdot \ldots \cdot (x-d_r), \text{ podemos suponer que $a = d_1$}
        $$
        Dado $d_i$ con $i > 1$, queremos ver que $d_i \in E$. Llamemos $b = d_i$ para algún $i$. Como $a$ y $b$ son raíces de $p \in K[x]$ irreducible, por el teorema \ref{thm:2.5} $\exists \sigma: K(a) \to K(b)$ tal que $\sigma(k) = k,\ \forall k \in K$.

        \begin{center}
            \begin{tikzpicture}[node distance=1cm and 2cm, auto]
                \node (A) {$\sigma: K(a)$};
                \node (B) [right= of A] {$K(b)$};
                \node (C) [below of =A] {$id: K$};
                \node (D) [below of =B] {$K$};
                \draw[->] (A) to node {}(B);
                \draw[->] (C) to node {}(D);
                \draw[<->] (A) to node {}(C);
                \draw[<->] (B) to node {}(D);
            \end{tikzpicture}\\
        \end{center}

        Como $K(a) \subseteq E$, en particular $E = K(a)(f)$. Por otro lado $K(b)(f) = E(b)$. Las raíces de $f$ en $E(b)$ son $b_1, \ldots, b_m$, así que:
        $$
            K(b)(f) = K(b)(b_1, \ldots, b_m) = K(b_1, \ldots, b_m) = E(b)
        $$
        Por el teorema \ref{thm:3.5} existe un isomorfismo $\tau: E \to E(b)$ que extiende a $\sigma$.

        \begin{center}
            \begin{tikzpicture}[node distance=1cm and 2cm, auto]
                \node (E) [above of =A]{$\tau: E = K(a)(f)$};
                \node (E2)[above of =B]{$E(b) = K(b)(f)$};
                \node (A) {$\sigma: K(a)$};
                \node (B) [right= of A] {$K(b)$};
                \node (C) [below of =A] {$id: K$};
                \node (D) [below of =B] {$K$};
                \draw[->] (A) to node {}(B);
                \draw[->] (C) to node {}(D);
                \draw[<->] (A) to node {}(C);
                \draw[<->] (B) to node {}(D);
                \draw[-] (A) to node {}(E);
                \draw[-] (B) to node {}(E2);
                \draw[->] (E) to node {}(E2);
            \end{tikzpicture}\\
        \end{center}

        Por el ejercicio $H2.15$, $|E:K| = |E(b):K|$ y $|K(a):K| = |K(b):K|$. Además $|E:K(a)| = |E(b):K(b)|$.
        Por tanto:
        $$
            |E:K| = |E:K(a)||K(a):K| = |E(b):K(b)||K(b):K| = |E(b):K| = |E(b):E||E:K| \text{ usando el teorema \ref{thm:transgr}.}
        $$
        Con lo que concluimos:
        $$
            |E(b):E| = 1 \implies E(b) = E \implies b\in E
        $$
    \end{itemize}
\end{proof}


\begin{ex}[H3.5]
    Decide si son normales las siguientes extensiones:
    \begin{enumerate}
        \item $\Q(\sqrt{5}i)/\Q$.\\\\

        Es normal ya que $\Q(\sqrt{5}i) = \Q(x^2+5)$ es el cuerpo de escisión de $x^2+5$.
        \item $\Q(\sqrt{5})/\Q$.\\\\

        Es normal ya que $\Q(\sqrt{5})/\Q = \Q(x^2-5)$ es el cuerpo de escisión de $x^2-5$.

        \item $\Q(\sqrt[4]{5})/\Q$.\\\\

        No es normal. Sabemos que $\Q(\sqrt[4]{5})/\Q$ es finita, y como $p(x) = x^4 - 5$ es irreducible en $\Q[x]$, tiene una raíz en $\Q(\sqrt[4]{5})$ pero no escinde, por la caracterización de normalidad (teorema \ref{thm:3.9}), concluimos con que $\Q(\sqrt[4]{5})/\Q$ no es normal.
    \end{enumerate}
\end{ex}

\begin{ex}[H3.6]
    Demuestra que $\Q(\sqrt[3]{2})/\Q$ no es normal. Encuentra una extensión normal que contenga $\Q(\sqrt[3]{2})$ como subcuerpo.\\\\

    Sabemos que no es normal porque sea $p(x) = x^3 - 2 \in \Q[x]$, se puede comprobar que $p(x)$ es irreducible y tiene una raíz en $\Q(\sqrt[3]{2})$ pero no se escinde. Por tanto, por el teorema \ref{thm:3.9}, el cuerpo no es normal.\\
    Una extensión normal que contiene a $\Q(\sqrt[3]{2})$ puede ser $\Q(\sqrt[3]{2}, \omega)$ con $\omega \in \C,\ o(\omega) = 3$.
\end{ex}

\begin{obs}
    Si $E/K$ es normal y $f\in K[x]$ no necesariamente irreducible que tiene una raíz en $E$, $f$ no tiene por que escindirse.
\end{obs}

%%%%%%%%%%%%%%%%%%%%%%%%%%%%%%%%%%%%%%%%%%%%%%%%%%%%%%%%%%%%%%%%%%%%%%%%%% 28/10
\section{El grupo de Galois de una extensi\'on}

\subsection{Acción de un grupo}

Vamos a ver un repaso de algunos conceptos de la acción de un grupo sobre un conjunto.

\begin{dfn}[Acción de un grupo]
    Una \textbf{acción} de un grupo $(G, \star)$ sobre un conjunto $X$ es una aplicación: $\phi: G \times X \to X$ que cumple:
    \begin{enumerate}
        \item $ \forall x \in X,\ \phi(e, x) = x$ donde $e$ es el elemento neutro del grupo.
        \item $ \forall x \in X,\ g, h \in G,\ \phi(g \star h, x) = \phi(g, \phi(h \star x))$
    \end{enumerate}
\end{dfn}
\begin{obs}$ $
    \begin{itemize}
        \item El cumplimiento de las dos condiciones hace que la aplicación $\phi(g, \cdot): X \to X$ sea biyectiva para cada $g \in G$. Con esto se puede dar una definición alternativa de la acción de un grupo:
        $$
            \phi: G \to \sdf{\text{funciones biyectivas } X \to X}
        $$
        \item Una notación alternativa usada para las acciones es:
        $$
            (g, x) \mapsto x \cdot g
        $$
        Con lo que las dos condiciones se reescriben como:
        \begin{enumerate}
            \item $x \cdot e = x$, con $e$ el neutro del grupo.
            \item $(x \cdot g) \cdot h = x \cdot (g \cdot h),\ \forall g,h \in G,\ x \in X$
        \end{enumerate}
    \end{itemize}
\end{obs}

\begin{pro}[Una acción es un homomorfismo]
    Sea $(G, \cdot)$ un grupo, $\Omega \neq \varnothing$ un conjunto. Si $G$ actúa sobre $\Omega$ (lo denotamos por $\Omega^{\actar G}$), la aplicación:
    $$
        \rho: G \to S_{\Omega};\ g \mapsto f \circ g: \Omega \to \Omega;\ \omega \mapsto \omega \cdot g
    $$
    es un homomorfismo de grupos, en particular $\ker(\rho) \normsub G$. Y tenemos que:
    $$
        \quo{G}{\ker \rho} \isom \rho(G) \leq S_{\Omega}
    $$
    Recíprocamente, todo homomorfismo $\rho: G \to S_\Omega $ define una acción de $G$ en $S_\Omega$.
\end{pro}

Además, sea $(G, \cdot)$ un grupo, $\Omega \neq \varnothing$ un conjunto. Si $G$ actúa sobre $\Omega$ ($\Omega^{\actar G}$), entonces define una relación binaria de equivalencia (r.b.e) de la forma:
$$
    \omega \sim \omega' \iff \exists g \in G:\ \omega \cdot g = \omega'
$$
Las clases de equivalencia definidas por esta relación se denominan $G$-órbitas.

\begin{dfn}[$G$-órbita de un elemento de un conjunto]
    Sea $G$ un grupo, la \textbf{órbita} de un elemento $\omega$ de un conjunto $\Omega$ es la clase de equivalencia:
    $$
        \omega^G = \theta_\omega = \sdf{\omega \cdot g \mid \forall g \in G} \subseteq \Omega
    $$
\end{dfn}

\begin{dfn}[Estabilizador de un elemento de un conjunto]
    Sea $G$ un grupo, el \textbf{estabilizador} de un elemento $\omega$ de un conjunto $\Omega$ es:
    $$
        G_\omega = \mathrm{St}(\omega) = \sdf{g \in G \mid \omega \cdot g = \omega} \leq G
    $$
    que es un subgrupo de $G$.
\end{dfn}

\begin{thm}[Teorema de la Órbita-Estabilizador]\label{thm:orbita}
    Sea $G$ un grupo que actúa sobre $\Omega$ y $\omega \in \Omega$ un elemento, entonces la longitud de la $G$-órbita de $\omega$ ($\left|\theta_\omega\right| = \vabs{G_\omega}$) es:
    $$
        \vabs{G_\omega} = \vabs{\theta_\omega} = \vabs{G:G_\omega}
    $$
\end{thm}

\begin{dfn}[Puntos fijos por la acción de un grupo]
    Sea $G$ un grupo y $\Omega$ un conjunto.
    \begin{itemize}
        \item Los puntos fijos por la acción de $G$ son los pertenecientes al conjunto:
        $$
            \Omega^G = \sdf{ \omega \in \Omega \mid \omega \cdot g = \omega,\ \forall g \in G}
        $$
        \item Los puntos fijos por la acción de $g \in G$ son los pertenecientes al conjunto:
        $$
            \Omega^g = \sdf{ \omega \in \Omega \mid \omega \cdot g = \omega}
        $$
    \end{itemize}
\end{dfn}

\subsection{Grupo de Galois}

En esta subsección, vamos a considerar $E$ como una extensión. Además, sea $\aut(E)$ el conjunto de automorfismos de $E$, podemos comprobar que $(\aut(E), \circ)$ es un grupo (no abeliano en general).\\

El grupo $(\mathrm{Aut}(E), \circ)$, actúa sobre $E$. ($e \cdot \sigma = \sigma(e) = e$).

\begin{dfn}[Grupo de Galois]
Sea $\sigma \in \aut(E)$, $E^\sigma = \sdf{e \in E \mid e \sigma = e}$ el conjunto de puntos fijos de $E$ por la acción de $\sigma$ (es subcuerpo de $E$). Sea $E/K$ una extensión, llamamos \textbf{grupo de Galois} de $E/K$ a:
$$
    \gal(E/K) = \sdf{\sigma \in \aut(E) \mid \sigma(k) = k \forall k \in K} = \sdf{\sigma \in \aut(E) \mid K \subseteq E^\sigma}
$$
Donde además, el \textit{grupo de Galois} es un cuerpo intermedio de $E/K$.
\end{dfn}

\begin{pro} % TODO: Dar nombre
Sea $L$ un cuerpo intermedio de $E/K$, entonces: $\gal(E/L) \leq \gal(E/K)$. En general $\left.\sigma\right|_L \notin \gal(L/K)$.
\end{pro}

\begin{obs}[Notación]
    Sea $\sigma \in \gal(E/K)$, decimos que es un $K$-automorfismo de $E$ si fija todo $K$ elemento a elemento.
\end{obs}

\begin{dfn}[$K$-isomorfismo]
    Sea $E/K$ una extensión, $L_1, L_2$ cuerpos intermedios, $\sigma: L_1 \to L_2$ es un $\mbf K$\textbf{-isomorfismo} si $\sigma: L_1 \to L_2$ es un isomorfismo y $\sigma(k) = k,\ \forall k \in K$.
\end{dfn}

Con esta notación vamos a reescribir el corolario \ref{cor:3.8}.
\begin{thm}\label{thm:3.8bis}
    Sea $E/K$ extensión normal, $L_1, L_2$ cuerpos intermedios. Si $\sigma: L_1 \to L_2$ es un $K$-isomorfismo, entonces: $\exists \tau \in \gal(E/K)$ tal que $\left. \tau \right|_{L_1} = \sigma$.
\end{thm}

\begin{cor}[al teorema \ref{thm:3.8bis}] \label{cor:3.10}
    Sea $E/K$ una extensión normal, $p \in K[x]$ un polinomio irreducible. Si $a$ y $b$ son raíces de $p$ en $E$, entonces $\exists \tau \in \gal(E/K):\ \tau(a) = b$.
\end{cor}
\begin{proof}
    Por el corolario \ref{cor:2.6}, existe un $K$-isomorfismo $\sigma: K(a) \to K(b)$. Por el teorema \ref{thm:3.8bis}, $\sigma$ se extiende a $\tau \in \gal(E/K)$.
\end{proof}

\begin{obs}
    $\gal(E/K)$ actúa sobre las raíces de $f \in K[x]$ en $E$.
\end{obs}

\begin{thm}[Grupo de Galois y raíces de un polinomio]\label{thm:3.11}
    Sea $E/K$ una extensión, $0 \neq f \in K[x]$, $\Omega = \sdf{a_1, \ldots, a_n}$ el conjunto de todas las raíces distintas de $f$ en $E$ (no quiere decir que $E$ contenga todas las raíces de $f$), con $n \leq 1$ y sea $L = K(\mbf a) \subseteq E$ un subcuerpo intermedio, entonces:
    \begin{itemize}
        \item[(a)] Si $a \in \Omega$ y $\sigma = \gal(E/K)$, entonces $\sigma(a) \in \Omega$. En particular, $\left.\sigma\right|_\Omega \in S_\Omega$ y $\sigma(L) = L$.
        \item[(b)] La aplicación $\psi: \gal(E/K) \to \gal(L/K);\ \sigma \mapsto \left.\sigma\right|_L$ es un homomorfismo de grupos con $\ker(\psi) = \gal(E/L) \normsub \gal(E/K)$.
        \item[(c)] Si $E/K$ es normal, entonces $\psi$ es sobreyectivo. En particular, $\quo{\gal(E/K)}{\gal(E/L)} \isom \gal(L/K)$.
        \item[(d)] La aplicación $\rho: \gal(E/K) \to S_\Omega;\ \sigma \mapsto \left.\sigma\right|_\Omega$ es un homomorfismo de grupos con $\ker(\rho) = \gal(E/L)$.
        \item[(e)] Si $E/K$ es normal y $f$ es irreducible, dados $a, b \in \Omega$ existe $\sigma \in \gal(E/K)$ tal que $\sigma(a) = b$. Por tanto:
        $$
            n = \left|\gal(E/K):\gal(E/K(a))\right|
        $$
    \end{itemize}
\end{thm}

\begin{obs}[al teorema \ref{thm:3.11}]
    $$
        \left(\gal(L/L) = \sdf{id}\right)
    $$
    En el teorema, (d) aplicado a $L/K$ implica que:
    $$
        \gal(L/K) \isom \rho(\gal(L/K)) \leq S_\Omega
    $$
\end{obs}

\begin{proof}$ $
    \begin{itemize}
        \item[(a)]
        $$
            0 = \sigma(0) = \sigma(f(a)) = \sigma(f)(\sigma(a)) = f(\sigma(a)) \implies \sigma(a) \in \Omega
        $$
        Como:
        $$
            \left. \sigma \right|_\Omega: \Omega \to \Omega \text{ es inyectiva,}
        $$
        entonces:
        $$
            \Omega \text{ es finita } \implies \left. \sigma \right|_\Omega \in S_\Omega \text{ es biyectiva,}
        $$
        por tanto:
        $$
            \sigma(L) = \sigma(K(a_1, \ldots, a_n)) = \sigma(K)(\sigma(a_1), \ldots, \sigma(a_n)) = K(a_1, \ldots, a_n) = L
        $$
        \item[(b)] Se deja como ejercicio.
        \item[(c)] Sea $\theta \in \gal(L/K)$, por el teorema \ref{thm:3.8bis}, $\exists \sigma \in \gal(E/K)$ tal que $\left. \sigma \right|_L = \theta$. En particular, $\psi$ es sobreyectiva. (La segunda conclusión se sigue del teorema de isomorfía).
        \item[(d)] Se deja como ejercicio partiendo de $\left. \sigma \right|_\Omega \in S_\Omega$.
        \item[(e)] Por el corolario \ref{cor:3.10} si $a, b \in \Omega$ cualesquiera, $\exists \sigma \in \gal(E/K)$ con $a\sigma = \sigma(a) = b$ y $\theta_a = \sigma$ por el teorema \ref{thm:orbita}.
        %%TODO: COMPLETAR demostración del cardinal. Santorum.
    \end{itemize}
\end{proof}

%%%%%%%%%%%%%%%%%%%%%%%%%%%%%%%%%%%%%%%%%%%%%%%%%%%%%%%%%%%%%%%%%%%%%%%%%% 29/10

\begin{cor}[Caracterización de normalidad a través de isomorfismos]\label{cor:3.12}
    Sea $E/K$ una extensión, $L$ un cuerpo intermedio:
    \begin{itemize}
        \item[(a)] Si $L/K$ es normal entonces $\tau \in \gal(E/K)$:
        $$
            \tau(L) = L \left(\rest{\tau}_L\right) \in \gal(L/K)
        $$
        \item[(b)] Si $E/K$ es normal:
        $$
            L/K \text{es normal} \iff \forall \tau \in \gal(E/K),\ \tau(L) = L
        $$
        En particular $\gal(E/L) \normsub \gal(E/K)$ y $\quo{\gal(E/K)}{\gal(E/L)} \isom \gal(L/K)$.
    \end{itemize}
\end{cor}

\begin{proof}$ $
    \begin{itemize}
        \item[(a)] $L = K(f) = K(a_1, \ldots, a_n)$. $f \in K[x] = K(a_1, \ldots, a_n)$. Con $\Omega = \sdf{a_i}$ todas las raíces distintas entre sí de $f$ en $L$.
        Por el teorema \ref{thm:3.11}(a), si $\sigma \in \gal(E/K)$ $\sigma(L)=L$.
        \item[(b)] Supongamos $E/K$ normal.
        \begin{itemize}
            \item[$(\implies)$] Directo de firma análoga al apartado (a).
            \item[$(\ \Longleftarrow \ )$] Suponemos que $\tau \in \gal(E/K),\ \tau(L) = L$ y queremos probar que $L/K$ es normal. Para ello, usamos el teorema \ref{thm:3.9}.\\

            Sea $p \in K[x]$ irreducible con una raíz $a \in L$m queremos ver que $p$ se escinde en $L$. Por el teorema \ref{thm:3.9}, $p$ se escinde en $E$ (pues $E/K$) es normal. Sea $b \in E$ una raíz de $p$, por el corolario \ref{cor:3.10}, $\exists \tau \in \gal(E/K)$ tal que $\tau(a) = b$ y además $a \in L \implies b = \tau(a) \in L$.\\

            Como $L$ tiene todas las raíces de $p$, entonces $p$ se escinde en $L$. La última parte se demuestra de forma directa aplicando el teorema \ref{thm:3.11}(c).
        \end{itemize}
    \end{itemize}
\end{proof}

\begin{obs}
    Puede parecer que el corolario \ref{cor:3.12}(b) no nos da una caracterización por que requiere estar dentro de una extensión normal, pero si $L/K$ es finita, existe una extensión $E/K$, $L \subseteq E$, tal que $E/K$ es normal (y es la menor con estas propiedades). $E$ es la \textit{clausura normal} de $L/K$.\\\\

    Por ejemplo, sea $L = K(a_1, \ldots, a_n)$ (por ejemplo adjuntando una $K$-base). Sea $a_i \in L$ es algebraico sobre $K$. Para constuir $E$ la clausura normal:
    $$
        p_i = Irr(K, a_i),\ f = \prod_{i=1}^{n} p_i, \text{ entonces } E = K(f)
    $$
\end{obs}

\begin{pro}[Finitud de las extensiones de Galois]\label{pro:ext-gal-finit}
    Sea $L/K$ una extensión, si $L/K$ es finita entonces $\gal(L/K)$ es finita.
\end{pro}

\begin{proof}
    Si $L/K$ es normal, entonces sea $L = K(a_1, \ldots, a_n)$ y $\Omega = \sdf{a_i}$ todas las raíces distintas entre sí de $f \in K[x]$. Por \ref{thm:3.11}(d), $\gal(L/K) \leq S_\Omega \isom S_n$, donde $S_n$ es el enésimo grupo simétrico.\\
    Sabemos que $\vabs{S_n} = n! \implies \text(por Lagrange) \vabs{\gal(L/K)} \divides \vabs{S_n} = n!$.\\
\end{proof}

\begin{eg}[Ejemplos de ejercicios de extensiones de Galois]$ $
    \begin{enumerate}
        \item Sea $E = \quo{\F_2[x]}{\gen{x^2 + x + 1}} = \sdf{\bar 0, \bar 1, \bar x, \bar x + \bar 1}$, además $E/\bar \F_2$, de hecho, $E = \bar \F_2(\bar x)$. Comprueba que $E = \bar\F_2(y^2 + y + 1)$ y calcula $\gal(E/\bar \F_2)$.\\\\

        Por el lema $\ref{lm:3.3}$ $\bar x$ es una raíz, y $y^2 + y + 1 \in \bar\F_2[y]$. Sabemos además que $y^2 + y + 1$ tiene como máximo dos raíces distintas en su cuerpo de escisión.\\
        $\gal(E/\bar\F_2)$ manda $\bar x$ en otra raíz de $y^2 + y + 1$. Es decir, sea $\sigma \in \gal(E/K), f\in K[x], a\in E$ una raíz de $f$, entonces:
        $$
            0 = \sigma(0) = \sigma(f(a)) = \sigma(f)(\sigma(a)) = f(\sigma(a))
        $$
        Sabemos además que $Frob \in \gal(E/\bar \F_2)$ donde $Frob$ es el automorfismo de Frobenius, ya que fija el cuerpo primo de $E$ que es $\bar \F_2$. Por tanto:
        $$
            Frob(\bar x) = \bar x^2 = \bar x + \bar 1 \in E \text{ es otra raíz de } y^2 + y + 1
        $$
        Por tanto, $E = \bar \F_2(\bar x, \bar x + \bar 1)$, es el cuerpo de escisión de $y^2 + y + 1$.\\

        Para calcular $\gal(E/\bar \F_2)$, por \ref{thm:3.11}(d): $\sdf{id, Frob} = \gal(E/\bar \F_2) \leq S_2$ Donde $\vabs{S_2} = 2 \implies \gal(E/\bar \F_2) = \sdf{id, Frob} = \gen{Frob}$.

        \item Sea $f(x) = x^3 - 2 \in \Q[x]$.
        \begin{itemize}
            \item[(a)] $E = \Q(f)$.
            \item[(b)] $\vabs{E: \Q}$
            \item[(c)] Calcula $\gal(E/\Q)$.
        \end{itemize}

        Vamos a ver la solución:
        \begin{itemize}
            \item[(a)] Las raíces de $f$ en $\C$ son $\sdf{\sqrt[3]{2}, \sqrt[3]{2} \omega, \sqrt[3]{2} \omega^2}$ donde:
            $$
                \omega = exp\left( \frac{2\pi}{3} \right) = -\frac{1}{2} + \frac{\sqrt 3}{2} i
            $$
            y por tanto:
            $$
                E = \Q(\sqrt[3]{2}, \sqrt[3]{2} \omega, \sqrt[3]{2} \omega^2) = \Q(\sqrt[3]{2}, \omega) = \Q(\sqrt[3]{2}, \sqrt{3} i)
            $$
            \item[(b)] Sabemos por el teorema \ref{thm:2.3} que $\vabs{\Q(\sqrt[3]{2}):\Q} = 3$ y $\vabs{\Q(\omega):\Q} = 2$. Por tanto:
            $$
                \Q \subseteq \Q(\sqrt[3]{2}) \subseteq \Q(\sqrt[3]{2})(\omega) \implies \vabs{\Q(\sqrt[3]{2}, \omega):\Q} = 6
            $$
            \item[(c)] Sea $L = \Q(\omega)$, sabemos que $\Q \subseteq L \subseteq E$. Por el teorema \ref{thm:3.11}(d):
            $$
                \gal(E/\Q) \leq S_3 \implies \vabs{\gal(E/\Q)} \divides 6
            $$
            Por el teorema \ref{thm:2.1}, sabemos que una $\Q$-base de $E$ es $\sdf{1, \sqrt[3]{2}, \sqrt[3]{4} \sqrt[3]{2} \omega, \sqrt[3]{4} \omega,  \omega}$. Luego $\tau \in \gal(E/\Q)$ queda determinado por $\tau(\sqrt[3]{2})$ y $\tau(\omega)$.\\
            Las distintas posibilidades de $\tau$ son:
            $$
                \tau(\omega) \in \sdf{\omega, \bar \omega} \text{ y } \tau(\sqrt[3]{2}) \in \sdf{\sqrt[3]{2}, \sqrt[3]{2} \omega, \sqrt[3]{2} \bar \omega}
            $$
            Donde con un rápido cálculo combinatorio vemos que hay como mucho $6$ posibles $\tau$.
            %% TODO: Completar de Santorum.
        \end{itemize}
    \end{enumerate}
\end{eg}

%%%%%%%%%%%%%%%%%%%%%%%%%%%%%%%%%%%%%%%%%%%%%%%%%%%%%%%%%%%%%%%%%%%%%%%%%% 31/10
\begin{eg}[Hallar grupo de Galois dado un polinomio y una extensión]
    Sea $f \in \Q[x]$, $f(x) = x^4 - 2$. Y sea $E = \Q(\sqrt[4] 2)$, entonces el conjunto $\Omega$ del teorema \ref{thm:3.11} es $\Omega = \sdf{\sqrt[4] 2, -\sqrt[4] 2}$. Sea $L = \Q(\Omega) = E$, por el teorema \ref{thm:3.11}(d), $\gal(E/\Q) \leq S_\Omega \isom S_2 \implies |\gal(E/\Q)| = 1$ o $2$ por el teorema de Lagrange de grupos. Además:
    $$
        \gal(E/\Q) = \aut(E) = \aut(\Q(\sqrt[4]{2}))
    $$
    Por \ref{cor:2.6}, $\exists K$-isomorfismo $\sigma$ tal que:
    $$
        \sigma: \Q(\sqrt[4]{2}) \to \Q(-\sqrt[4] 2) = E;\ \ \sqrt[4]{2} \mapsto -\sqrt[4]{4} \implies \sigma \in \gal(E/\Q) = \sdf{1, \sigma}
    $$
    Considerando las extensiones:
    \begin{center}
        \begin{tikzpicture}[node distance=1cm and 2cm, auto]
            \node (A) {$\Q(\sqrt[4]{2}) = E$};
            \node (B) [right= of A] {$E$};
            \node (C) [below of =A] {$\Q(\sqrt{2}) = L$};
            \node (D) [below of =B] {$L$};
            \node (E) [below of =C] {$\Q$};
            \node (F) [below of =D] {$\Q$};
            \draw[->] (A) to node {}(B);
            \draw[->] (C) to node {}(D);
            \draw[->] (E) to node {}(F);
            \draw[-] (A) to node {}(C);
            \draw[-] (B) to node {}(D);
            \draw[-] (C) to node {}(E);
            \draw[-] (D) to node {}(F);
        \end{tikzpicture}\\
    \end{center}
    Donde la extensión $L/\Q$ es normal pero $E/\Q$ no.  Además, sea $\theta: L \to L;\ \sqrt 2 \mapsto \sqrt 2$ no se extiende a $E$.\\

    Este ejemplo muestra que la hipótesis $E/K$ normal es necesaria en el teorema \ref{thm:3.8bis}.
\end{eg}

\section{Extensiones separables}

El siguiente teorema y proposición se vieron en el capítulo 1 pero vamos a reescribirlos con el lenguaje actual, es por ello que no se aporta demostración de los mismos.

\begin{thm}[Derivada formal y raíces multiples]\label{thm:3.13}\label{thm:3.4.1}$ $
    \begin{itemize}
        \item[(a)] Si $f \in K[x]$ tal que $mcd(f, f') = 1$ y $f' \neq 0$ donde $f'$ es la derivada formal, entonces todas las raíces de $f$ en su cuerpo de escisión son distintas.
        \item[(b)] Si $f \in K[x]$ irreducible y $f'$ su derivada formal tal que $f' \neq 0$, entonces $f$ es separable.
    \end{itemize}
\end{thm}
\begin{obs}
    Si $\car K = 0$ y $f\in K[x]$ con $\delta f > 1$, entonces $f' \neq 0$.
\end{obs}

\begin{eg}[Aplicación del contrarrecíproco del teorema \ref{thm:3.13}]
    Sea $f(x) = x^p -1$, $\car K = p$, vemos que $f' = 0$. Entonces no todas las raíces de $f$ en su cuerpo de escisión son distintas, de hecho $f = (x-1)^p$.
\end{eg}

\begin{pro}[Multiplicidad de raíces y derivada formal]\label{pro:3.4.2}
    Sea $E/K$ una extensión, $f \in K[x]$  un polinomio y $a \in E$ tal que $f(a) = 0$, y sea $m$ la multiplicidad de la raíz $a$, entonces:
    $$
        m > 1 \iff f(a)=0=f'(a)
    $$
\end{pro}

\begin{dfn}[Separabilidad. Polinomios, elementos y extensiones]$ $
    \begin{itemize}
        \item[($i$)] Sea $f \in K[x]$ un polinomio \textbf{irreducible}, decimos que es un \textbf{polinomio separable} si todas sus raíces en un cuerpo de escisión son distintas, es decir, $f' \neq 0$ (por el teorema \ref{thm:3.13}).
        \item[($ii$)] Sea $f \in K[x]$ un polinomio reducible, decimos que es un \textbf{polinomio separable} si sus factores irreducibles son separables.
        \item[($iii$)] Sea $E/K$ una extensión, $a \in E$ un elemento algebraico, decimos que $a$ es un \textbf{elemento separable} si $p = Irr(K, a)$ es separable.
        \item[($iv$)] Sea $E/K$ una extensión algebraica, decimos que es una \textbf{extensión separable} si todo $a \in E$ es separable.
    \end{itemize}
\end{dfn}

\begin{pro}[Divisores irreducibles de un polinomio separable]\label{pro:sepdiv}
    Sea $f \in K[x]$ un polinomio separable, entonces para todo $p$ irreducible tal que $p \divides f$, $p$ es separable. El recíproco es cierto por definición.
\end{pro}
\begin{proof}
    Si $p$ divide a $f$ quiere decir que todas sus raíces están contenidas en las de $f$, por tanto, si $f$ es separable, $p$ también lo es.
\end{proof}

\begin{pro}[Separabilidad de extensiones intermedias]
    Sea $E/K$ una extensión separable, y $L$ un cuerpo intermedio, entonces, $E/L$ y $L/K$ son separables.
\end{pro}
\begin{proof}
    Es fácil ver que $L/K$ es separable. Vamos a ver que $E/L$ es separable.\\
    Sean $a \in E$, $p = Irr(L, a)$. Y sea $q = Irr(K, a)$ que es separable por hipótesis. Como $p \divides q$ (en $L[x]$), sea $M$ un cuerpo de escisión de $q$ sobre $L$ (que sabemos que existe por el teorema \ref{thm:3.4}), entonces $q$ se escinde en $M \implies p$ se escinde en $M$ por \ref{lm:3.1}($ii$).\\

    Si $\alpha, \beta$ son raíces de $p$ en $M$, entonces $\alpha, \beta$ son raíces de $q$ en $M \implies \alpha \neq \beta$.
\end{proof}

\begin{lm}[Separabilidad en cuerpos de característica $0$]\label{lm:3.4.3}
    Todo polinomio $f \in K[x]$ irreducible (no constante) con $\car K = 0$ es separable y toda extensión $E/K$ algebraica es separable.
\end{lm}
\begin{proof}
    Como $f$ es irreducible, por el teorema \ref{thm:3.4.1}(b), si $f' \neq 0$ entonces es separable. Como estamos en característica $0$, $f'(x) = 0 \iff f(x) = cte$, con lo que hemos demostrado que $f(x)$ es separable si no es constante.
\end{proof}

\begin{cor}[Separabilidad en cuerpos de característica $p$]\label{cor:3.4.4}
    Sea $f \in K[x]$ irreducible con $\car K = p$:
    $$
        f(x) \text{ es separable } \iff f \notin K[x^p]
    $$
\end{cor}

Sin embargo el corolario no nos resuelve una pregunta. \textit{¿Podemos encontrar polinomios irreducibles que no sean separables?}, es decir, \textit{ ¿podemos encontrar un polinomio irreducible en $K[x^p]$?}.\\

Recordemos que un cuerpo $K$ es perfecto si el monomorfismo de Frobenius es un isomorfismo, es decir, $\varphi = Frob \in \aut(K)$, donde $\varphi(a) = a^p$.\\
En particular $\forall b \in K, \exists a \in K$ tal que $a^p = b$

\begin{thm}[Irreducibilidad sobre cuerpos perfectos]\label{thm:3.4.5}
    Sea $K$ un cuerpo perfecto, todo polinomio irreducible sobre $K$ es separable. En particular, todo polinomio es separable.
\end{thm}
\begin{proof}
    Sea $K$ perfecto, entonces:
    $$
        \forall b \in K,\ \exists a \in K:\ a^p = b.
    $$
    Sea $f\in K[x]$ irreducible. Sabemos por el corolario \ref{cor:3.4.4} que si $f \notin K[x^p]$, entonces $f$ es separable. Supongamos que $f \in K[x^p]$ de la forma:
    \begin{align*}
        f &= b_0 + b_1 x^{pn_1} + \cdots + b_k x^{pn_k}\\
          &= a_0^p + \cdots + a_k^p (x^{n_k})^p\\
          &= a_0^p + \cdots + (a_k x^{n_k})^p\\
          &= (a_0 + \cdots a_k x^{n_k})^p \text{ que no es irreducible }
    \end{align*}
\end{proof}

%%%%%%%%%%%%%%%%%%%%%%%%%%%%%%%%%%%%%%%%%%%%%%%%%%%%%%%%%%%%%%%%%%%%%%%%%% 04/11
\begin{eg}[Polinomio irreducible que no es separable]
    Vamos a buscar un polinomio $f$ irreducible que no sea separable. Además por el lema \ref{lm:3.4.3} y el teorema \ref{thm:3.4.5}, tenemos que buscar $f \in K[x]$ tal que $\car K = p$ con $K$ no perfecto.\\
    Por ejemplo, si consideramos $K = \F_p$, podemos ver que $K$ no es perfecto ya que $\varphi = Frob \in \mathrm{End}(\F_p(t))$.\\
    Por otra parte, $\varphi(\F_p(t)) = \F_p(t^p) \subsetneq \F_p(t)$. Vamos a considerar entonces la extensión $\F_p(t)/F_p(t^2)$, donde podemos ver que $\vabs{\F_p(t) : \F_p(t^2)} = p$ pues tenemos que $t$ es raíz de $x^p - t^p \in \F_p(t^p)[x]$.\\

    Tenemos que ver ahora que $x^p - t^p = (x - t)^p$ es irreducible. En otro caso, $x^p - t^p = g(x) \cdot h(x)$, con $\delta g,\ \delta h \leq p$. Por tanto:
    $$
        g(x) = (x - t)^n \text{ con } 0 < n < p
    $$
    Lo que es imposible por que el término independiente de $g$ es $t^n$ que no pertenece a $\F_p(t^p)$. Luego $Irr(\F_p(t^p), t) = x^p - t^p$.
\end{eg}

\begin{pro}[Raíces comunes y polinomios irreducibles mónicos]
    Sean $p \neq q \in K[x]$ irreducibles mónicos, entonces $p$ y $q$ no pueden tener una raíz en común.
\end{pro}

\begin{proof}
    Sea $E/K$ una extensión y $\alpha \in E$ la raíz común de $p$ y $q$, entonces $p = Irr(K, \alpha) = q$ por la unicidad del polinomio mónico.\\\\
    Otra forma de verlo es recordar que sean $f, g \in K[x]$ entonces $mcd_K(f, g) = mcd_E(f, g)$. Si $\alpha$ es una raíz común entonces $p \divides q$ y además $(x - \alpha) \divides mcd(p, q) > 1 \implies q \divides p$ y por tanto $p = q$.
\end{proof}

\begin{ex}[H3.??]
    Calcular las raíces distintas de $f \in K[x]$ separables en su cuerpo de escisión.\\\\
    Vamos a dar un esquema de la resolución del ejercicio:
    \begin{enumerate}
        \item Descomponer $f$ es factores irreducibles:
        $$
            f = a p_1^{e_1} p_2^{e_2} \ldots p_r^{e_r} \text{ donde los $p_i$ son irreducibles y distintos entre sí.}
        $$
        \item Por definición (por ser separable) cada $p_i$ es separable y además el número de raíces de $p_i$ coincide con $\delta p_i$.\\
        Con esto vemos que el número de raíces distintas de $f$ es igual a $\sum \delta p_i$.
    \end{enumerate}
\end{ex}


\begin{dfn}[Elemento primitivo]
    Sea $E/K$ una extensión simple, a un elemento de $\gamma \in E$ se le denomina \textbf{elemento primitivo} si $E = K(\gamma)$.
\end{dfn}

\begin{thm}[del Elemento Primitivo]\label{thm:elemprim}
    Toda extensión finita y separable es simple.
\end{thm}

\begin{proof}
    Se proporciona un esquema de la demostración. Sea $E/K$ finita y separable, entonces $E = K(\alpha_1, \ldots, \alpha_n)$ con $\alpha_i$ algebraicos y separables.
    \begin{itemize}
        \item Por inducción sobre $n$, basta probar el resultado cuando $n =2$, con $n = 1$ habríamos terminado pues $E = K(\alpha_1)$ y por tanto sería simple. Vamos a considerar el caso $n > 1$, es decir $L = K(\alpha_1, \ldots, \alpha_{n-1})$. Por inducción $L = K(\alpha)$ y entonces $E = K(\alpha_1, \ldots, \alpha_n) = K(\alpha_1, \alpha_{n-1})(\alpha) = K(\alpha)(\alpha_n) = K(\alpha, \beta)$
        \item Queremos probar que $K(\alpha, \beta)/K$ es simple. Esta parte es constructiva. Habría que conectar separabilidad y normalidad.
    \end{itemize}
\end{proof}


\begin{pro}[Irreducibilidad y acción del grupo de Galois] % Lo vimos como ejercicio
    Sea $f \in K[x]$ un polinomio, $E = K(f)$ y $G = \gal(E/K)$. Entonces:
    \begin{enumerate}
        \item Si $f$ es irreducible, $G$ actúa transitivamente (la acción define una sola G-órbita) sobre el conjunto $\Omega$ de raíces de $f$.
        \item Si $f$ no tiene raíces múltiples en $E$ y $G$ actúa transitivamente sobre $\Omega$, entonces $f$ es irreducible.
    \end{enumerate}
\end{pro}
\begin{proof}$ $
    \begin{enumerate}
        \item Se sigue inmediatamente del corolario \ref{cor:3.10}.
        \item Sea $f = a p_i \ldots p_r$ con $p_i$ primos irreducibles. Supongamos que $r > 1$, sea $\alpha \in E$ raíz de $p_1$, $\beta \in E$ raíz de $p_2$. Por hipótesis $\exists \sigma \in G$ tal que $\alpha \cdot \sigma = \sigma(\alpha) = \beta$. Entonces:
        $$
            0 = \sigma(p_1(\alpha)) = \sigma(p_1)(\sigma(\alpha)) = p_1 \beta \implies Irr(K, \alpha) = p_1 = Irr(K, \beta) = p_2
        $$
    \end{enumerate}
\end{proof}

\begin{thm}[Numero de extensiones de un isomorfismo]\label{thm:3.4.6}
    Sea $\sigma: K_1 \to K_2$ un isomorfismo, $f_1 \in K_1[x]$, $f_2 = \sigma(f_1) \in K_2[x]$ polinomios y $E_i = K(f_i)$ extensiones. Entonces:
    \begin{itemize}
        \item[(a)] $\sigma$ se extiende a un isomorfismo $\tau: E_1 \to E_2$. (Teorema \ref{thm:3.4.5}).
        \item[(b)] Si $f_1$ es separable entonces hay exactamente $\vabs{E_1:K_1}$ extensiones como en (a).
    \end{itemize}
\end{thm}

%%%%%%%%%%%%%%%%%%%%%%%%%%%%%%%%%%%%%%%%%%%%%%%%%%%%%%%%%%%%%%%%%%%%%%%%%% 05/11

\begin{proof}$ $
    \begin{itemize}
        \item[(a)] Teorema \ref{thm:3.4.5}.
        \item[(b)] Por inducción sobre $\vabs{E_1:K_1}$.
        \begin{itemize}
            \item $\vabs{E_1:K_1} = 1 \iff g_1$ se escinde en $E_1 \iff f_2$ se escinde en $E_2 \iff \vabs{E_2:K_2} = 1 \implies \tau = \sigma$
            \item Podemos suponer que $\vabs{E_1:K_1} > 1$, entonces:
            $$
                f_1 = p_1g_1 \in K_1[x], \text{ donde } p_1 \text{ es irreducible y } \delta p_1 = d \geq 2.
            $$
            por tanto:
            $$
                f_2 = p_2g_2, \text{ donde } p_2 = \sigma^\star(p_1) \text{ es irreducible y } \delta p_2 = \delta p_1 = d \geq 2.
            $$
            Como $f_1$ es separable, entonces $p_1$ es separable (por la proposición \ref{pro:sepdiv}). Como $E_1 = K_1(f_1)$, entonces $p_1$ se escinde en $E_1$ (por el lema \ref{lm:3.1}(ii)). Ahora, desarrollando:
            \begin{align*}
                p_1(x) = (x-\alpha_1) \cdots (x - \alpha_d) \in E_1[x]\\
                \text{Por (a)} \exists \tau:E_1 \to E_2 :\ \rest{\tau}_{K_1} = \sigma\\
                p_2 = \tau^\star(p_1) = (x - \tau(\alpha_1)) \cdots (x - \tau(\alpha_d))\\
                \beta_j = \tau(\alpha_j) \text{ donde las $\beta_j$ son las raíces distintas de } p_2 \text{ en } E_2.
            \end{align*}
            y por tanto $p_2$ es spearable.\\\\
            Sea $\alpha = \alpha_1$, por el teorema \ref{thm:2.5}, para cada $j = 1, \ldots, d$, existe un isomorfismo:
            $$
                \hat\sigma : K_1(\alpha) \to K_2(\beta_j)
            $$
            que extiende a $\sigma$. Como:
            $$
                \vabs{E_1:K_1(\alpha_1)} = \frac{\vabs{E_1:K_1}}{\vabs{K_1(\alpha):K_1}} = \frac{\vabs{E_1:K_1}}{d} < \vabs{E_1:K_1}.
            $$
            Por inducción, para cada $j = 1, \ldots, d$, existen $\vabs{E_1:K_1}$ extensiones exactamente de $\hat\sigma_j$.\\
            Es fácil ver que tenemos $\frac{\vabs{E_1:K_1}}{d} \cdot d$ extensiones de $\sigma$ a $E_1 \to E_2$. ¿Por qué son todas?.\\

            Si $\tau: E_1 \to E:2$ extiende a $\sigma$, en particular $\tau$ manda $\sigma$ en una raíz de $\tau^\star(p_1) = \sigma^\star(p_1) = p_2$. Es decir, $\tau(\alpha) = \beta_j$ luego, $\rest{\tau}_{K_1(\alpha)} = \hat\sigma_j$.
            Por tanto, $\sigma$ tiene exactamente $\vabs{E_1:K_1}$ extensiones a $E_1 \to E_2$.
        \end{itemize}
    \end{itemize}
\end{proof}

\begin{dfn}[Extensión de Galois]
    Sean $E, K$ cuerpos, diremos que la extensión $E/K$ es una \textbf{extensión de Galois} si $E/K$ es normal y separable.
\end{dfn}
\begin{obs}
    Si la extensión $E/K$ es de Galois, y sea $L$ un cuerpo intermedio, entonces $E/L$ es una extensión de Galois.
\end{obs}

\begin{cor}\label{cor:3.4.7}
    Sea $E/K$ una extensión de Galois, y $K \subseteq L \subseteq E$, $L/K$ es una extensión de Galois. Además, cada $\sigma \in \gal(L/K)$ tiene exactamente $\vabs{E:L}$ extensiones a $\gal(E/K)$. En particular:
    $$
        \vabs{\gal(E/K)} = \vabs{E:K}
    $$
\end{cor}
\begin{proof}
    Consecuencia del teorema \ref{thm:3.4.6}. Los detalles se dejan como ejercicio.
\end{proof}

\begin{ex}[H3.7]
    Sea $\xi$ la raíz primitiva quinta de $\1$. Demuestra o calcula:
    \begin{itemize}
        \item[(a)] $\Q(\xi)/\Q$ es normal.
        \item[(b)] $\vabs{\Q(\xi):\Q}$.
        \item[(c)] $\gal{\Q(\xi)/\Q}$.
    \end{itemize}

    \textbf{Solución:}\\\\
    \begin{itemize}
        \item[(a)] $\xi$ es raíz de $x^5 - 1$ y las raíces de $x^5 - 1$ son $\sdf{1, \xi, \xi^2, \xi^3, \xi^4}$, que son las raíces del quinto polinomio ciclotómico. Entonces $\Q(\xi) = \Q(x^5 - 1)$. Por tanto, el polinomio irreducible $Irr(\Q, \xi) = \phi_5(x) = x^4 + x^3 + x^2 + x + 1$.
        \item[(b)] $\vabs{\Q(\xi):\Q} = \delta \phi_5(x) = 4$
        \item[(c)] Sabemos que $\Q(\xi)/\Q$ es una extensión de Galois, (es normal por (a)  y separable por que $\car \Q = 0$). Además:
        $$
            \vabs{\gal(\Q(\xi)/\Q)} = 4, \text{ (usando el teorema \ref{thm:3.11}(d) sabemos que } \gal(\Q(\xi)/\Q) \leq S_4 \text{)}
        $$
        Por el corolario \ref{cor:3.10}, para $j = 1, 2, 3, 4$ sabemos que $\exists \sigma_j \in \gal(\Q(\xi)/\Q)$ tal que $\sigma_j(\xi) = \xi^j$, donde $\sigma_1 = id$. Y además, un $\sigma \in \gal(\Q(\xi)/\Q) = \sdf{id, \sigma_2, \sigma_3, \sigma_4}$ está determinado por la imagen de $\xi$. De hecho, $\gal(\Q(\xi)/\Q) = \gen{\sigma_2} \isom C_4$.
    \end{itemize}
\end{ex}

\section{Cuerpos finitos}

Vamos a empezar con un recordatorio de \textit{Teoría de Grupos}.

\begin{dfn}[Orden del elemento de un grupo]
    Sea $G$ un grupo finito, $g \in G$ un elemento, decimos que el \textbf{orden} de $g$ es:
    $$
        o(g) = n = \min\sdf{m \mid g^m = 1 } < \infty
    $$
\end{dfn}

\begin{thm}[Resultados de orden en un grupo]\label{thm:3.5.1}
    Sea $G$ un grupo finito, $g \in G$ un elemento del grupo:
    \begin{itemize}
        \item[(a)] $g^m = 1 \iff o(g) \divides m$.
        \item[(b)] $o(g) \divides \vabs{G}$. En particular $g^{\vabs{G}} = 1$.
        \item[(c)] Si $o(g) = n$, entonces $\gen{g} = \sdf{1, g, \ldots, g^{n-1}}$ y $\gen{g} = \gen{g^k} \iff mcd(g, k)=1$.
        \item[(d)] $G$ es cíclico $\iff \exists g \in G:\ o(g) = \vabs{G}$.
    \end{itemize}
\end{thm}

\begin{obs}
    Sea $K$ un  cuerpo finito y $F$ su cuerpo primo. Entonces:
    $$
        F \isom \F_p \text{ y además } \vabs{K:F} = n
    $$
    $$
        K \isom F^n \isom \F_p^n \implies \vabs{K} = p^n
    $$
\end{obs}

\begin{thm}[Clasificación de cuerpos finitos]\label{thm:3.5.2}$ $
    \item[(a)] Para cada primo $p$ y natural $n$, existe un cuerpo $\vabs{E} = p^n$.
    \item[(b)] Dos cuerpos $E$ y $L$ con el mismo número de elementos son isomorfos.
\end{thm}

\begin{obs}
    $$
        \F_{p^n} = GF(p^n)
    $$
\end{obs}

\begin{proof}$ $
    \begin{itemize}
        \item[(a)] Sea $F = \F_p$ y $f(x) = x^{p^n} - x \in F[x]$, entonces $E = F(f)$ (por el teorema \ref{thm:3.4}).\\
        Sea $U$ el conjunto de raíces de $f$ en $E$, y como $f'(x) = -1$, entonces $mcd(f, f') = 1$. Por el teorema \ref{thm:3.4.1}(a), obtenemos que $\vabs{U} = p^n$.\\

        \begin{enumerate}
            \item $F \subseteq U$ ya que $\forall a \in F$ tenemos que $a^p = a$ (por el pequeño teorema de Fermat).
            \item $U$ es un cuerpo:
            \begin{align*}
                \alpha, \beta \in U \text{ consideramos } \alpha^{p^n} &= \alpha, \beta^{p^n} = \beta \\
                (\alpha-\beta)^{p^n} &= \alpha^{p^n} - \beta^{p^n} = \alpha - \beta \\
                \text{Si } \beta &\neq 0 \implies \frac{\alpha^{p^n}}{\beta^{p^n}} = \frac{\alpha}{\beta}
            \end{align*}
            Y por tanto, $F \subseteq U \subseteq E = F(U) \implies E = F(U) = U$.
        \end{enumerate}
        \item[(b)] Supongamos que $L$ es otro cuerpo con $\vabs{L} = p^n$. Sea $D \subseteq L$ su cuerpo primo, $D \isom \F_p \isom F$. Ahora $L^\times = \uds L$ es un grupo de $p^n - 1$ elementos. En particular si $l \in L^\times$, entonces $l^{p^n-1} = 1$. (por \ref{thm:3.5.1}(b)) $\implies l^{p^n} = l$. Es decir, $L$ es un cuerpo de escisión de $f(x) = x^{p^n} - x$. Entonces tenemos el diagrama (que demuestra el isomorfismo):

        \begin{center}
            \begin{tikzpicture}[node distance=1cm and 2cm, auto]
                \node (A) {$E = F(f)$};
                \node (B) [right= of A] {$L=D$};
                \node (C) [below of =A] {$F$};
                \node (D) [below of =B] {$D$};
                \draw[->] (A) to node {}(B);
                \draw[->] (C) to node {$\isom$}(D);
                \draw[-] (A) to node {}(C);
                \draw[-] (B) to node {}(D);
            \end{tikzpicture}\\
        \end{center}
        con $f(n) = x^{p^n} -  x$.
    \end{itemize}
\end{proof}

\begin{obs}
    Sea $K$ un cuerpo finito, y $F$ su cuerpo primo, como consecuencia del teorema \ref{thm:3.5.2}, $K/F$ es una extensión de Galois y si $\vabs{K:F} = n$ entonces $\vabs{\gal(K/F)} = n$.
\end{obs}

%%%%%%%%%%%%%%%%%%%%%%%%%%%%%%%%%%%%%%%%%%%%%%%%%%%%%%%%%%%%%%%%%%%%%%%%%% 07/11

\begin{dfn}[Exponente de un grupo]
    Sea $G$ un grupo finito, definimos el \textbf{exponente} de $G$ como:
    $$
        \mathrm{exp}(G) = \mathrm{mcm}\sdf{o(g) \mid g \in G}
    $$
    que es el menor entero $e$ tal que:
    $$
        g^e = 1,\ \forall g \in G
    $$
\end{dfn}
\begin{obs}
    $$
        \mathrm{exp}(G) \divides \vabs{G}
    $$
\end{obs}

\begin{eg}[Ejemplos de exponentes de un grupo]$ $
    \begin{itemize}
        \item $\exp{C_2 \times C_2} = \mcm{1, 2} = 2$
        Donde $C_2 \times C_2 = \gen{a} \times \gen{b} = \sdf{1, a, b, ab}$ y tanto $o(a) = o(b) = o(ab) = 2$.
        \item $\exp{C_2 \times C_3} = \mcm{1, 2, 3, 6} = 6$.
        Donde $C_2 \times C_3 = \gen{c} \times \gen{d} = \sdf{1, c, d, d^2, cd, cd^2}$ y $o(cd^2) = o(cd) = 6$.
        \item $\exp{S_3} = \mcm{1, 2, 3} = 6$
        Donde $S_3$ es el grupo simétrico de elementos hasta $3-ciclos$.
    \end{itemize}
    En los dos primeros casos, existe un elemento cuyo orden es el mismo que el exponente del grupo, sin embargo esto no ocurre en general. Veremos que solo ocurre cuando el grupo es abeliano.
\end{eg}

\begin{lm}[Orden de elementos que conmutan entre sí]\label{lm:orden-conmut}$ $
    \begin{enumerate}
        \item Sea $H = \gen{h}$ cíclico y $\vabs{H} = n = o(h)$, para cada $d \divides n$ existe un elemento $g \in H$ tal que $d = o(g)$.
        \item Sean $g_i, g_j \in G$, si $g_i$ conmuta con $g_j$ ($g_ig_j = g_jg_i$), entonces $o(g_ig_j) \divides o(g_i) o(g_j)$. Si además, $o(g_i) \cdot o(g_j) = 1 \implies o(g_ig_j) = o(g_i) o(g_j)$
    \end{enumerate}
\end{lm}

\begin{proof}
    Se deja como ejercicio.
\end{proof}

\begin{lm}[Existencia de un elemento de orden el exponente de un grupo]\label{lm:3.5.5}
    Sea $G$ finito y abeliano, entonces $\exists g \in G$ tal que $o(g) = \exp{G}$.
\end{lm}
\begin{proof}
    Sea $e = \exp{G} = p_1^{\alpha_1} \cdots p_r^{\alpha_r}$ con $p_i$ primos y distintos entre sí y $\alpha_i \geq 1$. Por definición de exponente, para cada $j$ existe $h_j \in G$ tal que $p_j^{\alpha_j} \divides o(h_j)$.\\
    Por el lema \ref{lm:orden-conmut}(1), $\exists g_j \in G$ tal que $o(g_j) = p_j^{\alpha_j}$. Tomamos:
    $$
        g = \prod_{j=1}^r g_j \in G,\ \text{ y por el lema \ref{lm:orden-conmut}(2) } o(g) = e
    $$
\end{proof}

\begin{thm}[Subgrupo multiplicativo y subgrupo cíclico]\label{thm:3.5.6}
    Sea $K$ un cuerpo y $G \leq K^\times$ un subgrupo finito, entonces $G$ es cíclico. Enunciado informalmente:\\
    \begin{center}
        \textit{Todo subgrupo finito del grupo multiplicativo de un cuerpo es cíclico}
    \end{center}
\end{thm}

\begin{proof}
    Por el teorema \ref{thm:3.5.1}(d), queremos encontrar $g \in G$ tal que $o(g) = \vabs{G}$.\\
    Como $K^\times$ es abeliano y $G \leq K^\times$, entonces $G$ es abeliano (y finito por hipótesis). Sea $e = \exp G$, $\forall h \in G,\ h^e = 1$, entonces todo elemento $h \in G$ es raíz del polinomio $x^e - 1 \in K[x]$. Por tanto:
    $$
        \vabs{G} \leq e \leq \vabs{G} \text{ pues $e \divides \vabs{G}$ } \implies e = \vabs{G}
    $$
    Por el lema \ref{lm:3.5.5}, $\exists g \in G$ tal que $o(g) = e = \vabs{G}$.
\end{proof}

\begin{obs}
    Sea $K$ un cuerpo finito, en particular $K^\times$ es cíclico, luego $\exists \xi \in K\ :\ \gen{\xi} = K^\times$.
\end{obs}

\begin{ex}[H3.19]
    Sea $E = \quo{\F_3[x]}{\gen{x^2 + 1}}$ (en particular $x^2 = -1 = 2$), halla un generador de $E^\times$.\\\\

    \textbf{Solución}\\
    Vemos que:
    $$
        E = \sdf{0, 1, 2, x, 2x, x+1, x+2, 2x+1, 2x+2},\ \vabs{E} = 9,\ \vabs{E^\times}=8
    $$
    Buscamos un elemento $g \in E^\times$ tal que $o(g) = 8$. No hay un método general, solo ensayo y error.\\
    Vamos a probar algunos casos:
    \begin{itemize}
        \item[$(x)$]
        $$
            x^1 = x,\ x^2 = 2,\ x^3 = x\cdot x^2 = 2x,\ x^4 = 2 \cdot 2 = 1 \implies o(x) = 4 \implies x \text{ no es un generador.}
        $$
        \item[$(x+1)$]
        \begin{align*}
            (x+1)^1 &= x+1\\
            (x+1)^2 &= x^2 + 2x + 1 = 2x\\
            (x+1)^3 &= 2x(x+1) = 2x^2 + 2x = 2x + 1\\
            (x+1)^4 &= (2x)^2 = x^2 = 2\\
            (x+1)^5 &= 2x+2\\
            (x+1)^6 &= 2(2x) = x\\
            (x+1)^7 &= x(x+1) = x^2 + x = x + 2\\
            (x+1)^8 &= (x+2)(x+1) = x^2 + 2 = 1
        \end{align*}
    \end{itemize}
    \textit{Ejercicio extra}\\
    Halla o calcula:
    \begin{enumerate}
        \item $E/\F_3$ es simple.
        \item $E$ es el cuerpo de escisión de $y^2 + 1$ sobre $\F_3$.
        \item $E$ es el cuerpo de escisión sobre $\F_3$ de cualquier polinomio irreducible de grado $2$ en $\F_3[x]$.
    \end{enumerate}
    \textbf{Solución}\\ %TODO: Completar el ejercicio (2 últimos apartados)
    \begin{enumerate}
        \item Hemos visto que $\exists \xi \in E$ tal que $\gen{\xi} = E^\times$. Entonces $\sdf{\xi, \xi^2, \ldots, \xi^n} \subset E = \F_3(\xi)$ y por tanto $E/\F_3$ es simple.
    \end{enumerate}
\end{ex}

\begin{cor}[Teorema del Elemento Primitivo sobre cuerpos finitos]\label{cor:3.5.7}
    Sea $E/K$ una extensión finita sobre un cuerpo finito $K$, entonces $E/K$ es simple.
\end{cor}
\begin{proof}
    Aunque en el teorema del Elemento Primitivo (teorema \ref{thm:elemprim}) se pide que la extensión sea separable, en el corolario no hace falta pedirlo.\\
    Como $K$ es finito, entonces $K$ es perfecto y cualquier extensión algebraica sobre un cuerpo finito es automáticamente separable. Además, $E/K$ es algebraica por ser finita (\ref{thm:2.2}). Para ver que es finito:
    $$
        \vabs{E} = \vabs{K}^{\vabs{E:K}} < \infty, \implies E \text{ es finito}
    $$
    Además por el teorema \ref{thm:3.5.6}, $\exists \xi \in E:\ \gen{\xi} = E^\times \implies K(\xi) = E$.
\end{proof}

\begin{ex}[H3.17]
    Demuestra que para cada primo $p$ y cada natural $n \in \N$ positivo, existe un polinomio irreducible $f \in \F_p[x]$ de grado $n$.\\\\
    \textbf{Solución}\\
    Por el teorema \ref{thm:3.5.2}(a), $\exists$ un cuerpo $E$ tal que $\vabs{E} = p^n$.
    Sea $F \subseteq E$ su cuerpo primo, $F = \F_p$. Por \ref{cor:3.5.7} $\exists \xi \in E$ tal que $E = \F_p(\xi)$.\\
    Por el teorema \ref{thm:2.3}:
    $$
        n = \vabs{E:\F_p} = \delta Irr(\F_p, \xi) = \delta f(x)
    $$
    con lo que $f$ es un irreducible de grado $n$.
\end{ex}

\begin{ex}[H3.9]
    Sea $E = \Q(\sqrt[4]{2})$, $L = \Q(\sqrt 2)$, $K = \Q$. Comprueba que $E/K$ no es normal.\\\\

    \textbf{Solucion}\\
    Sabemos que $E/L$ y $L/K$ son extensiones normales, ya que $E = L(x^2-\sqrt{2}) = L(\pm \sqrt[4]{2}) = \Q(\sqrt[4]{2})$ y $L = K(x^2-2) = \Q(\pm \sqrt{2})$. Sin embargo, $E/K$ no es normal porque el polinomio irreducible sobre $K$ ($x^4-2$) tiene dos raíces en $E$ pero no se escinde en $E$. (teorema \ref{thm:3.9})
\end{ex}

\begin{ex}[Ejercicio de examen]
    Sea $f(x) = x^4 - 2 \in \Q[x]$. Determina:
    \begin{itemize}
        \item[(a)] $E=\Q(f)$.
        \item[(b)] $\vabs{E:\Q}$.
        \item[(c)] $\alpha \in E$ tal que $\Q(\alpha) = E$.
        \item[(d)] Decide si $\sigma: \Q(\sqrt 2) \to \Q(\sqrt 2);\ \sqrt{2} \mapsto -\sqrt{2}$ se puede extender a un isomorfismo de $\Q(\sqrt[4]{2})$ o de $E$.
    \end{itemize}

    \textbf{Solución}\\
    \begin{itemize}
        \item[(a)] Las raíces de $f$ son $\sdf{\pm \sqrt[4]{2}, \pm \sqrt[4]{2}i}$ entonces $E = \Q(\pm \sqrt[4]{2}, \pm \sqrt[4]{2}i) = \Q(\sqrt[4]{2}, i)$.
        \item[(b)] Sabemos que $\Q \subseteq \Q(i), \Q(\sqrt[4]{2}) \subseteq E$. Vamos a hallar el grado de la extensión $E/\Q(\sqrt[4]{2})$, que resulta ser $2 = \delta Irr(\sqrt[4]{2}, i) = \delta(x^2 + 1)$.\\
        Por otra parte calculamos $\delta Irr(\Q, \sqrt[4]{2}) = \delta(x^4-2) = 4$. Por tanto:
        $$
            \vabs{E:Q} = 2 \cdot 4 = 8
        $$
        \item[(c)] Hay dos formas:\\
            \begin{enumerate}
                \item Usando la demostración del Teorema del Elemento Primitivo (TEP) en el caso infinito:\\
                $E = \Q(i = \alpha, \sqrt[4]{2} = \beta)$, entonces $$\alpha_1 = i, \alpha_2 = -i, \beta_1 = \sqrt[4]{2}, \beta_2 = -\sqrt[4]{2}, \beta_3 = \sqrt[4]{2}i, \beta_4 = -\sqrt[4]{2}i$$.
                Tenemos que escoger $c \in \Q$ tal que:
                $$
                    c \neq \frac{\beta - \beta_j}{\alpha-\alpha_2} \text{ con } j \in \sdf{1, 2, 3, 4},\ \beta = \beta_1,\ \alpha = \alpha_1
                $$
                podemos ver que $c = -1$ funciona. Por la prueba del TEP: $$E = \Q(c\alpha - \beta) = \Q(-\alpha - \beta) = \Q(\alpha + \beta) = \Q(\sqrt[4]{2} + i)$$

                \item Usando todo lo que hemos visto:
                Sea $L = \Q(i)$, sabemos por el apartado (b) de este ejercicio que $\vabs{E:Q} = 8$, y además, $\vabs{L:\Q} = 2$ por lo que $\vabs{E:L} = 4 \implies x^4-2$. Sea $\alpha = \sqrt[4]{2} + i \in E$, por el teorema del elemento algebraico queremos ver que $\delta Irr(\Q, \alpha) = 8$.\\\\
                Además, $Irr(L, \sqrt[4]{2}) = x^4 - 2 = f$, y por tanto $p(x) = (x-i)^4-2 = f(x-i)$ donde $\alpha$ es una raíz de $p \in L[x],\ p \notin \Q[x]$. Por el ejercicio H1.32(a), $p$ es irreducible sobre $L$, es decir, $p(x) = (x-i)^4-2 = Irr(L, \alpha)$.
                %%TODO Pedir a Santorum.
            \end{enumerate}
        \item[(d)]
            Recordamos que si $E/K$ es normal, y $M/K$ es normal con $K \subseteq M \subseteq K$, entonces todo $\sigma \in \gal(M/K)$ se puede extender a $\gal(E/K)$ (Teorema \ref{thm:3.11} y corolario \ref{cor:3.12}).\\\\
            Sea $M = \Q(\sqrt[4]{2})$, podemos ver que $E/\Q$ es normal, $\Q(\sqrt{2})/\Q$ también lo es, pero $M/\Q$ no es normal. Por tanto, $\sigma$ no se puede extender a un isomorfismo de $M$. Podemos ver por reducción al absurdo:\\
            Sea $\tau: M \to M$ una extensión de sigma, entonces $\tau(\sqrt[4]{2}) = \sqrt[4]{2}$ o $\tau(\sqrt[4]{2}) = -\sqrt[4]{2}$ ya que $\tau$ lleva raíces sobre $\Q$ en raíces sobre $\Q$ entonces:
            $$
                -\sqrt{2} = \sigma(\sqrt 2) = \tau(\sqrt 2) = \tau((\sqrt[4]{2})^2) = \tau(\sqrt[4]{2})^2 = \pm \sqrt{2}
            $$
            y llegamos a una contradicción.\\\\

            Para ver que $\sigma$ se extiende a un isomorfismo $\omega$ de $E$ basta aplicar el corolario \ref{cor:3.12}, ya que las dos extensiones ($\Q(\sqrt 2)/\Q)$ y $E/\Q$ son normales).\\
            Para verlo de otra forma, sea $\omega \in \gal(E/\Q)$ tal que $\omega(\sqrt[4]{2}) = \sqrt[4]{2} i$ (sabemos que existe uno por que tanto $\sqrt[4]{2}$ como $\sqrt[4]{2} i$ son raíces del mimso polinomio irreducible) entonces:
            $$
                \omega(\sqrt 2) = \omega((\sqrt[4]{2})^2) = \omega(\sqrt[4]{2})^2 = (\sqrt[4]{2} i)^2 = -\sqrt{2} = \sigma(\sqrt 2)
            $$
    \end{itemize}
\end{ex}

\begin{ex}[H3.10]
    Todas las afirmaciones de este ejercicio son ciertas.
    \begin{itemize}
        \item[(f)] Sean $E/L$ y $L/K$ extensiones normales. Si todo $\sigma \in \gal(L/K)$ se extiende a $E$ ($\gal(E/K)$) entonces $E/K$ es normal.\\\\
        \textit{Nota}\\
        Este resultado se puede usar para ver que $\sigma \in \gal(\Q(\sqrt 2)/\Q)$ con $\sigma(\sqrt 2) = -\sqrt 2$ no se puede extender a $M = \Q(\sqrt[4]{2})$, por que si no $M/\Q$ sería normal y ya vimos que no lo es.
    \end{itemize}
\end{ex}

\begin{ex}[H3.4]
    Sea $K = \quo{\F_2[y]}{y^3 + y + 1}$, comprueba que es un cuerpo de escisión de $x^3 + x + 1$ y $x^3 + x^2 + 1$.\\\\
    \textbf{Solución}\\
    Podemos ver que $K = \sdf{0, 1, y, y+1, y^2, y^2+1, y^2+y, y^2+y+1}$. Por el lema \ref{lm:3.3}, sabemos que $y \in K$ es raíz de $x^3+x+1$. También sabemos que si $\sigma \in \gal(K/\F_2)$ entonces $\sigma(y)$ es raíz de $x^3+x+1$.\\
    Tomamos $\varphi = Frob \in \gal(K/\F_2)$ y $\varphi^2 \in \gal(K/\F_2)$. Entonces $\varphi(y) = y^2$ es otra raíz, y $\varphi^2(y) = \varphi(y^2) = y^4 = y \cdot y^3 = y^2+y$.\\
    Sabemos que $x^3+x+1$ y $x^3+x^2+1$ no comparten raíces en $K$. Comprobamos que $y+1$ es raíz de $x^3+x^2+1$ y vemos que $\varphi(y+1)=y^2+1$ y $\varphi(y^2+1) = y^2+y+1$ por lo que claramente $\F_2(y) = K$ y $\F_2(y+1)$
\end{ex}

\begin{ex}[H3.15]
    ¿Cuántas raíces tiene distintas $f(x) = x^{12} + 2x^6 + 1 \in \F_3[x]$ en su cuerpo de escisión?\\\\
    \textbf{Solución}\\
    Empezamos comprobando si tiene raíces en $\F_3$.
    \begin{align*}
        f(0) &= 0 + 0 + 1 \neq 0\\
        f(1) &= 1 + 2 + 1 \neq 0\\
        f(2) &= 1 + 2 + 1 \neq 0\\
    \end{align*}
    Y vemos que no, sin embargo esto no quiere decir que el polinomio sea irreducible. De hecho, como $f(x) \in \F_3[x^3]$ es reducible. Usando un poco la intuición podemos observar que $f(x)$ es el desarrollo del producto notable:
    $$
        f(x) = (x^6 + 1)^2
    $$
    Sin embargo, este polinomio sigue sin ser irreducible, ya que $x^6 + 1 \in \F_3[x^3]$. Además como estamos en un cuerpo de característica $p$, sabemos que $(x^k + 1)^p = x^{k\cdot p} + 1$ por que los términos intermedios tienen todos un coeficiente que es múltiplo de $3$. Por tanto:
    $$
        f(x) = x^{12} + 2x^6 + 1 = (x^6+1)^2 = ((x^2+1)^3)^2 = (x^2+1)^6
    $$
    Ahora consideramos $g(x) = x^2 + 1$, como $g(x)$ es irreducible sobre $\F_3$, sabemos que $\F_3(g)/\F_3$ es una extensión separable y por tanto $g(x)$ tiene dos raíces distintas en su cuerpo de escisión lo que implica que $f$ también.\\\\
    Por otra parte, si queremos hallar la descomposición de $f$ en el cuerpo de escisión, consideramos $\quo{\F_3[t]}{\gen{t^2+1}} = \quo{\F_3[t]}{\gen{g(t)}}$ que es un cuerpo con $9 = \vabs{\F_3} \cdot \delta(t^2+1)$ elementos. Ahora tendremos que ver que raíces del polinomio $\F_3[x]$ están presentes en $\quo{\F_3[t]}{\gen{t^2+1}}$.\\
    $$
        \quo{\F_3[t]}{\gen{t^2+1}} = \sdf{0, 1, 2, t, 1+t, 2+t, 2t, 1+2t, 2+2t}
    $$
    Como estamos cocientando por $t^2+1$, sabemos que $t^2+1 = 0 \implies t^2 = -1$ y por tanto $t$ es raíz de $x^2+1$. Además, $(2t)^2 = 4t^2 = -1$ y por tanto también es raíz de $x^2+1$. Entonces podemos descomponer de la forma:
    $$
        x^2 + 1 = (x-t)\cdot(x-2t)
    $$
    Por tanto, tendríamos:
    $$
        f(x) = (x-t)^6(x-2t)^6
    $$
\end{ex}
\begin{ex}[H3.18]
    Sea $f(x) = x^q - x \in \F_p[x]$ con $q = p^n$. Demuestra que:
    \begin{itemize}
        \item[(a)] Cualquier polinomio $g(x)$ irreducible en $\F_p[x]$ de grado $n$ divide a $f$.
        \item[(b)] El grado de todos los factores irreducibles de $f$ divide a $n$.
    \end{itemize}
    \textbf{Solución}\\
    \begin{itemize}
        \item[(a)] Sabemos que por ser $f(x)$ irreducible de grado $q$ entonces:
        $$
            K = \F_p(f) \isom \F_q = \sdf{\alpha_1, \ldots, \alpha_q}
        $$
        entonces $f(x) = \prod_{i=1}^{q} (x-\alpha_i)$.\\
        Por el teorema de Kronecker (lema \ref{lm:3.3}), sabemos que $L = \quo{\F_p[t]}{\gen{g}}$ es una extensión de $\F_p$ de grado $n$, que contiene una raíz de $g \in L[x]$. Además, $L$ contiene $p^{\delta(g)} = p^n = q$ elementos y por tanto es isomorfo a $K$ (por el teorema \ref{thm:3.5.2}).\\
        También sabemos que existe un isomorfismo que fija $\F_p$ y por tanto, como $g$ no cambia por el isomorfismo, también tiene una raíz en $K$.
        Es decir, $\exists \alpha K$ con $(x - \alpha) \divides g(x)$.\\
        % Además $K = \F_p(f) / \F_p$ es un cuerpo de escisión por lo que es una extensión normal. Como $g$ tiene una raíz en $K$ que es extensión normal $\implies$ $g$ se escinde en $K$. Por tanto:
        % $$
        %     g(x) = \prod_{i=1}^n (x-\alpha_i)
        % $$
        % y además:
        % $$
        %     f(x) = g \cdot \prod_{n+1}^q (x-\alpha_i)
        % $$
        % Como tenemos un algoritmo de la división en $\F_p[x]$, lo que nos dice que $\prod_{n+1}^q (x-\alpha_i) \in \F_p[x]$.
        Sabemos que $(x-\alpha)\divides g(x)$ y $(x-\alpha)\divides f(x)$ lo que quiere decir que $(x-\alpha)\divides mcd_{K[x]}(f, g) = mcd_{\F_p[x]}(f, q)$ y como $g(x)$ es irreducible, entonces $g(x) \divides f(x)$.
        \item[(b)]
        Sea $g \in \F_p[x]$ factor irreducible con $\delta(g) = d$. Sea $K = \F_p(f)$. Como $f$ se escinde en $K$ es obvio que $g$ se escinde en $K$. Sea $\alpha$ una raíz cualquiera de $g$. Vamos a construir un cuerpo intermedio $L$ de forma $\F_p \subseteq L = \F_p(\alpha) \subseteq K$. Sabemos que $\vabs{K:\F_p} = n$, y que por la elección de $\alpha$, $\vabs{L:\F_p} = d$. Entonces por transitividad de grados:
        $$
            \delta(g) = d \divides n
        $$
    \end{itemize}
\end{ex}

% \part{Tercer parcial}
% % !TeX root = ../galois.tex

\chapter{Teoría de Galois}

\section{Introducción}

A partir de este capítulo vamos a trabajar siempre con extensiones $E/K$ que sean de Galois, es decir, $E/K$ es separable y normal. En el teorema \ref{thm:3.11}, caracterizamos al grupo de Galois de una extensión, vamos a ver que ambas hipótesis (separabilidad y normalidad) son necesarias.

\begin{eg}[Grupo de Galois de una extensión no normal]
    Consideremos $\Q(\sqrt{1 + \sqrt 7})/\Q$. Podemos hallar que $\sqrt{1 + \sqrt 7}$ es raíz de $p(x) = x^4-2x^2-6 \in \Q[x]$, que es irreducible por el criterio de Einsestein para $p=2$.\\
    Entonces, las raíces de $p(x)$ son $\pm\sqrt{1 + \pm\sqrt 7}$, y su cuerpo de escisión es $E = \Q(p) = \Q(\sqrt{1 + \sqrt 7}, \sqrt{1 - \sqrt 7})$.\\
    Si suponemos que $\Q(\sqrt{1 + \sqrt 7}) = E(p)$ rápidamente llegamos a una contradicción con lo que concluimos que $\Q(\sqrt{1 + \sqrt 7})/\Q$ no es normal. Nos  preguntamos ahora acerca de $\gal(\Q(\sqrt{1 + \sqrt 7})/\Q)$.\\
    Por el teorema \ref{thm:3.11}, $\Omega = \sdf{\pm \sqrt{1 + \sqrt 7}}$ es el conjunto de raíces de $p$ en $\Q(\sqrt{1 + \sqrt 7})/\Q$, y por tanto $\gal(\Q(\sqrt{1 + \sqrt 7})/\Q) \leq S_2 \implies \gal(\Q(\sqrt{1 + \sqrt 7})/\Q) = \sdf{id, \sigma}$ con $\sigma(\sqrt{1 + \sqrt 7}) = -\sqrt{1 + \sqrt 7}$.\\\\
    De esta forma vemos que:
    $$
        \vabs{\gal(\Q(\sqrt{1 + \sqrt 7})/\Q)} = 2 \text{ pero } \vabs{\Q(\sqrt{1 + \sqrt 7}):\Q} = 4
    $$
\end{eg}

\begin{ex}[H4.3]$ $
    \begin{itemize}
        \item[(a)] Determina que $\F_3(t) / \F_3(t^3)$ no es de Galois.
        \item[(b)] Determina que $\C(t) / \C(t^3)$ es de Galois.
        \item[(c)] Calcula los grupos de Galois asociados.
    \end{itemize}
    \textbf{Solución}\\
    \begin{itemize}
        \item[(a)] Vamos a ver que la extensión no es separable. Sabemos que $\F_3(t) / \F_3(t^3)$ es algebraica y finita. Consideramos el polinomio irreducible $x^3-t^3 \in \F_3(t^3)[x]$, que tiene $t$ como raíz, entonces no es separable por que $x^3-t^3 = (x-t)^3$ y por tanto no es de Galois.\\
        Además, $\gal(\F_3(t)/\F_3(t^3)) = \sdf{id}$.
        \item[(b)] Sabemos que $\C(t) / \C(t^3)$ es separable por que estamos en un cuerpo de característica cero. Además, sea el polinomio $x^3-t^3 \in \C(t^3)[x]$, sus raíces son $\Omega = \sdf{t, t\omega, t\bar\omega}$ donde $\omega \in \C:\ o(\omega)=3$.\\
        Como $\Omega \subseteq \C(t)$, $\C(t)$ es el cuerpo de escisión de $p(x)$ y por tanto $\C(t)/\C(t^3)$ es separable y normal y por tanto de Galois.
        $$
            \text{Luego, } \vabs{\gal(\C(t) / \C(t^3))} = \vabs{\C(t) : \C(t^3)} = 3
        $$
        Con lo que hallamos:
        $$
            \gal(\C(t) / \C(t^3)) = \gen{\tau}:\ \tau(t) = t\omega,\ o(\tau) = 3
        $$
        \item[(c)] Resuelto en los apartados respectivos.
        \end{itemize}
\end{ex}

\begin{ex}[H4.2]
    Sea $E = \F_{p^n}$, demuestra que $E / \F_p$ es de Galois y calcula $\gal(E / \F_p)$.\\\\

    \textbf{Solución}\\
    Sabemos que $E / \F_p$ es separable por que $\F_p$ es perfecto y la extensión es algebraica. Tenemos que ver que $E / \F_p$ es normal.\\
    Además $E$ es el cuerpo de escisión de $x^{p^n} - x$ sobre $\F_p \implies E / \F_p$ es normal y por tanto es de Galois.\\\\
    Como $E / \F_p$ es de Galois $\implies \vabs{\gal(E / \F_p)} = \vabs{E : \F_p} = n$. Sea $\varphi = Frob \in \gal(E / \F_p)$, sabemos que $o(\varphi) \divides n$.
    Sea $d = o(\varphi)$, sabemos que $d \leq n$ y que:
    $$
        a = id(a) = \varphi^d(a) = a^{p^d} \forall\ a \in E \implies \forall a \in E,\ a^{p^d} - a = 0
    $$
    Por tanto, todo elemento en $E$ es raíz de $x^{p^d} - x \implies E$ es el cuerpo de escisión de $x^{p^d} - x$. Esto quiere decir que $E$ tiene $\delta(x^{p^d} - x)$ elementos, pero $E = \F_{p^n}$ que tiene $p^n$ elementos $\implies d = n$. Por tanto:
    $$
        \gal(E / \F_p) = \gen{\varphi}
    $$
\end{ex}

\begin{cor} % TODO: Dar nombre
    Sea $E/K$ una extensión de Galois y $L$ un cuerpo intermedio ($K \subseteq L \subseteq E$), entonces:
    $$
        L/K \text{ es de Galois } \iff \forall \tau \in \gal(E/K),\ \rest{\tau}_L \in \gal(E/K)
    $$
    Y además, en cualquiera de los supuestos equivalentes anteriores, todo $\sigma \in \gal(L/K)$ se extiende a $\tau \in \gal(E/K)$ (tiene exactamente $\vabs{E:L}$ extensiones), $\gal(E/L) \normsub \gal(E/K)$ y $ \quo{\gal(E/K)}{\gal(E/L)} \isom \gal(L/K)$.\\\\
    Podemos hacer un diagrama:
    % TODO
    % E          {1} = \gal(K/K)
    % |          |
    % |          |
    % L -------- N = \gal(E/L)
    % |          ||
    % |          ||
    % K          G = \gal(L/K)
\end{cor}

El objetivo de este capítulo es probar el teorema \ref{thm:fund-gal}:

Para ello vamos a hacer uso del corolario \ref{cor:3.12}. Además, necesitaremos un resultado sobre automorfismos de cuerpos en general.\\
Recordemos que si $E$ es un cuerpo, entonces $\aut(E)$ actúa sobre $E$ según:
$$
    \forall e \in E,\ \forall \sigma \in \aut(E),\ \ e \cdot \sigma = \sigma(e)
$$
y cumple las propiedades:
\begin{align*}
    e \cdot (\sigma \tau) = (e \cdot \sigma) \tau\\
    e \cdot (\sigma \tau) = \sigma \circ \tau (e) = \tau ( \sigma (e)) = \tau(e \cdot \sigma) = (e \cdot \sigma) \tau
\end{align*}
Además, si $G \leq \aut(E)$, entonces $G$ actúa sobre $E$, y recordemos la notación del conjunto de puntos fijos de $E$ por la acción de $G$:
$$
    E^G = \sdf{e \in E \mid e \cdot \sigma = e} = \sdf{e \in E \mid \sigma(e) = e} \subseteq E \text{ es un subcuerpo }
$$
Por tanto $E/E^G$ es una extensión.

\begin{thm}[de Artin]\label{thm:4.1}
    Sea $E$ un cuerpo, y sea $G \leq \aut(E)$ un subgrupo finito, entonces:
    $$
        \vabs{E:E^G} = \vabs{G}
    $$
\end{thm}
\begin{proof}
    Se proporciona un esquema de la prueba.\\
    Sea $F = E^G$, la prueba sigue por demostrar la doble igualdad:
    \begin{enumerate}
        \item $m = \vabs{E:F} \geq \vabs{G} = n$.\\
        Por reducción al absurdo suponemos $m < n$, entonces podemos encontrar $\tau_1, \ldots, \tau_n \in G$ distintos tal que:
        $$
            \sum_{j=1}^n a_j\tau_j = 0
        $$
        Y esto entra en contradicción con el teorema de Dedekind (teorema \ref{thm:4.2}). Entonces $m \geq n$.
        \item $\vabs{E:F} = m \leq n = \vabs{G}$\\
        Tendríamos que trabajar con sistemas lineales, soluciones y manipulaciones...
    \end{enumerate}
\end{proof}

\begin{thm}[de Dedekind]\label{thm:4.2}
    Sea $E$ un cuerpo y sean $\tau_1, \ldots, \tau_n \in \aut(E)$ distintos. Entonces $\tau_1, \ldots, \tau_n$ son $E-$linealmente independientes (en el espacio vectoreial de funciones de $E$).
\end{thm}
\begin{proof}
    Queremos probar que si $a_j \in E$:
    $$
        \sum_{j=1}^{n} a_j\tau_j = 0 \implies a_j = 0 \forall j
    $$
    Procedemos por inducción sonre $n$.
    \begin{itemize}
        \item[$n=1$]
        $$
            a_1\tau_1 = 0 \implies a_1\tau_1(1) = a_1 = 0
        $$
        \item[$n>1$]
        \begin{align}
            \forall x  \in E:\ \sum_{j=1}^n a_j\tau_j(x) = 0\\
            \forall y \in E:\ \sum_{j=1}^n a_j\tau_j(xy) = \sum_{j=1}^n a_j\tau_j(x)\tau_j(y) = 0\\
            \forall y \in E:\ \sum_{j=1}^n a_j\tau_j(x)\tau_n(y) = 0
        \end{align}
        Entonces restando:
        $$
            0 = \sum_{j=1}^n a_j(\tau_j(y) - \tau_n(y)) \tau_j(x) = \sum_{j=1}^{n-1} a_j(\tau_j(y)-\tau_n(y))\tau_j(x)
        $$
        Por inducción:
        $$
            a_j(\tau_j(y) - \tau_n(y)) = 0 \forall j \in \sdf{1, \ldots, n-1}
        $$
        Supongamos que $\exists a_j \neq 0$, entonces:
        $$
            \tau_j(y) = \tau_n(y) \implies \tau_j = \tau_n \text{ y es imposible por hipótesis}
        $$
    \end{itemize}
\end{proof}

\begin{thm}[de puntos fijados]\label{thm:4.3}
    Sea $E/K$ una extensión de Galois, y $G = \gal(E/K)$ su grupo de Galois $\implies E^G = K$.\\
    Es decir, si $a \in E$ tal que $\forall \tau \in G:\ \tau(a) = a \implies a \in K$.
\end{thm}
\begin{proof}
    Escribimos $F = E^G \subseteq E$. Suponemos $K \subseteq F$, tenemos entonces el diagrama:
    %TODO
    %   |-- E ----------|
    %   |   |           | |G|
    %|G||   F = E^G ----|
    %   |   |           | 1
    %   |-- K ----------|
    Por el teorema \ref{thm:3.4.6}, sabemos que $\vabs{G} = \vabs{E:K}$. Por el teorema de Artin (teorema \ref{thm:4.1}), $\vabs{E:E^G} = \vabs{E:F} \implies |F:K| = 1 \implies K = F = E^G$.
\end{proof}
\begin{obs}
    El recíproco también es cierto, se corresponde con el ejercicio H4.4. Basta ver que $\forall \alpha \in E,\ Irr(K, \alpha)$ se escinde en $E$ y todas sus raíces son distintas.
\end{obs}
\begin{cor}\label{cor:4.4}%TODO Dar nombre
    Sea $E/K$ una extensión de Galois, $G = \gal(E/K)$ su grupo de Galois y $L$ un cuerpo intermedio entonces:
    $$
        L/K \text{ es normal } \iff \gal(E/L) \normsub G
    $$
\end{cor}
\begin{proof}$ $
    \begin{itemize}
        \item[($\implies$)] Por el corolario \ref{cor:3.12}(b), entonces:
        $$
            (\star)\ L/K \text{ normal } \iff \forall \tau \in G: \tau(L) = L
        $$

        \item[($\Longleftarrow$)] Partimos de $\gal(E/L) \normsub G$, y queremos ver que $L/K$ es normal. Usando $(\star)$, basta ver que $\forall \tau \in G$, $\tau(L) = L$. Entonces:
        $$
            H \normsub G \iff \forall \sigma \in H,\ \forall \tau \in G,\ \sigma^\tau = \tau^{-1} \sigma \tau \in H
        $$
        entonces $\forall l \in L$:
        \begin{align*}
            \sigma^\tau(l) = \tau^{-1} \sigma \tau(l) = l\cdot \tau^{-1} \cdot \sigma \cdot \tau = l\\
            l \cdot (\tau^{-1} \cdot \sigma) = l \cdot \tau^{-1} \implies \tau^{-1}(l) \in E^H = L
        \end{align*}
        Y por tanto $E/L$ es de Galois por el teorema \ref{thm:4.3}. Como $E/L$ es de Galois $\implies \tau^{-1}(L) \subseteq L \implies L \subseteq \tau(L)$.\\
        Entonces, tenemos el diagrama:
        %   tau rest L : L \to (isom) tau(L)
        %           |                   | 1
        %           |                   |L cuerpo intermedio
        %           K \to ------------> K
        %
        Y por el ejercicio H2.15, $\vabs{L:K} = \vabs{\tau(L):K} = \vabs{\tau(L):L}\vabs{L:K} \implies \tau(L) = L$.
        Además, también tenemos el diagrama:
        %        E
        %      /   \
        %  \tau(L)  L
        %      \   /
        %        K
        Entonces:
        \begin{align*}
            H = \gal(E/L) \leq G = \gal(E/K),\ \tau \in G\\
            H \isom H^\tau = \tau^{-1} H \tau = \sdf{\tau^{-1} \sigma \tau \mid \sigma \in H}\\
            H^\tau = \gal(E/\tau(L))
        \end{align*}
    \end{itemize}
\end{proof}

\begin{obs}[Notación]
    Recordamos que usamos la notación de acción de grupos a derechas, es decir:\\
    Sean $f,\ g$ funciones $f,g: E \to E$ con $E$ un cuerpo, $f \circ g (e) = g (f (e))$ y sea el grupo de automorfismos $(\aut(E), \circ)$, entonces para $\sigma, \tau \in  \aut(E)$ tenemos que $\forall e \in E$:
    $$
        e \cdot \sigma = \sigma(e),\text{ y } e \cdot (\sigma \tau) = (e \cdot \sigma) \tau
    $$
\end{obs}

\begin{lm}[Lema]\label{lm:4.5} %%TODO: DAR NOMBRE
    Sean $K \subseteq L \subseteq E$ cuerpos, $H = \gal(E/L) \leq \gal(E/K)$, y sea $\tau \in G$ entonces:
    $$
        H^\tau = \gal(E/\tau(L))
    $$
\end{lm}
\begin{proof}$ $
    \begin{align*}
        %% TODO: Pedir a Santorum
    \end{align*}
\end{proof}

\section{Teorema fundamental de la teoría de Galois}
\begin{thm}[Fundamental de la Teoría de Galois]\label{thm:fund-gal}\label{thm:4.6}
    Sea $E/K$ una extensión de Galois con $G = \gal(E/K)$, y sean:
    $$
        \mathcal{K} = \sdf{L \mid K \subseteq L \subseteq E},\ \mathcal{J} = \sdf{H \mid H \leq G}
    $$
    Entonces:
    \begin{itemize}
        \item[(a)] Las aplicaciones:
        $$ f: \mathcal{K} \to \mtc{J};\ L \mapsto \gal(E/L)$$
        $$ g: \mtc{J} \to \mtc{K};\ H \mapsto E^H $$
        son biyectivas y $f^{-1} = g,\ g^{-1} = f$.
        \item[(b)] Bajo esta correspondencia:
        $$
            L/K \text{ es normal } \iff \gal(E/L) \normsub G
        $$
        \item[(c)] Si $L \longleftrightarrow H$, $\vabs{L:K}$ (grado) $ = \vabs{G:H}$ (índice)
    \end{itemize}
\end{thm}
\begin{proof}$ $
    \begin{itemize}
        \item[(a)] Sea $L$ un cuerpo intermedio, como $E/L$ es una extensión de Galois, si $H = f(L) = \gal(E/L)$ entonces:
        $$
            (f \circ g)(L) = g(f(L)) = g(H) = E^H =^{\text{(teorema \ref{thm:4.3})}} = L \implies (f \circ g) = id_{\mtc K}
        $$
        En particular, $g$ es sobreyectiva, si probamos que es inyectiva, entonces $g^{-1} = f$.\\
        Sean $H_1, H_2 \normsub G$ tales que:
        $$
            L = E^{H_1} = g(H_1) = g(H_2) = E^{H_2}
        $$
        Por el teorema de Artin (teorema \ref{thm:4.1}), entonces:
        \begin{align*}
            \vabs{E:E^{H_1}} &= \vabs{H_1}\\
            \vabs{E:E^{H_1}} &= \vabs{E:L} = \vabs{E:E^{H_2}} = \vabs{H_2}\\
            \vabs{E:L} &= \vabs{E:E^{\gal(E/L)}} = \vabs{\gal(E/L)}
        \end{align*}
        Como $E/L$ es de Galois, entonces
        % TODO: Pedir a Santorum
    \end{itemize}
\end{proof}

\begin{ex}[H4.5]
    Calcular el grupo de Galois del cuerpo de escisión de $x^12 - 1$ sobre $\Q$ y todas las subextensiones intermedias.\\\\
    \textbf{Solución}\\
    Comenzamos viendo que:
    \begin{align*}
        E &= \Q(x^12-1) = \Q(x^6+1)\\
          &=\Q(\pm 1, \pm i \pm \frac{\sqrt 3}{2} \pm \frac{1}{2}i, \pm\frac{1}{2} + \frac{\sqrt 3}{2})\\
          &= \Q(\sqrt 3, i)
    \end{align*}
    Además, podemos comprobar que $\vabs{E:\Q} = 4 = \vabs{G}$ y como $\vabs{G} = 2^m$, entonces $G = \gal(E/\Q)$ es normal.\\
    Recordamos que si un grupo $G$ es cíclico, entonces $\forall d \divides \vabs{G}$ existe un único subgrupo de índice $d$.\\
    Además, $G$ tiene dos grupos distintos de índice $2 \implies G = C_2 \times C_2$.
    Por el teorema \ref{thm:4.6}, existen tantas subextensiones como $H \leq G \isom C_2 \times C_2$. Es decir, tantos como subgrupos tiene $C_2 \times C_2$, que podemos calcularlos y hallar que son $5$.
\end{ex}
\begin{obs}
    Si nos fijamos en la prueba del teorema \ref{thm:3.5}, entonces:
    %
    % K_1(f_1) = E_1 ---------> E_2 = K_2(f_2)
    %          |                          |
    %          |                          |
    %          |                          |
    %  \sigma: K_1 ---------------------> K_2
    %
    Sea $f_2 = \sigma^\star(f_1)$, entonces $\sigma$ se puede extender a $\tau: E_1 \to E_2$. Para cada $p_1 \divides f_1$ factores irreducibles de $f_1$ y para cada par $\alpha$ raíz de $p_1$ en $E_1$ y $\beta$ raíz $\sigma^\star(p_1) = p_2 \divides f_2$ existe $\tau: E_1 \to E_2$ tal que:
    $$
        \tau(\alpha) = \beta
    $$.
\end{obs}

%%%%%%%%%%%%%%%%%%%%%%%%%%%%%%%%%%%%%%%%%%%%%%%%%%%%%%%%%%%%%%%%%%%%%%%%%% 26/11
%TODO: Completar el ejercicio H4.5

Vamos a recordar algunos resultados de teoría de grupos.
\begin{pro}[Resultados de teoría de grupos]$ $
    \begin{enumerate}
        \item Si $G$ es cíclico $\implies G$ es abeliano.
        \item Si $\vabs{G} = p^a$ con $p$ un primo, entonces $Z(G) > 1$.
        \item Si $\vabs{G} = p^2$ con $p$ un primo, entonces $G$ es abeliano, en particular $G \isom C_{p^2}$ o $G \isom C_p \times C_p$.
    \end{enumerate}
\end{pro}

\begin{eg}[Cálculo del grupo de Galois de una extensión dado un polinomio (1)]
    Sea $x^4-2$ un polinomio en $\Q[x]$, calcula $\gal(E/\Q)$ siendo $E$ su cuerpo de escisión.\\\\

    \textbf{Solución}\\
    Las raíces en $\C$ del polinomio son $\sdf{\pm \sqrt[4]{2}, \pm \sqrt[4]{2} i}$, por tanto $E = \Q(\pm \sqrt[4]{2}, \pm \sqrt[4]{2} i)  = \Q(\sqrt[4]{2}, i)$.
    Ahora calculamos el grado de la extensión, por transitividad de grados: $\vabs{E:\Q(\sqrt[4]{2})} = 2$ y $\vabs{\Q(\sqrt[4]{2}):\Q} = 4$ por tanto $\vabs{G} = \vabs{E:\Q} = 8$, con $G = \gal(E/\Q)$.\\

    Además, sabemos que $\Q(\sqrt[4]{2})/\Q$ no es normal, lo que implica que $G$ tiene un subgrupo que no es normal (por el teorema \ref{thm:4.6}) y entonces $G$ no puede ser abeliano. De hecho $\gal(E/\Q(\sqrt[4]{2})) = \sdf{\1_E, \tau \mid \tau(i) = -i}$. Por tanto $G \isom D_8$ o $G \isom Q_8$.
    Sin embargo, en $Q_8$ todo subgrupo es normal, por lo que $G \isom D_8$.\\

    Recordamos que $D_8 = \gen{a, b\ \mid a^4 = b^2 = 1,\ a^b = b^{-1} a b = a^{-1}}$. En general:
    $$
        D_{2n} = \gen{\alpha, \beta\ \mid a^n = 1 = \beta^2,\ \alpha^\beta = \alpha^{-1}}
    $$
    Para calcular el grupo de Galois vamos a considerar $L = \Q(\sqrt 2, i)$, y ya hemos visto que $L = \gal(L/\Q) \isom C_2 \times C_2 = \sdf{\1_L, \alpha_2, \alpha_3, \alpha_4}$.\\

    \begin{center}
        \begin{tabular}{|c||c|c|c|}
        \hline
                   & $\sqrt{2}$  & $i$  & ord \\ \hline \hline
        $id$       & $\sqrt{2}$  & $i$  & $1$ \\ \hline
        $\alpha_2$ & $\sqrt{2}$  & $-i$ & $2$ \\ \hline
        $\alpha_3$ & $-\sqrt{2}$ & $i$  & $2$ \\ \hline
        $\alpha_4$ & $-\sqrt{2}$ & $-i$ & $2$ \\ \hline
        \end{tabular}
    \end{center}

    Entonces sea $E = L(x^2-\sqrt{2})$, sabemos que $\alpha_i: L \to L$ es un ismorfismo, y que su extensión $\hat{\alpha}: E \to E$, puede llevar raíces de $x^2 - \sqrt{2}$ (si $\alpha(\sqrt 2) = \sqrt 2)$ o raíces de $x^2 - \sqrt{2}$ en raíces de $x^2 + \sqrt{2}$ (si$ \alpha(\sqrt 2) = -\sqrt 2$).\\

    Además ya hemos comentado que $G \isom D_8$. Podemos hacer una tabla de sus elementos:\\
    \begin{center}
        \begin{tabular}{|c||c|c|c|}
        \hline
                 & $\sqrt[4]{2}$   & $i$  & ord \\ \hline \hline
        $\1_E$     & $\sqrt[4]{2}$   & $i$  & $1$ \\ \hline
        $\tau_2$ & $-\sqrt[4]{2}$  & $i$  & $2$ \\ \hline
        $\tau_3$ & $\sqrt[4]{2}$   & $-i$ & $2$ \\ \hline
        $\tau_4$ & $-\sqrt[4]{2}$  & $-i$ & $2$ \\ \hline
        $\tau_5$ & $\sqrt[4]{2}$   & $i$  & $4$ \\ \hline
        $\tau_6$ & $-\sqrt[4]{2}i$ & $i$  & $4$ \\ \hline
        $\tau_7$ & $\sqrt[4]{2}$   & $-i$ & $2$ \\ \hline
        $\tau_8$ & $-\sqrt[4]{2}i$ & $-i$ & $2$ \\ \hline
        \end{tabular}
    \end{center}

Queremos encontrar una presentación del grupo. Sabemos que $\tau_5^{-1}$ tiene orden $4$ distinto de $\tau_5$, por tanto $\tau_5^{-1} = \tau_6$. Por otra parte:
\begin{align*}
    &o(\tau_3) = 2 \iff \tau_3^2 = \1_E \iff \tau_3^{-1} = \tau_3\\
    &\tau_5^{\tau_3} = \tau_3^{-1}\tau_5\tau_3 = \tau_3\tau_5\tau_3 \implies \\
    &\tau_5^{\tau_3}(i) = i,\ \tau_5^{\tau_3}(\sqrt[4]{2}) = -\sqrt[4]{2}i \implies\\
    &\tau_5^{\tau_3} = \tau_6 = \tau_5^{-1}
\end{align*}
Por tanto:
$$
    G = \gen{\tau_5,\ \tau_3 \mid \tau_5^4 = \tau_3^2 = 1,\ \tau_5^{\tau_3} = \tau_5^{-1} }
$$
\end{eg}
\begin{ex}[H4.7]
    Sea $p$ un primo, y $f(x) = x^p - 1 \in \Q[x]$.
    \begin{itemize}
        \item[(a)] Halla $E = \Q(f)$
        \item[(b)] Prueba que $E/\Q$ es simple de grado $p - 1$
        \item[(c)] Demuestra que $G = \gal(E/\Q)$ es cíclico encontrando un isomorfismo explícito entre $G$ y $F_p^\times$
        \item[(d)] Demuestra que si $p$ es impar, entonces $E$ contiene exactamente una extensión cuadrática de $\Q$ (es decir, una extensión de grado $2$ sobre $\Q$).
    \end{itemize}

    \textbf{Solución}\\\\
    \begin{itemize}
        \item[(a)] Las raíces de $f$ son las raíces $p$-ésimas de la unidad, por tanto, sea $\Omega$ el conjunto de raíces:
        $$
            \Omega = \sdf{1, \xi, \ldots, \xi^{p-1}},\ \xi = e^{\frac{2\pi i}{p}} \implies E = \Q(\xi)
        $$
        \item[(b)] $\Q(\xi)/\Q$ es simple pues ya hemos visto que $\xi$ la genera. Además, como $f$ es ciclotómico:
        $$
            \vabs{E:\Q} = \delta(Irr(\Q, \xi)) = \delta(x^{p-1} + \cdots + x + 1) = p-1
        $$
        \item[(c)] $\sigma \in \gal(\Q(\xi)/\Q)$ queda determinada por $\sigma(\xi)$. Además sabemos que $\vabs{\gal(\Q(\xi)/\Q)} = p-1$. Por tanto:\\
        \begin{center}
            \begin{tabular}{|c|c|}
            \hline
            $G = \gal(\Q(\xi)/\Q) $ & $ \xi $\\ \hline \hline
            $\1 = \sigma_1 $ & $ \xi $\\ \hline
            $\sigma_2 $ & $ \xi^2 $\\ \hline
            $\vdots $ & $ \vdots $\\ \hline
            $\sigma_{p-1} $ & $ \xi^{p-1} $\\ \hline
            \end{tabular}
        \end{center}
        Tenemos que encontrar un isomorfismo entre $\F_p^\times$ y $G$. Recordemos que $F_p^\times = \sdf{1, 2, \ldots, p-1}$. Entonces, sea $\rho: G \to F_p^\times$ nuestro isomorfismo, basta establecer $\rho(\sigma_j) = j$, entonces:
        $$
            \rho(\sigma_k \sigma_j) = \rho(\sigma_{kj}) = \overline{kj} = \rho(\sigma_k) \rho(\sigma_j)
        $$
        Y es fácil ver que $\rho$ es biyectiva. Con esto hemos hallado que $G \isom \F_p^\times \isom C_{p-1}$.

        \item[(d)]  Por el TFTG (\ref{thm:fund-gal}), $G$ tiene un único subgrupo $H$ con $\vabs{G:H} = 2 \iff $ se corresponde con la única subextensión de $\Q(\xi) / \Q$ de grado $2$ sobre $\Q$. Es decir:
        $$
            \Q(\xi)^H = \Q(\xi)^{\sigma_j} \text{ (habría que calcular un elemento $\sigma_j$ tal que $o(\sigma_j) = \frac{p-1}{2}$)}
        $$
    \end{itemize}
\end{ex}



\begin{pro}[Teoremas de Sylow]
    Sea $G$ un grupo de orden $\vabs{G} = p^a n$ con $p$ primo y $p \not \divides n$, entonces:
    \begin{enumerate}
        \item $\exists P \leq G$ subgrupo, de orden $\vabs{P} = p^a$. Decimos que $P \in Syl_p(G)$.
        \item Sean $P_1, P_2 \in Syl_p(G) \implies \exists g \in G: P_1^g = P_2$, es decir, $\exists g \in G: \sdf{g^{-1} \cdot x \cdot g \mid x \in P_1} = P_2$.
        \item $\vabs{Syl_p(G)} = \vabs{G:N_G(P)}$.
    \end{enumerate}
\end{pro}

\begin{eg}[Cálculo del grupo de Galois de una extensión dado un polinomio (2)]
    Partimos del polinomio: $f = x^3 - 2 \in \Q[x]$, vamos a hallar el grupo de Galois de su cuerpo de escisión.
    \begin{enumerate}
        \item Calculamos $E = \Q(f)$.
        $$
            E = \Q(\sqrt[3]{2}, \sqrt[3]{2} \omega, \sqrt[3]{2} \omega^2) = \Q(\sqrt[3]{2}, \omega) = \sqrt(\sqrt[3]{2}, \sqrt{3} i) \text{ ya que } \omega = -\frac{1}{2} + \frac{\sqrt{3}i}{2}
        $$
        Además, como $\car(\Q) = 0$ es separable y $E = \Q(f)$ es normal, por tanto $E$ es de Galois.
        \item Hallamos el orden de la extensión.\\
        Podemos escribir el diagrama:\\
        \begin{center}
            \begin{tikzpicture}[node distance=1cm and 2cm, auto]
                \node (A) [above of =B] {$E = M(\omega)$};
                \node (B) {$M = \Q(\sqrt[3]{2})$};
                \node (C) [below of =B] {$\Q$};
                \draw[-] (A) to node {2 ($x^2+x+1$)}(B);
                \draw[-] (B) to node {3 ($x^3-2$)}(C);
            \end{tikzpicture}\\
        \end{center}
        Por tanto, por transitividad de grados $\vabs{E:\Q} = 6$ y una $\Q$-base de $E$ sería:
        $$
            \sdf{1, \sqrt[3]{2}, \sqrt[3]{4}, \omega, \sqrt[3]{2}\omega, \sqrt[3]{4}\omega}
        $$
        \item Determinar la clase de isomorfía de $G = \gal(E/\Q)$.\\
        Podemos verlo de dos formas:\\
        Sabemos que $\vabs{G} = \vabs{E:\Q} = 6$. Además por el teorema \ref{thm:3.11}, $G \leq S_3$ y por órdenes llegamos a que $G \isom S_3$.\\ % TODO: Explicar mejor (?)
        Por otra parte, sabemos que $M/\Q$ no es normal, por tanto por el TFTG (\ref{thm:fund-gal}), $\gal(E/M) \leq G$ no es normal, y entonces $G$ no es abeliano $\implies G \isom S_3$.

        \item Describir $G$ y sus órdenes.\\
        \begin{center}
            \begin{tikzpicture}[node distance=1cm and 2cm, auto]
                \node (A) [above of =B] {$E = L(\sqrt[3]{2})$};
                \node (B) {$L = \Q(\omega)$};
                \node (C) [below of =B] {$\Q$};
                \draw[-] (A) to node {$3$}(B);
                \draw[-] (B) to node {$2$}(C);
            \end{tikzpicture}\\
        \end{center}
        Como $L/\Q$ es normal (ya que es de grado $2$), todo $\tau \in G$ es la extensión de algún $\alpha \in \gal(L/\Q)$. Además, $\gal(L/\Q) = \sdf{1, \sigma}:\ \sigma(\omega) = \omega^2$.\\
        Cada elemento $\alpha \in \gal(L/\Q)$ se extiende a $E$ de tres formas distintas ($\tilde\alpha(\sqrt[3]{2}) \in \sdf{\sqrt[3]{2}, \sqrt[3]{2}\omega, \sqrt[3]{2}\omega^2}$) $\implies x^3-2 \in L[x]$ es irreducible y fijado por $\gal(L/\Q)$.
        Además, teniendo en cuenta que $\omega^2+\omega+1 = 0$ y $\omega^3 = 1$ podemos rellenar la siguiente tabla:
        \begin{center}
            \begin{tabular}{|c||c|c|c|}
            \hline
            Isomorfismo                 & $\omega$   & $\sqrt[3]{2}$         & ord\\ \hline \hline
            $\1_L   \rightarrow \tau_1 = 1_E$ & $\omega$   & $\sqrt[3]{2}$         & $1$\\ \hline
            $\1_L   \rightarrow \tau_2$       & $\omega$   & $\sqrt[3]{2}\omega$   & $3$\\ \hline
            $\1_L   \rightarrow \tau_3$       & $\omega$   & $\sqrt[3]{2}\omega^2$ & $3$\\ \hline
            $\sigma \rightarrow \tau_4$       & $\omega^2$ & $\sqrt[3]{2}$         & $2$\\ \hline
            $\sigma \rightarrow \tau_5$       & $\omega^2$ & $\sqrt[3]{2}\omega$   & $2$\\ \hline
            $\sigma \rightarrow \tau_6$       & $\omega^2$ & $\sqrt[3]{2}\omega^2$ & $2$\\ \hline
            \end{tabular}
        \end{center}
        Por los órdenes, y por ser un grupo de orden $6$ podemos determinar que es isomorfo a $S_3$ ya que no es cíclico y solo existen dos grupos de orden $6$. Otra forma sería intentar hallar una presentación del grupo a partir de relaciones entre los elementos, de esta forma hallaríamos la presentación:
        $$
            G = \gen{\tau_2 \tau_4 \mid o(\tau_2) = 3,\ o(\tau_4)=2,\ \tau_2^{\tau_4} = \tau_2^{-1}} = \gen{\tau_2} \ast \gen{\tau_4}
        $$
        \item Describir las subextensiones de $E/\Q$ indicando cuáles de ellas son normales.\\
        Por el TFTG:\\
        \begin{center}
            \begin{tikzpicture}[node distance=1cm and 2cm, auto]
                \node (B) {$K$};
                \node (A) [above of =B] {$\sdf{\Q \subseteq K \subseteq E}$};
                \node (E) [right= of B] {$\gal(E/K)$};
                \node (C) [below of =B] {$E^H$};
                \node (D) [above of =E] {$\sdf{H \leq G}$};
                \node (F) [below of =E] {$H$};
                \draw[<->] (A) to node {$1:1$}(D);
                \draw[|->] (B) to node {}(E);
                \draw[<-|] (F) to node {}(C);
            \end{tikzpicture}\\
        \end{center}
        $$
            K/\Q \text{ es normal } \iff \gal(E/K) \normsub G
        $$
        \begin{center}
            \begin{tikzpicture}[node distance=1cm and 2cm, auto]
                \node (A) {$E$};
                \node (B) [right= of A] {$\1$};
                \node (C) [below of =A] {$E^H$};
                \node (D) [below of =B] {$H$};
                \node (E) [below of =C] {$\Q$};
                \node (F) [below of =D] {$G$};
                \draw[-] (A) to node {}(C);
                \draw[-] (C) to node {$n$}(E);
                \draw[-] (B) to node {}(D);
                \draw[-] (D) to node {$n$}(F);
                \draw[<->] (A) to node {$1:1$}(B);
                \draw[<->] (E) to node {$1:1$}(F);
            \end{tikzpicture}\\
        \end{center}
        Como $\vabs{G} = 6$, si $1 < H < G$ ($H$ subgrupo propio), entonces por el teorema de Lagrange para grupos, si:
        \begin{align*}
            \vabs{H} = 2 \implies H \in Syl_2(G)\\
            \vabs{H} = 3 \implies H \in Syl_3(G)\\
        \end{align*}
        Entonces si $\vabs{H} = 3 \implies \vabs{G:H} = 2 \implies H \normsub G$. Con esto hemos visto que las subextensiones de $E/\Q$ son normales si son triviales o se corresponden con un grupo en $Syl_3(G)$.
        \begin{center}
            \begin{tikzpicture}[node distance=1cm and 2cm, auto]
                \node (A) {$P = \gen{\tau_2}$};
                \node (B) [right= of A] {$ $};
                \node (C) [right= of B] {$\gen{\tau_4} = Q$};
                \node (D) [above of =B] {$G$};
                \node (E) [below of =B] {$1$};
                \draw[-, double] (D) to node {}(A);
                \draw[-] (D) to node {$N_G(Q)$}(C);
                \draw[-] (A) to node {$3$}(E);
                \draw[-] (C) to node {$2$}(E);
            \end{tikzpicture}\\
        \end{center}
        Sabemos que $P \in Syl_3(G)$ y $Q \in Syl_2(G)$. Como $Q \cancel \normsub G$, consideremos el normalizador $N_G(Q) < G$, entonces $\vabs{N_G(Q):Q}\vabs{G:Q} = 3 \implies N_G(Q) = Q$. Por tanto $\vabs{Syl_2(G)} = 3$:
        $$
            Syl_2(G) = \sdf{Q,\ Q^{\tau_2},\ Q^{\tau_2^{-1}}} \text{ y } Q \neq Q^{\tau_2} \neq Q^{\tau_2^{-1}} \neq Q
        $$
        Y entonces $1, G, P=\gen{\tau_2}, Q, Q^{\tau_2}, Q^{\tau_3}$ son todos los subgrupos de $G$ donde $Q = \gen{\tau_4}$. Además, $E/\Q$ tiene una subextensión normal que es de grado $2$ sobre $\Q$. Sea $L$ esta extensión afirmamos que $L/\Q$ es normal y $\vabs{L:\Q} = 2$, $L=\Q(\omega)$ y por el TFTG: $L = E^{P} = E^{\gen{\tau_2}} = E^{\tau_2}$.

        Para calcular las subextensiones de grado $3$ sobre $\Q$, vamos a calcular $M_1 = E^{Q}$, $M_2 = E^{Q^{\tau_2}}$, $M_3 = E^{Q^{\tau_2^{-1}}}$. Además por el lema \ref{lm:4.5}:
        $$
            M_2 = E^{Q^{\tau_2}} = t_2(M_1),\ M_3 = E^{Q^{\tau_2^{-1}}} = \tau_2^{-1}(M_1)
        $$
        Vamos a calcular $M_1$, recordamos que $\sdf{1, \alpha, \alpha^2, \omega, \alpha \omega, \alpha^2 \omega}$ es una $\Q$-base de $E$ donde $\alpha = \sqrt[3]{2}$:
        \begin{align*}
            M_1 &= E^Q = E^{\gen{\tau_4}} = E^{\tau_4}\\
            E^{\tau^4} &= \sdf{x \in E:\ \tau_4(x) = x}\\
            x &= a + b\alpha + c \alpha^2 + d\omega + e\alpha\omega + f\alpha^2\omega = \tau_4(x) =\\
              &= \cancel a + \cancel{b\alpha} + \cancel{c\alpha^2} + d\omega^2 + e\alpha\omega^2 + f\alpha^2\omega^2 = (\text{ ya que } \omega^2+\omega+1 = 0)\\
              &= d(-\omega -1) + e\alpha^2(-\omega - 1) + f\alpha^2(-\omega -1) =\\
              &= d + 2d\omega + e\alpha + 2e\alpha \omega + f\alpha^2 + 2f\alpha^2\omega = 0 \implies\\
              &\implies d = e = f = 0
        \end{align*}
        Con este desarrollo hallamos:
        $$
            M_1 = E^{\tau_4} = \sdf{a + b\alpha + c\alpha^2 \mid a,b,c \in \Q} = \Q(\alpha) = \Q(\sqrt[3]{2})
        $$
        Para hallar $M_2$ y $M_3$ podemos repetir el mismo procedimiento o podemos:
        $$
            M_2 = E^{Q^{\tau_2}} = \tau_2(E^Q) = \tau_2(M_1) = Q(\sqrt[3]{2} \omega)
        $$
        y recordemos que $\tau_2^{-1} = \tau_3$, por tanto:
        $$
            M_3 = E^{Q^{\tau_3}} = \tau_3(E^Q) = \tau_3(M_1) = Q(\sqrt[3]{2} \omega^2)
        $$
    \end{enumerate}
\end{eg}

\begin{pro}[Consecuencia del teorema \ref{thm:3.11}]
    Sea $E/K$ una extensión, $E = K(f)$, $f \in K[x]$ irreducible y separable, y $\delta(f) = n$. Entonces $G = \gal(E/K) \leq S_n$ es transitivo, es decir:
    $$
        \forall i,j \in \sdf{1, \ldots, n} \exists \delta \in G \mid \delta(i) = j
    $$
\end{pro}
\begin{proof}
    Sea $\Omega = \sdf{\alpha_1, \ldots \alpha_n}$ el conjunto de raíces distintas de $f$ en $E$. Sea el morfismo $\rho: G \to S_\Omega;\ \delta \mapsto \rho(\delta)(\alpha) = \delta(\alpha)$, entonces como $f$ es separable $\ker(\rho) = \gal(E/E) = 1$ (tomando $L=E$ en \ref{thm:3.11}).\\

    Por tanto, $G \leq S_\Omega \isom S_n$, y como $f$ es irreducible, dados $\alpha_i, \alpha_j \in \Omega$ podemos afirmar que $\exists \delta \in G \mid \delta(\alpha_i) = \alpha_j$.
\end{proof}

\begin{obs}
    Los subgrupos transitivos de $S_n$ son $D_{2n}$
\end{obs}

\begin{pro}[Consecuencia del Teorema del Elemento Primitivo (\ref{thm:elemprim})]
    Sea $E/K$ una extensión de Galois, existe $f \in K[x]$ irreducible tal que $E = K(f)$.
\end{pro}

\begin{proof}
    Como $E/K$ es normal por ser de Galois, entonces $E = K(f)$, para algún $f \in K[x]$. Si $f$ es irreducible, lo hemos demostrado. Si $f$ no es irreducible, entonces $E = K_1(\alpha_1, \ldots, \alpha_n)$ y sabemos que por el teorema \ref{thm:elemprim}, $\exists \alpha \in E: E = K(\alpha)$. Entonces basta tomar $p = Irr(K, \alpha)$, con lo que tendríamos $E = K(p)$ con $p$ irreducible.
\end{proof}

\begin{pro}[Consecuencia del lema \ref{lm:4.5} y del teorema \ref{thm:4.3}]
    Sea $E/K$ una extensión de Galois, $G = \gal(E/K)$ su grupo de Galois, $M$ un cuerpo intermedio y $H = \gal(E/M) \leq G$, entonces:
    $$
        \forall \tau \in G,\ \tau(M) = E^{H^\tau}
    $$
\end{pro}
\begin{proof}
    Por el lema \ref{lm:4.5}, $H^\tau = \gal(E/\tau{M})$. Además, $E/K$ de Galois $\implies E/\tau(M)$ de Galois $\implies$ (por el teorema \ref{thm:4.3}) $E^{H^{\tau}} = \tau(M)$
\end{proof}

\begin{obs}
    El teorema \ref{thm:4.3} nos dice que sean $E/K$ finito y $G = \gal(E/K)$, si $E/K$ de Galois, entonces $E^G=K$. Se puede ver que el recíproco también es cierto. En resumen:\\

    \begin{center}
        Sean $E/K$ finita y $G=\gal(E/K)$:
        $$
            E^G = K \iff E/K \text{ es de Galois}
        $$
    \end{center}
\end{obs}

\begin{ex}[H4.11]
    Sea $E/K$ una extensión de Galois con $G = \gal(E/K)$ cíclico de orden $n$. Demuestra que:
    \begin{itemize}
        \item[(a)] Para cada divisor $d$ de $n$ existe exactamente un cuerpo intermedio $L$ con $\vabs{E:L} = d$.
        \item[(b)] Si $L_1$ y $L_2$ son dos cuerpos intermedios, entonces $L_1 \subseteq L_2$ si, y solo si, $E:L_2$ divide a $E:L_1$.
        \item[(c)] Demuestra que la hipótesis de $E/K$ de Galois puede relajarse a que $E/K$ sea finita.
    \end{itemize}

    \textbf{Solución}
    \begin{itemize}
        \item[(a)] Basta representar las correspondencias entre subcuerpos intermedios y subgrupos de Galois para comprobarlo.

        \begin{center}
            \begin{tikzpicture}[node distance=1cm and 2cm, auto]
                \node (A) {$E$};
                \node (B) [right= of A] {$\1$};
                \node (C) [below of =A] {$L = E^H$};
                \node (D) [below of =B] {$H$};
                \node (E) [below of =C] {$K$};
                \node (F) [below of =D] {$G$};
                \draw[-] (A) to node {$d$}(C);
                \draw[-] (C) to node {}(E);
                \draw[-] (B) to node {}(D);
                \draw[-] (D) to node {}(F);
                \draw[<->] (C) to node {$1:1$}(D);
            \end{tikzpicture}\\
        \end{center}
        \item[(b)]
        \begin{itemize}
            \item[$(\implies)$] Es trivial.
            \item[$(\Longleftarrow)$]
                \begin{align*}
                    d = &\vabs{E:L_2} = \vabs{H_2},\ H_2 = \gal(E/L_2)\\
                        &\vabs{E:L_1} = \vabs{H_1},\ H_1 = \gal(E/L_1) \implies d \divides \vabs{alpha_1} \divides n
                \end{align*}
                $H_1$ tiene un único subgrupo de orden $d$. Como $H_2$ es el único subgrupo de orden $d$ de $G \implies $ tienen que coincidir $\implies H_2 \subseteq H_1 \implies L_1 = E^{H_1} \subseteq L_2 = E^{H_2}$
        \end{itemize}
        \item[(c)]
            Basta aplicar que recíproco del teorema \ref{thm:4.3} (se resuelve en el ejercicio H4.4)
    \end{itemize}
\end{ex}

\begin{ex}[H4.8]
    Sea $E = \Q(\xi)$ donde $\xi = e^{\frac{2 \pi i}{7}}$. Muestra que $E$ es una extensión de Galois de $\Q$. Encuentra todos los subcuerpos intermedios de la extensión $E/\Q$, los subgrupos de $\gal(E/\Q)$ que les corresponden y determina qué subcuerpos intermedios se corresponden con extensiones normales de $\Q$.\\\\

    \textbf{Solución}\\
    $\xi$ es una raíz $p$-ésima de la unidad, con $p=7$. Es decir, $E = \Q(\xi) = \Q(x^7-1) = \Q(x^6 + x^5 + x^4 + x^3 + x^2 + x + 1)$. Por tanto, $\vabs{\Q(\xi):\Q} = 6$ y $G = \gal(\Q(\xi)/\Q) \isom \F_7^\times$ que es cíclico.\\
    $$G = \sdf{\varphi_j \mid 1 \leq j \leq 6} \mapsto \bar{j},\ \varphi(\xi) = \xi^j$$

    Buscamos $\gen{\varphi} = G$, que es lo mismo que buscar un generador de $\F_7^\times$ (de orden $6$). Fácilmente podemos ver que $3$ es un generador de $\F_7^\times$. Tomamos $\varphi = \varphi_3$, y como $G = C_6 \implies \varphi_2 = \varphi^2,\ \varphi_6 = \varphi^3$.\\
    Los subgrupos son:
    $$
        1,\ \gen{\varphi^2}, \gen{\varphi^3}, G
    $$
    Que se corresponden con las subextensiones (por el TFTG):
    $$
        E,\ L=E^{\varphi^2}, M=E^{\varphi^3}, \Q
    $$
    Vamos a calcular $L$:
    \begin{align*}
        L=E^{\varphi^2} &= \sdf{x = a_0 + a_1 \xi + \cdots a_5 \xi^5,\ a_i \in \Q \mid \varphi_2(x) = x}\\
        \varphi_2(x)    &= a_0 + a_1\xi^2 + a_2\xi^4 + a_3\xi^6 + a_4\xi + a_5\xi^3 =\\
                        &= a_0 + a_1\xi^2 + a_2\xi^4 + a_3(-1 -\xi \cdots -\xi^5) + a_5\xi^3 = x \iff \\
                        &\iff \xi(a_1-a_4-a_4) + \xi^2(a_2-a_1+a_3) + \xi^3(a_3-a_5+a_3) + \xi^4(a_4-a_2+a_3) + \xi^5(a_5-a_3) = 0
    \end{align*}
    %%TODO: REHACER FOTOS SANTORUM
    Con estas restricciones hallamos $L =\Q(\sqrt{-1})$.
\end{ex}

% % !TeX root = ../galois.tex

\chapter{Aplicaciones}

\section{Introducción. Radicales y resolubilidad.}

En este capítulo vamos a ver el \textit{Gran Torema de Galois}, que nos dice que un polinomio es \textit{resoluble por radicales} si y sólo si su grupo de Galois es \textit{resoluble}. Como es evidente, vamos a tener que introducir una serie de conceptos nuevos para desarrollar esta parte de la teoría.

\begin{dfn}[Radical]
    $R/K$ es \textbf{radical} si existe una cadena de subcuerpos $K = K_0 \subseteq K_1 \subseteq \cdots \subseteq K_{n-1} \subseteq K_n = R$, de modo que $K_j = K_{j-1}(a_j),\ (a_j)^{n_j} \in K_{j-1}$.

    \begin{center}
        \begin{tikzpicture}[node distance=1cm and 2cm, auto]
            \node (R)  {$R = K_n$};
            \node (Kj) [below of =R] {$K_j = K_{j-1}(a_j)$};
            \node (K2) [below of =Kj] {$K_2 = K_1(a_2)$};
            \node (K1) [below of =K2] {$a_2^{n_2} \in K_1 = K_0(a_1),\ x^{n_2}-a_2^{n_2}$};
            \node (K0) [below of =K1] {$a_1^{n_1} \in K = K_0 $};
            \draw[-] (R) to node {}(Kj);
            \draw[-] (Kj) to node {}(K2);
            \draw[-] (K2) to node {}(K1);
            \draw[-] (K1) to node {}(K0);
        \end{tikzpicture}\\
    \end{center}
\end{dfn}

\begin{eg}[Ejemplos de radicales]
    \begin{enumerate}[a)]
        \item Ejemplo 1:\\$$\Q(e^{\frac{2\pi i}{n}})/\Q,\ \ \ \ (e^{\frac{2\pi i}{n}})^n = 1 \in \Q$$
        \item Ejemplo 2:\\

        \begin{center}
            \begin{tikzpicture}[node distance=1cm and 2cm, auto]
                \node (A) {$ K_2 = \Q(\sqrt[3]{2}, i) = K_1(i) $};
                \node (B) [below of =A]{$ K_1 = \Q(\sqrt[3]{2}) = K_0(i)$};
                \node (C) [below of =B] {$ K_0 = \Q $};
                \draw[-] (A) to node {$i^2 = -1 \in K_1$}(B);
                \draw[-] (B) to node {$(\sqrt[3]{2})^3 = 2 \in K_0$}(C);
            \end{tikzpicture}\\
        \end{center}
        \item Ejemplo 3:\\

        \begin{center}
            \begin{tikzpicture}[node distance=1cm and 2cm, auto]
                \node (A) {$ \Q(\sqrt[3]{1+\sqrt 7})$};
                \node (B) [below of =A]{$ \Q(\sqrt 7)$};
                \node (C) [below of =B] {$ \Q $};
                \draw[-] (A) to node {$(\sqrt[3]{1+\sqrt 7})^3 = 1 + \sqrt{7} \in \Q(\sqrt 7)$}(B);
                \draw[-] (B) to node {$(\sqrt 7)^2 = 7 \in \Q $}(C);
            \end{tikzpicture}\\
        \end{center}
    \end{enumerate}
\end{eg}

\begin{dfn}[Resolubilidad por radicales]
    $f \in K[x]$ es \textbf{resoluble por radicales} si existe alguna extensión radical $R/K$ tal que $K(f) \subseteq R$.
\end{dfn}

\begin{eg}[Ejemplos de resolubilidad por radicales]
    \begin{enumerate}[a)]
        \item $x^p-1$. $\Q(x^p-1) = \Q(e^{\frac{2\pi i}{p}})$.
        \item $x^4-2$. $\Q(\sqrt[4]{2}, i)$.
        \item $x^4+4x^2+2$. $\Q(\sqrt{2+\sqrt{2}})$.
    \end{enumerate}
\end{eg}

\section{Grupos resolubles}
Como se anticipó en la sección anterior, el \textit{Gran Teorema de Galois} utiliza el concepto de grupo resoluble, para ello vamos a definir los conceptos necesarios.

\begin{dfn}[Grupo simple]
    Sea $G$ un grupo, es \textbf{simple} si no tiene subgrupos normales.
\end{dfn}

\begin{obs}$ $
    \begin{itemize}
        \item $G$ abeliano es simple si $G \isom C_p$.
        \item Sea $G_m$ el subgrupo máximo normal de $G$. Entonces $G_m \normsub G$ y no hay normales entre $G_m$ y $G \implies \quo{G}{G_m}$ es simple.
    \end{itemize}
\end{obs}

\begin{obs}[Clasificación de grupos simples]
    Si $G$ es finito y simple, entonces:
    \begin{itemize}
        \item $G \isom C_p$, si son abelianos.
        \item $G \isom A_n: n\geq5$.
        \item $G$ es de tipo \textit{Lie}.
        \item 26 esporádicos.
    \end{itemize}
\end{obs}

Los grupos finitos simples tienen relevancia en teoría de grupos ya que informalmente son los \textit{bloques} que forman el resto de grupos finitos. En cierta forma un grupo simple es a un grupo finito lo que un número primo a un número natural. De la misma forma que se puede descomponer en factores primos a un número natural, podemos escribir una descomposición de un grupo en subgrupos simples.

\begin{dfn}[Serie de descomposición]
    Sea $G$ un grupo finito, y $\sdf{H_j}_{j=1}^{j=n}$ una sucesión de subgrupos tales que:
    $$
        \1 = H_0 \normsub H_1 \normsub \cdots \normsub H_n = G
    $$
    se llama a dicha cadena \textbf{serie de descomposición} del grupo si $\forall j,\ \quo{H_{j+1}}{H_j}$ es un grupo simple y $n$ es la longitud de la serie.
\end{dfn}

\begin{obs}
    Aunque queda fuera del temario, todas las series de descomposición de un mismo grupo $G$ son equivalentes salvo por permutaciones e isomorfismos.
\end{obs}

\begin{dfn}[Grupo resoluble]
    Un grupo $G$ es \textbf{resoluble} si existe una serie de descomposición de la forma:
    $$
        \sdf{e} = G_0 \normsubeq G_1 \normsubeq G_2 \normsubeq \cdots \normsubeq G_{n-1} \normsubeq G_n = G
    $$
    tal que $\quo{G_j}{G_{j-1}}$ es abeliano.
\end{dfn}

\begin{pro}[Resolubilidad de grupos abelianos]
    Todo grupo $G$ abeliano es resoluble, es decir:
    $$
        G \text{ abeliano} \implies G \text{ resoluble.}
    $$
\end{pro}
\begin{proof}
    Como $\sdf{1} \subseteq G$, y $\sdf{1} \normsub G$ dado que $x \cdot 1 \cdot x^{-1} \in \sdf{1}$ y además $\quo{G}{\sdf{1}} \isom G$ es abeliano, hemos encontrado una serie:
    $$
        \sdf{1 = e} \normsub G
    $$
    tal que $\quo{G_j}{G_{j-1}}$ es abeliano.
\end{proof}
\begin{eg}$ $
    \begin{enumerate}[i)]
        \item $A_4$ es resoluble. Sea $V = \gen{(12)(34), (13)(24)}$:\\

        \begin{center}
            \begin{tikzpicture}[node distance=1cm and 2cm, auto]
                \node (A) {$ A_4 $};
                \node (B) [below of =A]{$ V $};
                \node (C) [below of =B] {$ \sdf{1} $};
                \draw[-, double] (A) to node {$3$}(B);
                \draw[-] (B) to node {$4$}(C);
            \end{tikzpicture}\\
        \end{center}
        \item $D_{2n}$ es resoluble.
        \item $S_3$ y $S_4$ son resolubles.
    \end{enumerate}
\end{eg}

\begin{obs}[Recuerdo de teoría de grupos]
    Sea $G$ un grupo, $N\normsubeq G$ y $H \leq G$ subgrupos y sea $M = H \cap N$, entonces:
    $$
        M \normsub H \text{ y } \quo{H}{M} \isom \quo{HN}{N}
    $$
    \begin{center}
        \begin{tikzpicture}[node distance=1cm and 2cm, auto]
            \node (A) {$N$};
            \node (B) [right= of A] {$ $};
            \node (C) [right= of B] {$H$};
            \node (D) [above of =B] {$G$};
            \node (E) [below of =B] {$N\cap H = M$};
            \node (F) [below of =E] {$\1$};
            \draw[-, double] (D) to node {}(A);
            \draw[-] (D) to node {}(C);
            \draw[-] (A) to node {}(E);
            \draw[-] (C) to node {}(E);
            \draw[-] (E) to node {}(F);
        \end{tikzpicture}\\
    \end{center}
\end{obs}

\section{Resolubilidad y extensiones de Galois}

En esta sección vamos a ver resultados que relacionen los conceptos de un radical $R/K$ con extensiones de Galois, ya que en general un radical no es una extensión de Galois, como por ejemplo $\Q(\sqrt[4]{2})/\Q$ y
$\Q(\sqrt{1+\sqrt 7})/\Q$.

\begin{thm}[Kummer]\label{thm:5.1}
    Sean $E, K$ cuerpos de la forma $\Q \subseteq K \subseteq E \subseteq \C$, si $K$ contiene $\xi \in \C$ con $o(\xi) = n$, entonces son equivalentes:
    \begin{enumerate}[a)]
        \item $E = K(a)$ con $a \in E$ tal que $a^n \in K$.
        \item $E/K$ es una extensión de Galois con $\gal(E/K)$ cíclico de orden un divisor de $n$.
    \end{enumerate}
\end{thm}

\begin{proof}$ $
    \begin{enumerate}
        \item[a) $\implies$ b)] Tenemos que $E = K(a)$ con $a^n \in K$. Podemos suponer que $a \notin \K$. Sea $f(x) = x^n - a^n \in K[x]$, las raíces de $f$ son $\sdf{a\xi^j \mid 0 \leq j \leq n}$.
        Como $E = K(a) = K(a\xi^j\mid 0 \leq j \leq n) = K(x^n - a^n)$, entonces $E/K$ es de Galois.\\
        Sea $\sigma \in \gal(E/K)$, $\sigma(a) = a\xi^{\kappa_\sigma}$ donde $0 \leq \kappa_\sigma \leq n$. Definimos:\\
        \begin{center}
            \begin{tikzpicture}[node distance=1cm and 2cm, auto]
                \node (B)                {$\sigma$};
                \node (A)  [above of =B] {$\rho: \gal(E/K)$};
                \node (E)  [right= of B] {$\xi^{\kappa_\sigma}$};
                \node (D)  [above of =E] {$\gen{\xi}$};
                \draw[->]  (A) to node {}(D);
                \draw[|->] (B) to node {}(E);
            \end{tikzpicture}\\
        \end{center}
        Y entonces:
        \begin{itemize}
            \item $\rho(\1_E) = \rho^0 = 1$
            \item $\rho(\sigma\tau) = \xi^{\kappa_\sigma+\kappa_\tau} \xi^{\kappa_\sigma} * \xi^{\kappa_\tau} = \rho(\sigma)\rho(\tau)$.
            Donde $\sigma\tau(a) = \tau(a \xi^{\kappa_\sigma}) = a \xi^{\kappa_\tau} \xi^{\kappa_\sigma} = \xi^{\kappa_\sigma+\kappa_\tau}$
            \item $\rho$ es inyectiva porque $\rho(\sigma) = 1 \implies \sigma(a) = a \implies \sigma = 1_E$ y por tanto $\gal(E/K) \leq \gen{\xi}$.
        \end{itemize}
        \item[b) $\implies$ a)] Como $E/K$ es de Galois, $\gal(E/K) = \gen{\tau}$ y sea $n$ el orden del grupo de Galois, entonces: $o(\tau) = d\divides n$ y $\vabs{E:K} = d$.

        Sea $\omega \in\gen{\xi}$ tal que $o(\omega) = d$ (por ejemplo $\omega = \xi^{\frac{n}{d}})$, basta encontrar un $a \in E$ tal que $\tau(a) = a\omega$ y $a \notin K$. Vamos a ver por qué:\\
        Supongamos que existe tal $a$, entonces:
        $$
            \tau(a^d) = (\tau(a))^d = a^d \omega^d = a^d \implies a^d \in E^\gen{\tau} = K.
        $$
        En particular, $a^n = (a^d)^{\frac{n}{d}} \in K$. Sea $p = Irr(K, a) \in K[x]$, $p \divides x^d - a^d$.
        Por otro lado, $a, a\omega = \tau(a), \ldots, a\omega^{d-1} = \tau^{d-1}(a)$ son raíces de $p$ distintas $\implies \delta(p) \geq d \implies p(x) = x^d - a^d$, por tanto:
        $$
            \vabs{K(a):K} = \delta(p) = d = \vabs{E:K} \implies E = K(a)
        $$

        Finalmente, solo nos queda buscar el $a$ que mencionamos arriba. Por el \hyperref[thm:4.2]{teorema de Dedekind}, $\sdf{1, \tau, \ldots, \tau^{d-1}}$ es linealmente independiente y entonces:
        $$
            \1 + \frac{1}{\omega} \tau + \cdots + \frac{1}{\omega^{d-1}}\tau^{d-1} \neq 0
        $$
        Sea $b \in E$ tal que $0 \neq b + \frac{1}{\omega} \tau(b) + \cdots + \frac{1}{\omega^{d-1}}\tau^{d-1}(b) = a$ y hemos encontrado el $a$ que buscábamos. Puede comprobarse que $\tau(a) = \omega a = a \omega$.
    \end{enumerate}
\end{proof}

El siguiente grafo resume el resultado del teorema:\\
\begin{center}
    \begin{tikzpicture}[edge quotes mid, node distance=1cm and 2cm, auto]
        \node (A) [above of =C]{$E(\xi) = \gen{K(\xi), E}$};
        \node (B) [] {$K(\xi)$};
        \node (C) [right= of B] {$ $};
        \node (D) [right= of C] {$E = K(a)$};
        \node (E) [below of =C] {$K(\xi)\cap E$};
        \node (F) [below of =E] {$K$};
        \draw[-] (A) to node [sloped]{abeliana}(B);
        \draw[-] (A) to node {}(D);
        \draw[-] (B) to node {}(E);
        \draw[-] (D) to node {}(E);
        \draw[-] (B) to node [sloped, below]{abeliana}(F);
    \end{tikzpicture}\\
\end{center}

\begin{ex}[H5.2]
    Sea $E/K$ una extensión de Galois, $K \subseteq L, M \subseteq E$ subcuerpos intermedios. $\gen{L, M} = \bigcap_{H \subseteq E} H$ tal que   $L,M \subseteq H$. Prueba que:
    $$
        \gal(E/L) \cap \gal(E/M) = \gal(E/\gen{L, M})
    $$
    \textbf{Solución}\\\\

    Vamos a resolverlo por doble inclusión. Por definición $\gal(E/\gen{L, M}) \subseteq \gal(E/L) \cap \gal(E/M)$.\\
    Además $L,M\subseteq E^{\gal(E/L) \cap \gal(E/M)} \subseteq E$ y $\gen{L, M} \subseteq E^{\gal(E/L) \cap \gal(E/M)}$. Por tanto:
    $$
        \gal(E/E^{\gal(E/L) \cap \gal(E/M)}) \subseteq \gal(E/\gen{L, M})
    $$
    y $\gal(E/L) \cap \gal(E/M) = \gal(E/E^{\gal(E/L) \cap \gal(E/M)})$ por el \hyperref[thm:fund-gal]{TFTG}.
\end{ex}

\begin{ex}[H5.4]
    Por completar. %TODO
\end{ex}

\section{Gran Teorema de Galois}

\textit{NOTA: Solo están redactados los resultados teóricos de esta sección. Todos los ejercicios de esta sección están por redactar (se corresponden con H5.8, H5.10, H5.11, H5.12, H5.5, H5.6). También falta algún ejemplo, y las pruebas de los resultados.}

\begin{obs}[Recuerdo de teoría de grupos]
    Sea $G$ un grupo y $N \normsubeq G$, $H \leq G$ subgrupos, entonces:
    \begin{enumerate}[a)]
        \item $NH \leq G$, $N \cap H \normsubeq H$ y además $$\quo{NH}{N} \isom \quo{H}{N\cap H}$$.
        \item Si además $H \normsub G$, entonces $NH \normsub G$, $H  \normsubeq NH$ y además:
        $$
            \quo{G}{NH} \isom \quo{\quo{G}{N}}{\quo{NH}{N}}
        $$
    \end{enumerate}
\end{obs}

\begin{dfn}[Extensión abeliana]
    Una extensión $E/K$ es \textbf{abeliana} si $E/K$ es de Galois y $\gal(E/K)$ es un grupo abeliano.
\end{dfn}

\begin{lm}[Extensión de Galois y resoluble]\label{lm:5.2}
    Sea $\Q \subseteq K$, $L \subseteq \C$ y $L/K$ una extensión. Supongamos que existe una sucesión $\sdf{K_j}_{j=0}^{j=m}$ tal que $K = K_0$, $L = K_m$ monótona creciente por inclusión y $K_j/K_{j-1}$ es abeliana. Entonces si $E/K$ es de Galois, con $E \subseteq L$, entonces $\gal(E/K)$ es resoluble.
\end{lm}
\begin{proof}
    A completar.
\end{proof}
\begin{thm}[Gran Teorema de Galois]\label{thm:5.3}\label{thm:GTG}
    Sea $\Q \subseteq K$ un cuerpo, $f \in K[x]$ un polinomio y $E = K(f) \subseteq \C$ su cuerpo de escisión.
    $$
        f \text{ es resoluble por radicales } \iff \gal(f) = \gal(E/K) \text{ es resoluble }
    $$
\end{thm}
\begin{proof}
    A completar.
\end{proof}

\begin{lm}[Generación del grupo simétrico]\label{lm:5.4}
    Para $n \leq 2$, $S_n = \gen{(12)(12\cdots n)}$
\end{lm}

\begin{thm}[de Abel-Ruffini]
    Las ecuaciones polinómicas \textbf{generales} $f(x) = a_n x^n + \cdots + a_1 x + a_0 = 0$ no son resolubles por radicales cuando $\delta(f)\geq 5$. Es decir, no es posible encontrar las soluciones de $f$ de forma general aplicando finitamente la suma, resta, multiplicación y radicación.
\end{thm}

\begin{proof}
    A completar.
\end{proof}

\begin{obs}
    Equivalentemente, el teorema enuncia que no existe una fórmula general para hallar las raíces de $f(x)$ con $\delta(f) \geq 5$ aplicada a sus coeficientes. Para $\delta(f) \leq 4$ si es posible, por ejemplo con $\delta(f) = 2$ tenemos:
    $$
        x = \frac{-b \pm \sqrt{b^2-4ac}}{2a}
    $$
    Existen también fórmulas para polinomios de tercer y cuarto grado.
\end{obs}

\part{Apéndices}

%\renewcommand{\chaptername}{Apéndice}
%\renewcommand{\thechapter}{\Alph{chapter}}
%\setcounter{chapter}{0}

%% !TeX root = ../apuntes-ea.tex

\chapter{Ejercicios}

\section{Hoja 1}

\begin{ex}[H1.2]
	Sean $a,b,c \in G = (-1,1)$. Probamos las propiedades de los grupos.
	\begin{itemize}
		\item \textbf{Asociatividad:}
		\begin{align*}
			(a \ast b) \ast c = \left(\frac{a+b}{1+ab}\right) \ast c = \frac{\left(\frac{a+b}{1+ab}\right) + c}{1 + \left(\frac{a+b}{1+ab}\right)c} = \frac{\frac{a+b+c+abc}{1+ab}}{\frac{1+ab+ac+bc}{1+ab}} = \frac{a+b+c+abc}{1+ab+ac+bc} \\
			a \ast (b \ast c) = a \ast \left(\frac{b+c}{1+bc}\right) = \frac{a+\left(\frac{b+c}{1+bc}\right)}{1+a\left(\frac{b+c}{1+bc}\right)} = \frac{\frac{a + b + c + abc}{1+bc}}{\frac{1+ab+ac+bc}{1+bc}} = \frac{a+b+c+abc}{1+ab+ac+bc}
		\end{align*}
		\item \textbf{Elemento neutro:} es el $0$ ya que $x\ast 0 = \frac{x+0}{1+x\cdot 0} = \frac{x}{1} = x$ y además $0 \ast x = \frac{0+x}{1+0\cdot x} = \frac{x}{1} = x$
		\item \textbf{Elemento inverso:} la ecuación
		\begin{align*}
			x \ast \inv{x} = 0 \iff \frac{x + \inv{x}}{1+x\inv{x}} = 0 \iff \inv{x} = -x
		\end{align*}
		siempre tiene solución y ocurre lo mismo para la ecuación $\inv{x} \ast x = 0 \iff \inv{x} = -x$
		\item \textbf{Clausura:} tenemos que probar que si $x,y \in (-1,1)$ entonces $x \ast y \in (-1,1)$. Consideramos $f(x,y) = x \ast y = \frac{x+y}{1+xy}$. Derivando tenemos que $\nabla f(x,y) = (\frac{1}{(1+xy)^2},\ \frac{1}{(1+xy)^2}) \neq 0, \forall x,y \in [-1,1]\times[-1,1]$. Si el máximo no se alcanza en ningún sitio de dentro del cuadrado $(-1,1)\times(-1,1)$ se tendrá que alcanzar en el borde.
		\begin{itemize}
			\item Fijado $x = 1$ tenemos que $f(1,y) = \frac{1+y}{1+y} = 1 \implies f(1, -1 < y < 1) < 1$ porque si $f(1, -1 < y < 1)$ tomara un valor mayor que $1$ habría un máximo en $(-1,1)\times(-1,1)$ y esto no puede ser pues $\nabla f$ no se anula en el cuadrado.
			\item Fijado $x = -1$ tenemos que $f(-1, y) = \frac{y-1}{1-y} = -1 \implies f(-1, -1 < y < 1) > -1$ por la misma razón que antes.
			\item Hacemos lo mismo fijando la $y$ y variando la $x$.
		\end{itemize}
		En el borde (que no está incluido) se alcanzan máximo y mínimo que acotan a $f$ en el cuadrado:
		\begin{align*}
			-1 < f(x,y) = x \ast y < 1,\quad \forall x, y \in G
		\end{align*}
	\end{itemize}
\end{ex}

\begin{ex}[H1.3]Hallar los inversos de los siguientes elementos, cada uno en su grupo correspondiente:
	\begin{enumerate}
		\item $o(\overline{11})$ en $\uds{\Z^\ast/23\Z}$ es $22$ porque $23 \cdot 22 \equiv 1 \mod 23$
		\item $o(\overline{5})$ en $\uds{\Z^\ast/31\Z}$ es $3$ porque $5 \cdot 3 \equiv 1 \mod 31$
	\end{enumerate}
\end{ex}

\begin{ex}[H1.33]
	\label{ex:h1.33}
	Sea $G$ un grupo. Suponed que existe un único $a \in G$ de orden 2. Demostrad que $a \in Z(G)$.
\end{ex}

\begin{proof}
	Recordamos que $a \in Z(G) \iff ga = ag,\ \forall g \in G$. Definimos el isomorfismo de conjugación $\phi_g (x) = gx\inv{g}$ para algún $g$. Como $\phi_g$ es isomorfismo lleva elementos de orden $n$ en elementos de orden $n$. Entonces $\phi_g(a) = a$ ya que $a$ es el único elemento de orden 2. Por tanto $ga\inv{g} = a \implies ga = ag \implies a \in Z(G)$.
\end{proof}

\section{Hoja 2}

\begin{ex}[H2.1]
	\label{ex:h2.1}
	Se considera el tercer grupo diédrico $D_3$. Se pide hallar lo siguiente:
	\begin{enumerate}
		\item Las clases de conjugación de cada uno de sus elementos.
		\begin{proof}
			Las clases dan una partición del grupo. Si un elemento pertenece a una clase, entonces la clase de ese elemento también es la clase a la que pertenece.
			\begin{itemize}
				\item $cl(e) = \{e\}$
				\item $cl(B)$? Sabemos que $|cl(B)| = [G:C(B)]$. Sabemos que $\gen{B} =\{1, B, B^2\} \subset C(B)$ luego $|C(B)| \geq 3$. Si hubiera más elementos en $C(B)$ tendríamos que $|C(B)| = 6$ pues $C(B) < D_3$. Esto no ocurre porque sabemos que $B$ no conmuta con todos los demás elementos. Por ejemplo $BA \neq AB$. Por tanto $|C(B)| = 3 \implies |cl(B)| = [D_3:C(B)]= 6/3 = 2$. Es claro que $B \in cl(B)$. Además, como $cl(B)$ contiene elementos transformados por el isomorfismo conjugación sabemos que el otro elemento que hay tiene orden 3. El único elemento que queda de orden 3 es $B^2 \implies cl(B) = \{B, B^2\}$.
				
				\item $cl(A)?$ Sabemos que $A$ no conmuta con todos ($A \not\in Z(D_3)$) luego $|C(A)| < 6$. Sabemos que $\gen{A} = \{1, A\} < C(A)$. Además, como $C(A)$ es un (sub)grupo sabemos que no puede haber más elementos porque si los hubiera, $|\gen{A}| \divides |C(A)| \implies C(A) \geq 6$ pero ya hemos visto que no puede ser. Es decir que $|cl(A)| = [D_3:D(A)] = 6 / 2 = 3$. Por tanto $cl(A)$ incluye los 3 elementos que nos quedan: $cl(A) = \{A, AB, AB^2\}$.
			\end{itemize}
		\end{proof}
	
		\item Los elementos de $\text{Int}(D_3)$.
		\item Los centralizadores $C_{D_3}(x)$ para cada $x \in D_3$
		\item Los normalizadores $N(H)$ para cada $H < D_3$.
	\end{enumerate}
\end{ex}

\begin{ex}[H2.2]
	
	\begin{proof}
		Obtenidas las clases en el ejercicio \nameref{ex:h2.1} se verifica que $|D_3| = |cl(e)| + |cl(B)| + |cl(A)| = 1 + 2 + 3 = 6$
	\end{proof}
	
\end{ex}

\begin{ex}[H2.6]
	Sea $G$ un grupo. ¿Verdadero o falso?
	\begin{enumerate}
		\item $H < G$ y $H$ conmutativo implica $H \normsub G$.
		\item $H < G$ y $|H| = 2$ implica $H \normsub G$.
		\item Si $\varphi: G \to G_1$ es un homomorfismo de grupos, entonces $\ima \varphi \normsub G$
		\item Si $H \normsub K$ y $K \normsub G$ entonces $H \normsub G$
		\item Si $H \normsub G$ y $|H| = m$ entonces $H$ es el único subgrupo de $G$ de orden $m$.
		\item Si $H \normsub G$ entonces $H < Z(G)$.
		\begin{proof}[FALSO]
			Contraejemplo: En $G = D_4$ tomamos $H = \gen{B^2} = \{1, B, B^2, B^3\} \not\subset Z(D_4) = \{1, B^2\}$.
		\end{proof}
	\end{enumerate}
\end{ex}

\begin{ex}[H2.10]
	\begin{proof}
		Fijado $n$ y definida $\alpha_n : G \to G,\ x \mapsto x^n$ tenemos que $\alpha_n$ es un homomorfismo de grupos. Además podemos expresar $H_2 = \ker \alpha_n \implies H_2 \normsub G$. Además también tenemos que $H_1 = \ima \alpha_n < G$. Veamos que $H_1 \normsub G$. Es decir, que $gH_1\inv{g} = H_1,\ \forall g \in G$. Para ello tomamos $x_1^n \in H_1$ y lo conjugamos $gx_1^n\inv{g} = (gx_1\inv{g})^n$ por ser $\alpha$ homomorfismo de grupos. En particular $(gx_1\inv{g})^n \in \ima \alpha \implies (gx_1\inv{g})^n \in H_1 \implies (gx_1\inv{g})^n = x_2^n$ para algún $x_2 \in H_1 \implies H_1 \normsub G$.
	\end{proof}
\end{ex}

\begin{ex}[H2.13] Si $A$ es un grupo abeliano con $n$ elementos y $k$ es un entero primo con $n$, demostrad que la aplicación $\varphi : A \to A$ definida por $\varphi(a) = a^k$ es un isomorfismo.
	\begin{itemize}
		\item $\varphi$ homomorfismo de grupos.
		\begin{proof}
			\begin{align*}
				\varphi(a)\varphi(b) = a^kb^k = (ab)^k = \varphi(ab)
			\end{align*}
		\end{proof}
		\item $\varphi$ biyectiva $\iff \varphi$ inyectiva ya que dominio y codominio coinciden
		\begin{proof}
			$\ker \varphi = \{a \in A \mid \varphi(a) = a^k = 1\}$. Probaremos que $a^k = 1 \iff a = 1$ y por tanto que $\ker \varphi = \{1\} \implies \varphi$ inyectiva. Sabemos que $a^k = 1 \iff o(a^k) = 1$. Sea $t = o(a) \divides n$. Distinguimos dos casos
			\begin{itemize}
				\item Si $t = 1$ entonces $a = 1$ y ya está
				\item Si $t > 1$ entonces $o(a^k) = \frac{t}{mcd(k, t)} = \frac{t}{1} > 1$ contradicción. Luego necesariamente $t = o(a) = 1$.
			\end{itemize}
		\end{proof}
	\end{itemize}
\end{ex}

\begin{ex}[H2.22]
	\label{ex:h2.22}
	Demostrad que si $G$ es un grupo no conmutativo y tiene orden $p^3$ ($p$ un número primo) entonces $Z(G)$ tiene orden $p$.
	
	\begin{proof}
		Sabemos que $Z(G) < G \implies |Z(G)| \divides |G| \implies |Z(G)| \in \{1, p, p^2, p^3\}$
		\begin{itemize}
			\item $|Z(G)| \neq p^3$ porque en tal caso $G$ sería conmutativo
			\item $|Z(G)| \neq 1$ porque $G$ es un p-grupo y por tanto su centro no es el trivial.
			\item Si $|Z(G)| = p^2$ entonces $|G/Z(G)| = p \implies G/Z(G)$ es cíclico lo que no es posible si $G$ no es abeliano.
		\end{itemize}
	
		Por descarte concluimos que $|Z(G)| = p$.
	\end{proof}
\end{ex}


\begin{ex}[H2.25]
	Sabemos que $\autom{\Z/12\Z} \isom \uds{\Z^\star/12\Z}$ donde $\Z^\star/12\Z$ es el grupo multiplicativo $(\{\overline{1}, \overline{2}, \overline{3}, \dots,  \overline{11}\}, \cdot)$. Queda $\uds{\Z^\star/12\Z} = (\{\overline{1}, \overline{5}, \overline{7}, \overline{11}\}, \cdot )$ y además da la casualidad que $\forall x \in \uds{\Z^\star/12\Z},\ o(x) = 2$ (todos los elementos son su propio inverso) por lo que no tenemos restricciones al definir $f: \Z/2\Z \to \autom{\Z/12\Z}$:
	\begin{align*}
		f:\Z/2\Z & \to \autom{\Z/12\Z} \isom \uds{\Z^\star/12\Z} \\
		e = \overline{0} &\mapsto 1 \\
		\overline{1} &\mapsto \{\overline{1}, \overline{5}, \overline{7}, \overline{11}\}
	\end{align*}
\end{ex}

\begin{ex}[H2.26]
	Sea $|G_1| = m, |G_2| = n,\ mcd(m, n) = 1$. Si $f:G_1 \to G_2$ es h. de g. sabemos que $o(f(a)) \divides o(a),\ \forall a \in G_1$. Además $o(a) \divides m \land o(f(a)) \divides n$ por el teorema de Lagrange (\ref{thm:lagrange}).
	\begin{align*}
		\begin{cases}
		o(a) \divides m \land o(f(a)) \divides n \\
		mcd(m, n) = 1 \\
		o(f(a)) \divides o(a)
		\end{cases} \implies o(a) = o(f(a)) = 1, \forall a \in G_1
	\end{align*}
	Por lo que solo puede haber un homomorfismo entre ellos y además es el trivial $f(a) = e_{G_2}$.
\end{ex}

\begin{ex}[H2.19]
	Definimos una función $f:[0, 2\pi] \subset \R \to \mathbb{S}^1, \alpha\mapsto \cos \alpha + i \sin \alpha$. Esta función tiene la propiedad de que $f(\alpha) \cdot f(\alpha') = \cos(\alpha + \alpha') + i\sin(\alpha + \alpha') = f(\alpha + \alpha')$ y por tanto es un h. de g.\footnote{A la izquierda (en $\R$) sumamos pero a la derecha (en $\mathbb{S}^1$) multiplicamos.} entre $R$ y $\mathbb{S}^1$.
	
	Un elemento de $\cos \alpha + i \sin \alpha \in \mathbb{S}^1$ es de torsión $\iff \exists n \mid (\cos \alpha + i \sin \alpha )^n = 1$. Ahora bien $(\cos \alpha + i \sin \alpha )^n = \cos n\alpha + i \sin n\alpha = 1 \iff n\alpha = k 2\pi$. 
\end{ex}

\section{Hoja 4}

\begin{ex}[H4.11] Hallar los subgrupos de Sylow de $S_5$.
	Sabemos que $|S_5| = 5! = 2^3\cdot3\cdot 5$ y por \nameref{thm:sylow1} tenemos lo siguiente:
	\begin{itemize}
		\item $\exists P_2,\ |P_2| = 2^3 = 8$.
		\item $\exists P_3,\ |P_3| = 3$. Además en $S_5$ hay $\binom{5}{3}2! =20 $ 3-ciclos y en cada $gP_3\inv{g} = \{1, a, a^2 \mid o(a) = 3\}$ hay 2 elementos de orden 3 distintos. Además, $g_1P_3\inv{g_1} \cap g_2P_3\inv{g_2} = \{e\}$ porque si su intersección fuera más grande entonces serían el mismo subgrupo (porque son cíclicos). Es por esto que tenemos que repartir 40 elementos dando 2 a cada 3-grupo con lo que obtenemos $n_3 = 20 / 2 = 10$ 3-subgrupos de Sylow en $S_5$.
		\item $\exists P_5,\ |P_5| = 5$. Además en $S_5$ hay $4! = 24$ 5-ciclos (elementos de orden 5) y en cada $gP_5\inv{g} = \{1, a, a^2, a^3, a^4 \mid o(a) = 5\}$ tenemos $4$ elementos de orden $5$ distintos. Además, $g_1P_5\inv{g_1} \cap g_2P_5\inv{g_2} = \{e\}$ porque si su intersección fuera más grande entonces serían el mismo. Así, tenemos $24$ 5-ciclos a repartir entre los diferentes $gP_5\inv{g}$ dando 4 5-ciclos a cada 1. Por tanto tenemos $n_5 = 24/4 = 6$ 5-subgrupos de Sylow en $S_5$.
	\end{itemize}
\end{ex}

\begin{ex}[H4.18]
	Demostrar que todo grupo de orden $|G| = 5^3 \cdot 7^3$ tiene un subgrupo normal de orden 125.
	
	\begin{proof}
		\nameref{thm:sylow1} $\implies \exists P_5 < G,\ |P_5| = 125$. \nameref{thm:sylow3} $\implies n_5 \divides 7^3 \land n_5 \equiv 1 \mod 5$ es decir $n_5 \in \{1, 7, 49, 343\} \land n_5 \in \{1, 6, 11, \dots\}$. Como ni $49$ ni $343$ son conguentes con $1$ módulo 5 tenemos que $n_5 = 1 \implies P_5 \normsub G$.
	\end{proof}
\end{ex}

\begin{ex}[H4.20] Hallad todos los grupos abelianos de órdenes 36, 64, 96 y 100.
	\begin{enumerate}
		\item $|G| = 36 = 2^2\cdot3^2$
		\begin{proof}
			\nameref{thm:sylow1} $\implies \exists P_2,\ |P_2| = 4 \land \exists P_3,\ |P_3| = 9$. Además $G$ abeliano $\implies P_2, P_3 \normsub G \implies G \isom P_2 \times P_3$. Estudiamos los grupos de orden 4 y de orden 9
			\begin{itemize}
				\item $|P_2| = 4$ entonces $P_2 \isom \Z/2\Z \times \Z/2\Z \lor P_2 \isom \Z/4\Z$
				\item $|P_3| = 9$ entonces $P_3 \isom \Z/3\Z \times \Z/3\Z \lor P_3 \isom \Z/9\Z$
			\end{itemize}
			Como $G \isom P_2 \times P_3$ tenemos 4 posibles grupos abelianos de orden 36.
		\end{proof}
	\end{enumerate}
\end{ex}

\begin{ex}[H4.22]
	Hallar todos los grupos abelianos de orden 175.
	
	\begin{proof}
		$|G| = 5^2\cdot 7$. Por el \nameref{thm:sylow1} tenemos que $\exists P_5, P_7 < G$ con $|P_5| = 25,\ |P_7| = 7$ y además por ser $G$ abeliano tenemos que $P_5, P_7 \normsub G \implies G \isom P_5 \times P_7$. Estudiamos los grupos de orden 25 y de orden 7:
		\begin{itemize}
			\item $|P_5| = 25 \land P_5$ abeliano $\implies P_5 \isom \Z/25\Z \lor P_5 \isom \Z/5\Z \times \Z/5\Z$. En ambos casos $P_5$ es producto directo de cíclicos pues $\Z/n\Z$ es cíclico.
			\item $|P_7| = 7 \land P_7$ abeliano $\implies P_7 \isom \Z/7\Z$. Ocurre lo mismo que con $P_5$.
		\end{itemize}
		Concluimos que $G \isom \Z/25\Z \times \Z/7\Z \lor G \isom \Z/5\Z \times \Z/5\Z \times \Z/7\Z$. Los dos casos son abelianos por ser producto directo de grupos cíclicos.
	\end{proof}
\end{ex}

\begin{ex}[H4.23]
	¿Cuántos elementos de orden 3 puede tener un grupo abeliano de orden 36?
	
	\begin{proof}
		$|G| = 36 = 2^2 3^2$. \nameref{thm:sylow1} $\implies \exists P_2, P_3 < G,\ |P_2| = 4,\ |P_3| = 9$. Estudiamos los grupos de órdenes 4 y 9:
		\begin{itemize}
			\item $|P_3| = 9 \implies P_3 \isom \Z/9\Z \lor P_3 \isom \Z/3\Z \times \Z/3\Z$
			% TODO: terminar
		\end{itemize}
	\end{proof}
\end{ex}

%\renewcommand{\chaptername}{Apéndice}
%\renewcommand{\thechapter}{\Alph{chapter}}
%\setcounter{chapter}{0}

%% !TeX root = ../apuntes-ea.tex

\chapter{Ejercicios}

\section{Hoja 1}

\begin{ex}[H1.2]
	Sean $a,b,c \in G = (-1,1)$. Probamos las propiedades de los grupos.
	\begin{itemize}
		\item \textbf{Asociatividad:}
		\begin{align*}
			(a \ast b) \ast c = \left(\frac{a+b}{1+ab}\right) \ast c = \frac{\left(\frac{a+b}{1+ab}\right) + c}{1 + \left(\frac{a+b}{1+ab}\right)c} = \frac{\frac{a+b+c+abc}{1+ab}}{\frac{1+ab+ac+bc}{1+ab}} = \frac{a+b+c+abc}{1+ab+ac+bc} \\
			a \ast (b \ast c) = a \ast \left(\frac{b+c}{1+bc}\right) = \frac{a+\left(\frac{b+c}{1+bc}\right)}{1+a\left(\frac{b+c}{1+bc}\right)} = \frac{\frac{a + b + c + abc}{1+bc}}{\frac{1+ab+ac+bc}{1+bc}} = \frac{a+b+c+abc}{1+ab+ac+bc}
		\end{align*}
		\item \textbf{Elemento neutro:} es el $0$ ya que $x\ast 0 = \frac{x+0}{1+x\cdot 0} = \frac{x}{1} = x$ y además $0 \ast x = \frac{0+x}{1+0\cdot x} = \frac{x}{1} = x$
		\item \textbf{Elemento inverso:} la ecuación
		\begin{align*}
			x \ast \inv{x} = 0 \iff \frac{x + \inv{x}}{1+x\inv{x}} = 0 \iff \inv{x} = -x
		\end{align*}
		siempre tiene solución y ocurre lo mismo para la ecuación $\inv{x} \ast x = 0 \iff \inv{x} = -x$
		\item \textbf{Clausura:} tenemos que probar que si $x,y \in (-1,1)$ entonces $x \ast y \in (-1,1)$. Consideramos $f(x,y) = x \ast y = \frac{x+y}{1+xy}$. Derivando tenemos que $\nabla f(x,y) = (\frac{1}{(1+xy)^2},\ \frac{1}{(1+xy)^2}) \neq 0, \forall x,y \in [-1,1]\times[-1,1]$. Si el máximo no se alcanza en ningún sitio de dentro del cuadrado $(-1,1)\times(-1,1)$ se tendrá que alcanzar en el borde.
		\begin{itemize}
			\item Fijado $x = 1$ tenemos que $f(1,y) = \frac{1+y}{1+y} = 1 \implies f(1, -1 < y < 1) < 1$ porque si $f(1, -1 < y < 1)$ tomara un valor mayor que $1$ habría un máximo en $(-1,1)\times(-1,1)$ y esto no puede ser pues $\nabla f$ no se anula en el cuadrado.
			\item Fijado $x = -1$ tenemos que $f(-1, y) = \frac{y-1}{1-y} = -1 \implies f(-1, -1 < y < 1) > -1$ por la misma razón que antes.
			\item Hacemos lo mismo fijando la $y$ y variando la $x$.
		\end{itemize}
		En el borde (que no está incluido) se alcanzan máximo y mínimo que acotan a $f$ en el cuadrado:
		\begin{align*}
			-1 < f(x,y) = x \ast y < 1,\quad \forall x, y \in G
		\end{align*}
	\end{itemize}
\end{ex}

\begin{ex}[H1.3]Hallar los inversos de los siguientes elementos, cada uno en su grupo correspondiente:
	\begin{enumerate}
		\item $o(\overline{11})$ en $\uds{\Z^\ast/23\Z}$ es $22$ porque $23 \cdot 22 \equiv 1 \mod 23$
		\item $o(\overline{5})$ en $\uds{\Z^\ast/31\Z}$ es $3$ porque $5 \cdot 3 \equiv 1 \mod 31$
	\end{enumerate}
\end{ex}

\begin{ex}[H1.33]
	\label{ex:h1.33}
	Sea $G$ un grupo. Suponed que existe un único $a \in G$ de orden 2. Demostrad que $a \in Z(G)$.
\end{ex}

\begin{proof}
	Recordamos que $a \in Z(G) \iff ga = ag,\ \forall g \in G$. Definimos el isomorfismo de conjugación $\phi_g (x) = gx\inv{g}$ para algún $g$. Como $\phi_g$ es isomorfismo lleva elementos de orden $n$ en elementos de orden $n$. Entonces $\phi_g(a) = a$ ya que $a$ es el único elemento de orden 2. Por tanto $ga\inv{g} = a \implies ga = ag \implies a \in Z(G)$.
\end{proof}

\section{Hoja 2}

\begin{ex}[H2.1]
	\label{ex:h2.1}
	Se considera el tercer grupo diédrico $D_3$. Se pide hallar lo siguiente:
	\begin{enumerate}
		\item Las clases de conjugación de cada uno de sus elementos.
		\begin{proof}
			Las clases dan una partición del grupo. Si un elemento pertenece a una clase, entonces la clase de ese elemento también es la clase a la que pertenece.
			\begin{itemize}
				\item $cl(e) = \{e\}$
				\item $cl(B)$? Sabemos que $|cl(B)| = [G:C(B)]$. Sabemos que $\gen{B} =\{1, B, B^2\} \subset C(B)$ luego $|C(B)| \geq 3$. Si hubiera más elementos en $C(B)$ tendríamos que $|C(B)| = 6$ pues $C(B) < D_3$. Esto no ocurre porque sabemos que $B$ no conmuta con todos los demás elementos. Por ejemplo $BA \neq AB$. Por tanto $|C(B)| = 3 \implies |cl(B)| = [D_3:C(B)]= 6/3 = 2$. Es claro que $B \in cl(B)$. Además, como $cl(B)$ contiene elementos transformados por el isomorfismo conjugación sabemos que el otro elemento que hay tiene orden 3. El único elemento que queda de orden 3 es $B^2 \implies cl(B) = \{B, B^2\}$.
				
				\item $cl(A)?$ Sabemos que $A$ no conmuta con todos ($A \not\in Z(D_3)$) luego $|C(A)| < 6$. Sabemos que $\gen{A} = \{1, A\} < C(A)$. Además, como $C(A)$ es un (sub)grupo sabemos que no puede haber más elementos porque si los hubiera, $|\gen{A}| \divides |C(A)| \implies C(A) \geq 6$ pero ya hemos visto que no puede ser. Es decir que $|cl(A)| = [D_3:D(A)] = 6 / 2 = 3$. Por tanto $cl(A)$ incluye los 3 elementos que nos quedan: $cl(A) = \{A, AB, AB^2\}$.
			\end{itemize}
		\end{proof}
	
		\item Los elementos de $\text{Int}(D_3)$.
		\item Los centralizadores $C_{D_3}(x)$ para cada $x \in D_3$
		\item Los normalizadores $N(H)$ para cada $H < D_3$.
	\end{enumerate}
\end{ex}

\begin{ex}[H2.2]
	
	\begin{proof}
		Obtenidas las clases en el ejercicio \nameref{ex:h2.1} se verifica que $|D_3| = |cl(e)| + |cl(B)| + |cl(A)| = 1 + 2 + 3 = 6$
	\end{proof}
	
\end{ex}

\begin{ex}[H2.6]
	Sea $G$ un grupo. ¿Verdadero o falso?
	\begin{enumerate}
		\item $H < G$ y $H$ conmutativo implica $H \normsub G$.
		\item $H < G$ y $|H| = 2$ implica $H \normsub G$.
		\item Si $\varphi: G \to G_1$ es un homomorfismo de grupos, entonces $\ima \varphi \normsub G$
		\item Si $H \normsub K$ y $K \normsub G$ entonces $H \normsub G$
		\item Si $H \normsub G$ y $|H| = m$ entonces $H$ es el único subgrupo de $G$ de orden $m$.
		\item Si $H \normsub G$ entonces $H < Z(G)$.
		\begin{proof}[FALSO]
			Contraejemplo: En $G = D_4$ tomamos $H = \gen{B^2} = \{1, B, B^2, B^3\} \not\subset Z(D_4) = \{1, B^2\}$.
		\end{proof}
	\end{enumerate}
\end{ex}

\begin{ex}[H2.10]
	\begin{proof}
		Fijado $n$ y definida $\alpha_n : G \to G,\ x \mapsto x^n$ tenemos que $\alpha_n$ es un homomorfismo de grupos. Además podemos expresar $H_2 = \ker \alpha_n \implies H_2 \normsub G$. Además también tenemos que $H_1 = \ima \alpha_n < G$. Veamos que $H_1 \normsub G$. Es decir, que $gH_1\inv{g} = H_1,\ \forall g \in G$. Para ello tomamos $x_1^n \in H_1$ y lo conjugamos $gx_1^n\inv{g} = (gx_1\inv{g})^n$ por ser $\alpha$ homomorfismo de grupos. En particular $(gx_1\inv{g})^n \in \ima \alpha \implies (gx_1\inv{g})^n \in H_1 \implies (gx_1\inv{g})^n = x_2^n$ para algún $x_2 \in H_1 \implies H_1 \normsub G$.
	\end{proof}
\end{ex}

\begin{ex}[H2.13] Si $A$ es un grupo abeliano con $n$ elementos y $k$ es un entero primo con $n$, demostrad que la aplicación $\varphi : A \to A$ definida por $\varphi(a) = a^k$ es un isomorfismo.
	\begin{itemize}
		\item $\varphi$ homomorfismo de grupos.
		\begin{proof}
			\begin{align*}
				\varphi(a)\varphi(b) = a^kb^k = (ab)^k = \varphi(ab)
			\end{align*}
		\end{proof}
		\item $\varphi$ biyectiva $\iff \varphi$ inyectiva ya que dominio y codominio coinciden
		\begin{proof}
			$\ker \varphi = \{a \in A \mid \varphi(a) = a^k = 1\}$. Probaremos que $a^k = 1 \iff a = 1$ y por tanto que $\ker \varphi = \{1\} \implies \varphi$ inyectiva. Sabemos que $a^k = 1 \iff o(a^k) = 1$. Sea $t = o(a) \divides n$. Distinguimos dos casos
			\begin{itemize}
				\item Si $t = 1$ entonces $a = 1$ y ya está
				\item Si $t > 1$ entonces $o(a^k) = \frac{t}{mcd(k, t)} = \frac{t}{1} > 1$ contradicción. Luego necesariamente $t = o(a) = 1$.
			\end{itemize}
		\end{proof}
	\end{itemize}
\end{ex}

\begin{ex}[H2.22]
	\label{ex:h2.22}
	Demostrad que si $G$ es un grupo no conmutativo y tiene orden $p^3$ ($p$ un número primo) entonces $Z(G)$ tiene orden $p$.
	
	\begin{proof}
		Sabemos que $Z(G) < G \implies |Z(G)| \divides |G| \implies |Z(G)| \in \{1, p, p^2, p^3\}$
		\begin{itemize}
			\item $|Z(G)| \neq p^3$ porque en tal caso $G$ sería conmutativo
			\item $|Z(G)| \neq 1$ porque $G$ es un p-grupo y por tanto su centro no es el trivial.
			\item Si $|Z(G)| = p^2$ entonces $|G/Z(G)| = p \implies G/Z(G)$ es cíclico lo que no es posible si $G$ no es abeliano.
		\end{itemize}
	
		Por descarte concluimos que $|Z(G)| = p$.
	\end{proof}
\end{ex}


\begin{ex}[H2.25]
	Sabemos que $\autom{\Z/12\Z} \isom \uds{\Z^\star/12\Z}$ donde $\Z^\star/12\Z$ es el grupo multiplicativo $(\{\overline{1}, \overline{2}, \overline{3}, \dots,  \overline{11}\}, \cdot)$. Queda $\uds{\Z^\star/12\Z} = (\{\overline{1}, \overline{5}, \overline{7}, \overline{11}\}, \cdot )$ y además da la casualidad que $\forall x \in \uds{\Z^\star/12\Z},\ o(x) = 2$ (todos los elementos son su propio inverso) por lo que no tenemos restricciones al definir $f: \Z/2\Z \to \autom{\Z/12\Z}$:
	\begin{align*}
		f:\Z/2\Z & \to \autom{\Z/12\Z} \isom \uds{\Z^\star/12\Z} \\
		e = \overline{0} &\mapsto 1 \\
		\overline{1} &\mapsto \{\overline{1}, \overline{5}, \overline{7}, \overline{11}\}
	\end{align*}
\end{ex}

\begin{ex}[H2.26]
	Sea $|G_1| = m, |G_2| = n,\ mcd(m, n) = 1$. Si $f:G_1 \to G_2$ es h. de g. sabemos que $o(f(a)) \divides o(a),\ \forall a \in G_1$. Además $o(a) \divides m \land o(f(a)) \divides n$ por el teorema de Lagrange (\ref{thm:lagrange}).
	\begin{align*}
		\begin{cases}
		o(a) \divides m \land o(f(a)) \divides n \\
		mcd(m, n) = 1 \\
		o(f(a)) \divides o(a)
		\end{cases} \implies o(a) = o(f(a)) = 1, \forall a \in G_1
	\end{align*}
	Por lo que solo puede haber un homomorfismo entre ellos y además es el trivial $f(a) = e_{G_2}$.
\end{ex}

\begin{ex}[H2.19]
	Definimos una función $f:[0, 2\pi] \subset \R \to \mathbb{S}^1, \alpha\mapsto \cos \alpha + i \sin \alpha$. Esta función tiene la propiedad de que $f(\alpha) \cdot f(\alpha') = \cos(\alpha + \alpha') + i\sin(\alpha + \alpha') = f(\alpha + \alpha')$ y por tanto es un h. de g.\footnote{A la izquierda (en $\R$) sumamos pero a la derecha (en $\mathbb{S}^1$) multiplicamos.} entre $R$ y $\mathbb{S}^1$.
	
	Un elemento de $\cos \alpha + i \sin \alpha \in \mathbb{S}^1$ es de torsión $\iff \exists n \mid (\cos \alpha + i \sin \alpha )^n = 1$. Ahora bien $(\cos \alpha + i \sin \alpha )^n = \cos n\alpha + i \sin n\alpha = 1 \iff n\alpha = k 2\pi$. 
\end{ex}

\section{Hoja 4}

\begin{ex}[H4.11] Hallar los subgrupos de Sylow de $S_5$.
	Sabemos que $|S_5| = 5! = 2^3\cdot3\cdot 5$ y por \nameref{thm:sylow1} tenemos lo siguiente:
	\begin{itemize}
		\item $\exists P_2,\ |P_2| = 2^3 = 8$.
		\item $\exists P_3,\ |P_3| = 3$. Además en $S_5$ hay $\binom{5}{3}2! =20 $ 3-ciclos y en cada $gP_3\inv{g} = \{1, a, a^2 \mid o(a) = 3\}$ hay 2 elementos de orden 3 distintos. Además, $g_1P_3\inv{g_1} \cap g_2P_3\inv{g_2} = \{e\}$ porque si su intersección fuera más grande entonces serían el mismo subgrupo (porque son cíclicos). Es por esto que tenemos que repartir 40 elementos dando 2 a cada 3-grupo con lo que obtenemos $n_3 = 20 / 2 = 10$ 3-subgrupos de Sylow en $S_5$.
		\item $\exists P_5,\ |P_5| = 5$. Además en $S_5$ hay $4! = 24$ 5-ciclos (elementos de orden 5) y en cada $gP_5\inv{g} = \{1, a, a^2, a^3, a^4 \mid o(a) = 5\}$ tenemos $4$ elementos de orden $5$ distintos. Además, $g_1P_5\inv{g_1} \cap g_2P_5\inv{g_2} = \{e\}$ porque si su intersección fuera más grande entonces serían el mismo. Así, tenemos $24$ 5-ciclos a repartir entre los diferentes $gP_5\inv{g}$ dando 4 5-ciclos a cada 1. Por tanto tenemos $n_5 = 24/4 = 6$ 5-subgrupos de Sylow en $S_5$.
	\end{itemize}
\end{ex}

\begin{ex}[H4.18]
	Demostrar que todo grupo de orden $|G| = 5^3 \cdot 7^3$ tiene un subgrupo normal de orden 125.
	
	\begin{proof}
		\nameref{thm:sylow1} $\implies \exists P_5 < G,\ |P_5| = 125$. \nameref{thm:sylow3} $\implies n_5 \divides 7^3 \land n_5 \equiv 1 \mod 5$ es decir $n_5 \in \{1, 7, 49, 343\} \land n_5 \in \{1, 6, 11, \dots\}$. Como ni $49$ ni $343$ son conguentes con $1$ módulo 5 tenemos que $n_5 = 1 \implies P_5 \normsub G$.
	\end{proof}
\end{ex}

\begin{ex}[H4.20] Hallad todos los grupos abelianos de órdenes 36, 64, 96 y 100.
	\begin{enumerate}
		\item $|G| = 36 = 2^2\cdot3^2$
		\begin{proof}
			\nameref{thm:sylow1} $\implies \exists P_2,\ |P_2| = 4 \land \exists P_3,\ |P_3| = 9$. Además $G$ abeliano $\implies P_2, P_3 \normsub G \implies G \isom P_2 \times P_3$. Estudiamos los grupos de orden 4 y de orden 9
			\begin{itemize}
				\item $|P_2| = 4$ entonces $P_2 \isom \Z/2\Z \times \Z/2\Z \lor P_2 \isom \Z/4\Z$
				\item $|P_3| = 9$ entonces $P_3 \isom \Z/3\Z \times \Z/3\Z \lor P_3 \isom \Z/9\Z$
			\end{itemize}
			Como $G \isom P_2 \times P_3$ tenemos 4 posibles grupos abelianos de orden 36.
		\end{proof}
	\end{enumerate}
\end{ex}

\begin{ex}[H4.22]
	Hallar todos los grupos abelianos de orden 175.
	
	\begin{proof}
		$|G| = 5^2\cdot 7$. Por el \nameref{thm:sylow1} tenemos que $\exists P_5, P_7 < G$ con $|P_5| = 25,\ |P_7| = 7$ y además por ser $G$ abeliano tenemos que $P_5, P_7 \normsub G \implies G \isom P_5 \times P_7$. Estudiamos los grupos de orden 25 y de orden 7:
		\begin{itemize}
			\item $|P_5| = 25 \land P_5$ abeliano $\implies P_5 \isom \Z/25\Z \lor P_5 \isom \Z/5\Z \times \Z/5\Z$. En ambos casos $P_5$ es producto directo de cíclicos pues $\Z/n\Z$ es cíclico.
			\item $|P_7| = 7 \land P_7$ abeliano $\implies P_7 \isom \Z/7\Z$. Ocurre lo mismo que con $P_5$.
		\end{itemize}
		Concluimos que $G \isom \Z/25\Z \times \Z/7\Z \lor G \isom \Z/5\Z \times \Z/5\Z \times \Z/7\Z$. Los dos casos son abelianos por ser producto directo de grupos cíclicos.
	\end{proof}
\end{ex}

\begin{ex}[H4.23]
	¿Cuántos elementos de orden 3 puede tener un grupo abeliano de orden 36?
	
	\begin{proof}
		$|G| = 36 = 2^2 3^2$. \nameref{thm:sylow1} $\implies \exists P_2, P_3 < G,\ |P_2| = 4,\ |P_3| = 9$. Estudiamos los grupos de órdenes 4 y 9:
		\begin{itemize}
			\item $|P_3| = 9 \implies P_3 \isom \Z/9\Z \lor P_3 \isom \Z/3\Z \times \Z/3\Z$
			% TODO: terminar
		\end{itemize}
	\end{proof}
\end{ex}

\chapter{\'{I}ndices}

%%TODO por qué no podemos utilizar only= para hacer whitelisting en lugar de blacklisting? que alguien lo arregle parfavar
\renewcommand{\listtheoremname}{Lista de definiciones}
\listoftheorems[ignore={thm,eg,pro,cor,obs,lm,ex,th_ex}]

\renewcommand{\listtheoremname}{Lista de teoremas}
\listoftheorems[onlynamed,ignore={dfn,eg,cor,obs,ex,th_ex}]

\renewcommand{\listtheoremname}{Lista de ejemplos}
\listoftheorems[onlynamed,ignore={dfn,thm,pro,cor,obs,lm,ex,th_ex}]

\renewcommand{\listtheoremname}{Lista de ejercicios}
\listoftheorems[ignore={dfn,thm,pro,cor,obs,lm,eg}]

\printglossaries

\bibliographystyle{alpha}

\end{document}
